\documentclass[12pt,a4paper,twoside]{article}

% === Chinese language support ===
\usepackage{ctex}

% === Page layout ===
\usepackage[a4paper,top=25mm,bottom=25mm,left=28mm,right=22mm,
            headheight=15pt,headsep=8mm,footskip=12mm]{geometry}

% === Core packages ===
\usepackage{graphicx}
\usepackage{xcolor}
\usepackage{booktabs}
\usepackage{longtable}
\usepackage{tabularx}
\usepackage{array}
\usepackage{multirow}
\usepackage{amsmath}
\usepackage{amssymb}
\usepackage{enumitem}
\usepackage{fancyhdr}
\usepackage{titlesec}

\usepackage{hyperref}
\usepackage{listings}
\usepackage{float}
\usepackage{caption}
\usepackage{subcaption}
\usepackage{microtype}
\usepackage{natbib}
\usepackage{setspace}
\usepackage{soul}
\usepackage{pifont}
\usepackage{wrapfig}
\usepackage{tikz}
\usetikzlibrary{shapes,arrows,positioning,calc}
\usepackage{tcolorbox}
\tcbuselibrary{skins,breakable,listings}

% === Color palette ===
\definecolor{darkblue}{RGB}{0,51,102}
\definecolor{medblue}{RGB}{0,86,163}
\definecolor{teal}{RGB}{0,128,128}
\definecolor{lightblue}{RGB}{220,235,250}
\definecolor{lightgreen}{RGB}{232,248,232}
\definecolor{lightyellow}{RGB}{255,253,220}
\definecolor{lightred}{RGB}{255,235,235}
\definecolor{codebg}{RGB}{245,245,245}
\definecolor{warningbg}{RGB}{255,248,230}
\definecolor{warningborder}{RGB}{204,153,0}
\definecolor{tipbg}{RGB}{230,245,255}
\definecolor{tipborder}{RGB}{0,102,204}
\definecolor{dangerbg}{RGB}{255,235,235}
\definecolor{dangerborder}{RGB}{204,0,0}
\definecolor{successbg}{RGB}{232,248,232}
\definecolor{successborder}{RGB}{0,128,0}

% === Hyperref setup ===
\hypersetup{
  colorlinks=true,
  linkcolor=darkblue,
  citecolor=teal,
  urlcolor=medblue,
  pdftitle={中药多维数据分析筛选免疫激活抗肿瘤药物全流程可落地操作手册},
  pdfauthor={K-Dense Web},
  pdfsubject={TCM Drug Screening SOP},
  pdfkeywords={TCM, NPC, network pharmacology, molecular docking, immunotherapy},
  bookmarks=true,
  bookmarksnumbered=true,
  pdfstartview=FitH
}

% === Section title formatting ===
\titleformat{\section}[block]
  {\normalfont\Large\bfseries\color{darkblue}}
  {\thesection}{1em}{}
  [{\color{darkblue}\titlerule[1.5pt]}]

\titleformat{\subsection}[block]
  {\normalfont\large\bfseries\color{medblue}}
  {\thesubsection}{1em}{}

\titleformat{\subsubsection}[block]
  {\normalfont\normalsize\bfseries\color{teal}}
  {\thesubsubsection}{1em}{}

\titlespacing*{\section}{0pt}{18pt}{10pt}
\titlespacing*{\subsection}{0pt}{14pt}{8pt}
\titlespacing*{\subsubsection}{0pt}{10pt}{6pt}

% === Header/Footer ===
\pagestyle{fancy}
\fancyhf{}
\fancyhead[LE]{\small\color{darkblue}\leftmark}
\fancyhead[RO]{\small\color{darkblue}中药多维数据分析筛选免疫激活抗肿瘤药物操作手册}
\fancyfoot[C]{\small\thepage}
\fancyfoot[LE]{\small\color{gray}K-Dense Web · [REDACTED-EMAIL]}
\fancyfoot[RO]{\small\color{gray}Version 3.0 · 2026-03-05}
\renewcommand{\headrulewidth}{0.5pt}
\renewcommand{\footrulewidth}{0.3pt}

% === Custom tcolorbox environments ===
\newtcolorbox{warningbox}[1][]{
  breakable,
  colback=warningbg,
  colframe=warningborder,
  title={\textbf{[!] 避坑指南}},
  fonttitle=\bfseries\color{white},
  coltitle=white,
  colbacktitle=warningborder,
  boxrule=0.8pt,
  arc=3pt,
  #1
}

\newtcolorbox{tipbox}[1][]{
  breakable,
  colback=tipbg,
  colframe=tipborder,
  title={\textbf{科研进阶技巧}},
  fonttitle=\bfseries\color{white},
  coltitle=white,
  colbacktitle=tipborder,
  boxrule=0.8pt,
  arc=3pt,
  #1
}

\newtcolorbox{qcbox}[1][]{
  breakable,
  colback=successbg,
  colframe=successborder,
  title={\textbf{[OK] 质控标准}},
  fonttitle=\bfseries\color{white},
  coltitle=white,
  colbacktitle=successborder,
  boxrule=0.8pt,
  arc=3pt,
  #1
}

\newtcolorbox{infobox}[2][]{
  breakable,
  colback=lightblue,
  colframe=medblue,
  title={\textbf{#2}},
  fonttitle=\bfseries\color{white},
  coltitle=white,
  colbacktitle=medblue,
  boxrule=0.8pt,
  arc=3pt,
  #1
}

\newtcolorbox{templatebox}[1][]{
  breakable,
  colback=lightyellow,
  colframe=orange!70,
  title={\textbf{【模板替换位】}},
  fonttitle=\bfseries,
  coltitle=darkblue,
  colbacktitle=orange!30,
  boxrule=0.8pt,
  arc=3pt,
  #1
}

\newtcolorbox{greenbox}[1][]{
  breakable,
  colback=lightgreen,
  colframe=teal,
  boxrule=0.8pt,
  arc=3pt,
  #1
}

% === Code listings ===
\lstset{
  basicstyle=\ttfamily\footnotesize,
  breaklines=true,
  breakatwhitespace=false,
  numbers=left,
  numberstyle=\tiny\color{gray},
  numbersep=8pt,
  frame=single,
  rulecolor=\color{gray!50},
  backgroundcolor=\color{codebg},
  showstringspaces=false,
  keywordstyle=\color{blue!80}\bfseries,
  commentstyle=\color{green!50!black}\itshape,
  stringstyle=\color{red!70},
  tabsize=2,
  xleftmargin=15pt,
  xrightmargin=5pt
}

% === List spacing ===
\setlist{nosep,leftmargin=*,topsep=4pt,parsep=2pt,itemsep=3pt}
\setlist[enumerate]{label=\arabic*.}
\setlist[itemize]{label=\textbullet}

% === Caption setup ===
\captionsetup{
  font=small,
  labelfont={bf,color=darkblue},
  margin=10pt
}

% === Table of contents ===
\setcounter{tocdepth}{3}

% === Spacing ===
\setlength{\parindent}{2em}
\setlength{\parskip}{4pt}
\onehalfspacing

% === Custom commands ===
\newcommand{\software}[1]{\textbf{\textit{#1}}}
\newcommand{\database}[1]{\textbf{\textsf{#1}}}
\newcommand{\gene}[1]{\textit{#1}}
\newcommand{\clickpath}[1]{\texttt{#1}}
\newcommand{\moduleoutput}[1]{\par\vspace{6pt}\noindent\textcolor{teal}{\textbf{$\rightarrow$ 本模块输出：}}#1\vspace{6pt}}

\begin{document}
\sloppy

% ============================================================
% TITLE PAGE
% ============================================================
\begin{titlepage}
  \begin{center}
    \vspace*{15mm}
    {\color{darkblue}\rule{\textwidth}{2pt}}
    \vspace{8mm}
    
    {\fontsize{22}{28}\selectfont\bfseries\color{darkblue}
    中药多维数据分析筛选免疫激活抗肿瘤药物\\[6pt]
    全流程可落地操作手册}
    
    \vspace{4mm}
    {\large\color{medblue}\itshape
    Traditional Chinese Medicine Multi-Dimensional Data Analysis\\
    for Immune-Activating Anti-Tumor Drug Screening:\\
    Full-Process Actionable Operation Manual}
    
    \vspace{8mm}
    {\color{darkblue}\rule{\textwidth}{1pt}}
    
    \vspace{12mm}
    \includegraphics[width=0.95\textwidth]{../figures/graphical_abstract_v2.png}
    
    \vspace{10mm}
    
    \begin{tabular}{ll}
      \textbf{核心肿瘤模型：} & 鼻咽癌 (NPC)，CNE1/CNE2 细胞系 \\
      \textbf{研究终点：} & 直接抗肿瘤活性 + 免疫激活双终点 \\
      \textbf{手册版本：} & Version 3.0 \\
      \textbf{发布日期：} & 2026 年 3 月 5 日 \\
      \textbf{适用场景：} & 硕士科研 · SCI 论文 · 课题申报 · 专利申请 \\
    \end{tabular}
    
    \vspace{12mm}
    {\large\bfseries K-Dense Web 研究团队}\\
    \vspace{2mm}
    {\normalsize\texttt{[REDACTED-EMAIL]}}
    
    \vspace{8mm}
    {\color{darkblue}\rule{\textwidth}{2pt}}
  \end{center}
\end{titlepage}

% ============================================================
% FRONT MATTER
% ============================================================
\setcounter{page}{1}
\pagenumbering{roman}

% --- Disclaimer ---
\section*{免责声明与使用规范}
\addcontentsline{toc}{section}{免责声明与使用规范}

本手册所有操作流程、参数设置与数据库资源均基于截至 2026 年 3 月的公开发表文献与权威数据库规范，仅供科研学习与教学参考。本手册不构成医疗建议，不可用于临床诊疗决策。所有湿实验操作必须遵循所在机构的生物安全规范与伦理委员会批准要求。研究者在使用本手册前，需自行评估所在机构的软硬件条件、伦理合规状态及数据使用权限。

\begin{itemize}
  \item \textbf{版权声明：}本手册内容受著作权保护，禁止以营利为目的的商业复制与传播。
  \item \textbf{引用规范：}使用本手册方法时，应同时引用对应章节所注明的原始数据库与工具论文。
  \item \textbf{数据库访问：}本手册标注的数据库 URL 于发布时有效，因数据库更新导致的 URL 变动，以官方公告为准。
  \item \textbf{细胞系鉴定：}CNE1/CNE2 细胞系已被 ICLAC 列为存在 HeLa 细胞污染的细胞系，研究者在发表论文时需参照所在期刊的细胞系鉴定要求执行 STR 鉴定，并在论文中予以说明。
\end{itemize}

\clearpage

% --- How to Use ---
\section*{手册使用指南}
\addcontentsline{toc}{section}{手册使用指南}

\subsection*{手册结构概览}

本手册采用「\textbf{9 大核心模块}」的线性逻辑结构，前一模块的输出作为后一模块的直接输入，实现全流程无断点闭环：

\begin{center}
\begin{tikzpicture}[
  node distance=0.7cm,
  box/.style={rectangle,draw=darkblue,fill=lightblue,rounded corners=3pt,
              text width=3cm,align=center,font=\small\bfseries,minimum height=0.8cm},
  arrow/.style={->,>=stealth,darkblue,line width=1pt}
]
\node[box] (m1) {模块一\\前期准备};
\node[box,right=of m1] (m2) {模块二\\数据库构建};
\node[box,right=of m2] (m3) {模块三\\靶点集筛选};
\node[box,right=of m3] (m4) {模块四\\虚拟筛选};
\node[box,right=of m4] (m5) {模块五\\多组学验证};
\node[box,below=1.2cm of m5] (m6) {模块六\\湿实验验证};
\node[box,left=of m6] (m7) {模块七\\数据分析};
\node[box,left=of m7] (m8) {模块八\\成果输出};
\node[box,left=of m8] (m9) {模块九\\质控风控};
\draw[arrow] (m1)--(m2);
\draw[arrow] (m2)--(m3);
\draw[arrow] (m3)--(m4);
\draw[arrow] (m4)--(m5);
\draw[arrow] (m5)--(m6);
\draw[arrow] (m6)--(m7);
\draw[arrow] (m7)--(m8);
\draw[arrow] (m8)--(m9);
\end{tikzpicture}
\end{center}

\subsection*{标注体系说明}

本手册采用统一标注规范，重点信息一眼可辨：

\begin{itemize}
  \item \textbf{加粗文字}：核心参数、关键操作、质控红线、必选要求
  \item 【方头括号】：模板替换位、场景适配调整说明
  \item \begin{warningbox}[title={}]引用块格式：避坑指南、报错解决方案、风险提示\end{warningbox}
  \item \begin{tipbox}[title={}]直角引号：科研进阶技巧、优化方案\end{tipbox}
  \item \texttt{等宽字体}：代码命令、文件路径、软件操作路径
\end{itemize}

\subsection*{软件版本固定说明}

本手册所有操作步骤严格对应以下固定软件版本，版本更迭可能导致界面或参数变化：

\begin{center}
\begin{tabular}{llll}
\toprule
\textbf{软件} & \textbf{固定版本} & \textbf{官方下载} & \textbf{用途} \\
\midrule
Cytoscape & 3.10.1 & cytoscape.org & 网络构建与可视化 \\
AutoDock Vina & 1.2.3 & autodock.scripps.edu & 分子对接 \\
R & 4.3.3 & cran.r-project.org & 生信分析 \\
RDKit & 2023.09.1 & rdkit.org & 化学信息学 \\
PyMOL & 2.5 & pymol.org & 结合模式可视化 \\
GraphPad Prism & 9.0 & graphpad.com & 统计分析与绘图 \\
\bottomrule
\end{tabular}
\end{center}

\clearpage

% --- Table of Contents ---
\tableofcontents
\clearpage

\listoffigures
\listoftables
\clearpage

\pagenumbering{arabic}
\setcounter{page}{1}

% ============================================================
% GRAPHICAL ABSTRACT
% ============================================================
\section*{图形摘要 (Graphical Abstract)}
\addcontentsline{toc}{section}{图形摘要}

\begin{figure}[H]
  \centering
  \includegraphics[width=\textwidth]{../figures/graphical_abstract_v2.png}
  \caption{中药多维数据分析筛选免疫激活抗肿瘤药物全流程操作手册（Version 3.0）——5 阶段核心流程总览图。
  本图展示从 70 味中药候选库出发，经数据整合（TCMSP/BATMAN-TCM，OB $\geq$ 30\%/DL $\geq$ 0.18 ADME 筛选）、虚拟筛选（Cytoscape 网络药理学 + AutoDock Vina 分子对接，阈值 $\leq$ $-$8.0 kcal/mol），至体外验证（双臂设计：T 细胞系统——Jurkat 技能培训→小鼠脾脏 T 细胞主验证→人 PBMC 临床转化；共培养系统——Transwell 分离直接杀伤与免疫介导杀伤），再经分子机制验证（qRT-PCR/WB/CCK-8 IC50），最终成果转化（SCI 论文/专利/课题申报）的完整五阶段科研闭环。
  核心药效终点为「\textbf{直接抗肿瘤活性}（肿瘤细胞增殖抑制、凋亡促进）」与「\textbf{免疫激活活性}（T 细胞活化、巨噬细胞 M1 极化、细胞因子分泌上调）」双终点评价。
  \textit{Graphical abstract (v2.0) of the full-process actionable operation manual for screening immune-activating anti-tumor drugs from TCM multi-dimensional data, featuring the dual-arm T-cell validation system and Transwell co-culture design.}}
  \label{fig:graphical_abstract}
\end{figure}

\clearpage

% ============================================================
% MODULE 1: 研究前期准备与基础规范
% ============================================================
\section{模块一：研究前期准备与基础规范}
\label{sec:module1}

\begin{infobox}{模块一概览}
\textbf{操作目的：}建立标准化研究框架，明确研究问题、伦理边界、工具体系与预实验方案，确保后续所有操作具备可溯源性与合规性。\\
\textbf{本模块输出：}技术路线图（PDF）、伦理审批文件清单、工具配置确认表、预实验方案文档。\\
\textbf{预计耗时：}2--4 周（含伦理审批）。
\end{infobox}

% --- 1.1 研究方案设计 ---
\subsection{研究方案设计标准化流程}
\label{sec:module1_design}

\subsubsection{操作目的}

在开展任何数据库检索与实验操作前，需先建立明确的科学问题与标准化研究框架，避免因研究目标模糊导致后期数据分析方向混乱、实验返工。

\subsubsection{分步操作流程}

\begin{enumerate}
  \item \textbf{确立科学问题（Scientific Question Formulation）}

    采用 PICO（Population，Intervention，Comparison，Outcome）框架规范化表述研究问题：
    \begin{itemize}
      \item \textbf{P（研究对象）：}【模板替换位：鼻咽癌（Nasopharyngeal Carcinoma, NPC）患者 / 瘤细胞株；核心细胞系 CNE1、CNE2】
      \item \textbf{I（干预因素）：}中药【模板替换位：单味药名/复方名/活性单体名】的多维数据分析筛选结果
      \item \textbf{C（对照设置）：}阳性对照药（顺铂 cisplatin / 5-氟尿嘧啶 5-FU）；空白溶剂对照（DMSO $\leq$ 0.1\%）
      \item \textbf{O（研究终点）：}①直接抗肿瘤活性（IC50、细胞凋亡率、周期阻滞）；②免疫激活活性（IFN-$\gamma$ 分泌、CD8\textsuperscript{+} T 细胞活化、M1 巨噬细胞极化）
    \end{itemize}

    \textbf{标准化科学问题表述示例：}
    \begin{greenbox}
    「基于多维数据分析（网络药理学 + 分子对接 + 多组学验证）方法，从中药【XX】中筛选能够同时激活免疫系统并直接杀伤鼻咽癌（CNE1/CNE2）细胞的活性成分，并在体外验证其双重药效机制。」
    \end{greenbox}

  \item \textbf{绘制技术路线图（Technical Roadmap）}

    需执行以下操作绘制标准化技术路线图：
    \begin{enumerate}
      \item 打开 \software{Microsoft PowerPoint 2021} 或 \software{Xmind Zen 2024}，新建空白画布。
      \item 按以下 9 阶段绘制矩形流程框，每框内注明「阶段名称 → 核心工具 → 主要输出」：
        \begin{center}
        \small
        \begin{tabular}{lll}
        \toprule
        \textbf{阶段} & \textbf{核心工具} & \textbf{主要输出} \\
        \midrule
        Phase 1: 中药数据库构建 & TCMSP、ETCM & 成分-靶点映射表 \\
        Phase 2: 双维度靶点集筛选 & GeneCards、GEO、ImmPort & 靶点交集 Venn 图 \\
        Phase 3: 网络药理学分析 & Cytoscape 3.10.1、STRING & PPI 网络、核心靶点列表 \\
        Phase 4: 分子对接 & AutoDock Vina 1.2.3、PyMOL 2.5 & 对接评分排序表 \\
        Phase 5: 多组学验证 & R 4.3.3、TCGA、GEO & 差异表达 + 预后分析 \\
        Phase 6: 抗肿瘤湿实验 & CCK-8、流式细胞仪 & IC50、凋亡率 \\
        Phase 7: 免疫激活实验 & ELISA、流式细胞仪 & 细胞因子浓度、免疫标志物 \\
        Phase 8: 机制验证 & qRT-PCR、Western Blot & 靶点 mRNA/蛋白表达 \\
        Phase 9: 成果转化 & -- & 论文/专利/课题申报 \\
        \bottomrule
        \end{tabular}
        \end{center}
      \item 以带箭头的实线连接各阶段框，标注每个箭头旁的「输入文件」与「输出文件」名称。
      \item 导出为 A3 尺寸 PDF，存档于项目文件夹 \texttt{/docs/tech\_roadmap\_v1.pdf}。
    \end{enumerate}

  \item \textbf{制定纳入与排除标准（Inclusion \& Exclusion Criteria）}

    \begin{itemize}
      \item \textbf{中药纳入标准（任一满足）：}
        \begin{enumerate}
          \item 收录于《中国药典》2025 版（ChP 2025）且药用部位有效成分数据可查
          \item 已有同行评审文献（PubMed 检索，近 10 年内）报告对【模板替换位：鼻咽癌/目标瘤种】或免疫功能的调节活性
          \item 在 \database{TCMSP}（tcmsp-e.com）中可检索到 $\geq$ 3 种活性成分记录
        \end{enumerate}
      \item \textbf{中药排除标准（任一满足即排除）：}
        \begin{enumerate}
          \item 含有毒性成分（如马兜铃酸、雷公藤碱）且无脱毒处理记录
          \item ChP 2025 未收录且国内临床使用无文献依据
          \item 成分数据库收录成分 $<$ 3 种，无法进行有效网络药理学分析
        \end{enumerate}
      \item \textbf{活性成分纳入标准：}口服生物利用度 \textbf{OB $\geq$ 30\%} 且类药性 \textbf{DL $\geq$ 0.18}（TCMSP 默认金标准阈值 \cite{ru2014tcmsp}）
    \end{itemize}

  \item \textbf{设定研究终点（Endpoint Definition）}

    本研究设定双重主要终点：
    \begin{itemize}
      \item \textbf{主要终点 A（抗肿瘤终点）：}候选成分对 CNE1/CNE2 细胞的半数抑制浓度 \textbf{IC50 $\leq$ 50 $\mu$mol/L}（参照文献金标准阈值）
      \item \textbf{主要终点 B（免疫激活终点）：}候选成分处理后，THP-1 巨噬细胞 M1 极化比例（CD80\textsuperscript{+}CD86\textsuperscript{+}\%）或 Jurkat T 细胞 IL-2 分泌量较对照组升高 \textbf{$\geq$ 20\%}（P $<$ 0.05）
      \item \textbf{次要终点：}分子对接结合自由能 \textbf{$\leq$ $-$7.0 kcal/mol}；核心靶点在 TCGA/GEO 公共数据集中差异表达显著（$|$log2FC$|$ $\geq$ 1，FDR $<$ 0.05）
    \end{itemize}
\end{enumerate}

\begin{warningbox}
\begin{itemize}
  \item 不可在未明确科学问题前直接进行数据库检索，否则后期无法确定靶点交集的逻辑依据，导致整个分析框架缺乏科学性。
  \item 技术路线图必须在开始任何操作前完成，并经导师签字确认后方可执行，否则论文答辩或课题申报时可能因研究方向模糊遭到否决。
  \item IC50 阈值需根据正式纳入的对照药（顺铂在 CNE1 细胞中的 IC50 约为 2--8 $\mu$mol/L）进行相对评估，不应将 50 $\mu$mol/L 视为绝对通过标准。
\end{itemize}
\end{warningbox}

% --- 1.2 伦理与合规规范 ---
\subsection{伦理与合规规范}
\label{sec:module1_ethics}

\subsubsection{中药基原鉴定规范}

需执行以下操作进行中药基原鉴定与规范化登记：
\begin{enumerate}
  \item 按《中国药典》2025 版（第一部）核对中药的中文学名、拉丁学名、药用部位与产地规范。
  \item 对于复方中药，逐一列出每味药的基原鉴定信息，建立「中药基原信息登记表」（格式：中文名 / 拉丁学名 / 科属 / 药用部位 / 产地 / ChP 2025 收录页码 / 批号）。
  \item 实验用中药材样品需附权威鉴定机构（如中国中医科学院中药研究所）出具的鉴定报告，或使用符合 ChP 2025 标准的标准品（NIFDC 批号）。
\end{enumerate}

\subsubsection{细胞实验伦理要求}

CNE1/CNE2 为商业细胞系，不涉及人体研究伦理，但须满足：
\begin{itemize}
  \item 细胞系来源合法（ATCC / DSMZ / 中国典型培养物保藏中心 CCTCC），保留购买合同与 Certificate of Analysis (CoA)。
  \item 细胞实验在 \textbf{BSL-2（生物安全二级）}实验室内操作，持有效生物安全培训证书。
  \item \textbf{细胞系 STR 鉴定：}新复苏细胞系须在传代 5 代内完成 STR 鉴定，鉴定结果存档。
\end{itemize}

\begin{warningbox}
CNE1 和 CNE2 细胞系已被国际细胞系认证与鉴定委员会（ICLAC）确认存在 HeLa 细胞污染（ICLAC-00596）。在 2015 年后发表的 SCI 论文中，部分高影响力期刊（如 \textit{Cancer Research}、\textit{Oncogene}）要求作者在使用已知污染细胞系时提供额外说明或替换为认证细胞系（推荐使用：HK-1、6-10B、HNE2 或 C666-1）。研究者应事先查阅目标期刊的细胞系政策。
\end{warningbox}

\subsubsection{数据使用合规性}

\begin{itemize}
  \item \database{TCGA} 数据：用于学术研究免费，商业使用须遵循 TCGA 数据访问政策（dbGaP 批准）。
  \item \database{GEO} 数据集：下载前确认数据集许可证类型（CC BY 或 CC0），在论文中引用原始发表文献。
  \item \database{PDB} 蛋白质结构：以 CC0 公开授权，可自由下载使用，引用时需注明 PDB ID 与原始论文。
\end{itemize}

\subsubsection{学术不端规避要点}

\begin{enumerate}
  \item \textbf{数据透明化：}所有分析代码须上传 GitHub 或 Zenodo，保留版本控制记录（git log）。
  \item \textbf{图片规范处理：}禁止对 Western Blot 条带进行拼接、亮度对比度过度调整，图片修改须记录修改日志。
  \item \textbf{重复性验证：}所有关键实验须进行 $\geq$ 3 次独立生物学重复（biological replicates），并标注「独立实验」与「技术重复」的区别。
  \item \textbf{AI 辅助声明：}若使用 ChatGPT、文心一言等 AI 工具辅助写作，须在投稿时按目标期刊要求进行声明（参考 ICMJE 2024 指南）。
\end{enumerate}

% --- 1.3 必备工具清单 ---
\subsection{必备工具与数据库清单}
\label{sec:module1_tools}

\subsubsection{第一类：中药数据库}

\begin{longtable}{p{3cm}p{3.5cm}p{2cm}p{2cm}p{4cm}}
\caption{核心中药数据库清单} \label{tab:tcm_databases} \\
\toprule
\textbf{数据库} & \textbf{官方 URL} & \textbf{收费} & \textbf{版本} & \textbf{核心用途} \\
\midrule
\endfirsthead
\toprule
\textbf{数据库} & \textbf{官方 URL} & \textbf{收费} & \textbf{版本} & \textbf{核心用途} \\
\midrule
\endhead
\midrule
\multicolumn{5}{r}{\small 续下页} \\
\endfoot
\bottomrule
\endlastfoot
TCMSP & tcmsp-e.com & 免费 & 2.3 & 活性成分 ADME 筛选（OB/DL 阈值标准来源）\cite{ru2014tcmsp} \\
\addlinespace
ETCM & etcm.alihealth.cn & 免费 & v2.0 & 中药-成分-靶点-表型网络，含毒性数据 \\
\addlinespace
BATMAN-TCM & bionet.ncpsb.org.cn/batman-tcm & 免费 & v2.0 & 成分-靶点相互作用预测，置信度评分 \cite{liu2024batman} \\
\addlinespace
TCMID & bidd.nus.edu.sg/tcmid & 免费 & v2.01 & 中药方剂-成分-药效数据库 \\
\addlinespace
HIT 2.0 & hit2.badd-cadd.cn & 免费 & v2.0 & 草药抑制剂靶点数据库 \\
\addlinespace
SymMap & symmap.org & 免费 & v2 & 中药症状-靶点-疾病关联 \\
\addlinespace
HERB & herb.ac.cn & 免费 & v1.0 & 整合中药-成分-疾病-靶点全数据 \\
\end{longtable}

\subsubsection{第二类：肿瘤与免疫靶点数据库}

\begin{longtable}{p{3cm}p{3.5cm}p{2cm}p{4.5cm}}
\caption{核心肿瘤与免疫靶点数据库清单} \label{tab:target_databases} \\
\toprule
\textbf{数据库} & \textbf{官方 URL} & \textbf{收费} & \textbf{核心用途} \\
\midrule
\endfirsthead
\toprule
\textbf{数据库} & \textbf{官方 URL} & \textbf{收费} & \textbf{核心用途} \\
\midrule
\endhead
\midrule
\multicolumn{4}{r}{\small 续下页} \\
\endfoot
\bottomrule
\endlastfoot
GeneCards & genecards.org & 免费(学术) & 疾病-基因关联，含综合得分排序 \\
\addlinespace
OMIM & omim.org & 免费 & 遗传病与基因关联标准数据库 \\
\addlinespace
DisGeNET & disgenet.org & 免费 & 疾病-基因关联，含证据分级 \\
\addlinespace
DrugBank & drugbank.ca & 免费(学术) & 药物-靶点-药理学全数据 \\
\addlinespace
ImmPort & immport.org & 免费 & 免疫相关基因标准基因集（核心免疫激活靶点来源）\\
\addlinespace
InnateDB & innatedb.com & 免费 & 先天免疫调控基因与通路 \\
\addlinespace
TIMER 2.0 & timer.cistrome.org & 免费 & 肿瘤免疫浸润分析 \\
\addlinespace
TCGA & portal.gdc.cancer.gov & 免费 & 泛癌基因组数据，NPC 代用 HNSC 数据集 \\
\addlinespace
GEO & ncbi.nlm.nih.gov/geo & 免费 & NPC 专属数据集（GSE12452、GSE53819、GSE61218）\\
\end{longtable}

\subsubsection{第三类：生信分析工具}

\begin{longtable}{p{3cm}p{2.5cm}p{2cm}p{6cm}}
\caption{核心生信分析工具清单} \label{tab:bioinformatics_tools} \\
\toprule
\textbf{工具} & \textbf{版本} & \textbf{收费} & \textbf{安装与核心用途} \\
\midrule
\endfirsthead
\toprule
\textbf{工具} & \textbf{版本} & \textbf{收费} & \textbf{安装与核心用途} \\
\midrule
\endhead
\bottomrule
\endlastfoot
R & 4.3.3 & 免费 & cran.r-project.org；生信核心计算平台 \\
\addlinespace
RStudio & 2024.04 & 免费 & posit.co；R 语言 IDE \\
\addlinespace
Bioconductor & 3.18 & 免费 & 随 R 安装，提供 limma、clusterProfiler 等 \\
\addlinespace
STRING & online / app & 免费 & string-db.org；PPI 网络构建 \cite{szklarczyk2021string} \\
\addlinespace
DAVID & online & 免费 & david.ncifcrf.gov；GO/KEGG 富集分析 \\
\addlinespace
Metascape & online & 免费 & metascape.org；综合富集分析 \\
\addlinespace
Python & 3.11 & 免费 & python.org；数据处理、批量对接脚本 \\
\addlinespace
RDKit & 2023.09.1 & 免费 & rdkit.org；化学信息学处理 \\
\end{longtable}

\subsubsection{第四类：分子对接软件}

\begin{longtable}{p{3cm}p{2.5cm}p{2cm}p{6cm}}
\caption{分子对接软件清单} \label{tab:docking_tools} \\
\toprule
\textbf{软件} & \textbf{版本} & \textbf{收费} & \textbf{安装与核心用途} \\
\midrule
\endfirsthead
\toprule
\textbf{软件} & \textbf{版本} & \textbf{收费} & \textbf{安装与核心用途} \\
\midrule
\endhead
\bottomrule
\endlastfoot
AutoDock Vina & 1.2.3 & 免费 & autodock.scripps.edu；主要对接引擎 \cite{eberhardt2021autodock4} \\
\addlinespace
AutoDockTools & 1.5.7 & 免费 & MGLTools；受体预处理与格点设置 \\
\addlinespace
PyMOL & 2.5 & 付费/学术免费 & pymol.org；结合模式可视化 \\
\addlinespace
UCSF ChimeraX & 1.7 & 免费 & rbvi.ucsf.edu/chimerax；受体预处理备选 \\
\addlinespace
OpenBabel & 3.1.1 & 免费 & openbabel.org；配体格式转换 \\
\addlinespace
PubChem & online & 免费 & pubchem.ncbi.nlm.nih.gov；配体 3D 结构下载 \\
\addlinespace
RCSB PDB & online & 免费 & rcsb.org；受体晶体结构下载 \\
\end{longtable}

% --- 1.4 预实验准备规范 ---
\subsection{预实验准备规范}
\label{sec:module1_pilot}

\subsubsection{试剂配制标准操作}

\begin{enumerate}
  \item \textbf{DMSO 储液配制（中药成分溶液）：}
    \begin{enumerate}
      \item 精确称量中药活性成分（电子天平，精度 0.01 mg），记录批次与纯度（$\geq$ 98\%，HPLC 级）。
      \item 溶于分析纯 DMSO（sigma-aldrich，目录号 276855），配制 \textbf{10 mmol/L} 储液。
      \item 用涡旋仪（1000 rpm，3 min）充分溶解，分装为 200 $\mu$L/管，储存于 $-$80$^{\circ}$C。
      \item \textbf{工作浓度配制：}使用前从 $-$80$^{\circ}$C 取出储液，室温解冻 5 min，用完全培养基梯度稀释至所需工作浓度，确保 \textbf{DMSO 终浓度 $\leq$ 0.1\%}（v/v）。
    \end{enumerate}

  \item \textbf{仪器校准清单（实验开始前 1 周完成）：}
    \begin{itemize}
      \item 超净工作台：开机紫外灭菌 $\geq$ 30 min，酒精灯火焰灭菌操作台面
      \item CO$_2$ 培养箱：校准温度 37.0 $\pm$ 0.5$^{\circ}$C，CO$_2$ 浓度 5.0 $\pm$ 0.1\%
      \item 酶标仪（ELISA）：校准空白本底值与线性范围（OD 值 0.1--3.0）
      \item 流式细胞仪：以标准荧光微球（检测当天）执行 QC 质控，确认 CV $<$ 5\%
      \item 离心机：校准转速（rpm）与实际 g 值对应关系
    \end{itemize}
\end{enumerate}

\begin{qcbox}
\textbf{模块一质控合格标准：}
\begin{itemize}
  \item 技术路线图已导师审核签字 \ding{52}
  \item 中药基原信息登记表已完成，每味药有 ChP 2025 页码记录 \ding{52}
  \item 细胞系 STR 鉴定结果已存档（图片 + 数据文件）\ding{52}
  \item 所有数据库账号已注册并测试可正常登录 \ding{52}
  \item 软件版本已安装确认（Cytoscape 3.10.1、AutoDock Vina 1.2.3、R 4.3.3）\ding{52}
  \item DMSO 储液已配制，浓度验证无误 \ding{52}
\end{itemize}
\end{qcbox}

\moduleoutput{技术路线图 PDF、工具配置清单（Excel）、中药基原信息登记表、伦理合规文件夹}

% ============================================================
% MODULE 2: 中药多维基础数据库构建全流程
% ============================================================
\section{模块二：中药多维基础数据库构建全流程}
\label{sec:module2}

\begin{infobox}{模块二概览}
\textbf{操作目的：}构建标准化的中药化学成分多维数据库，为后续靶点预测与虚拟筛选提供高质量输入数据。\\
\textbf{前置输入：}模块一输出的中药基原信息登记表、纳入/排除标准。\\
\textbf{本模块输出：}中药-成分映射表（Excel）、ADME/T 筛选后活性成分列表（含 SMILES）、去重标准化成分数据库（CSV）。\\
\textbf{预计耗时：}1--2 周。
\end{infobox}

\subsection{中药入选与排除标准制定}
\label{sec:module2_selection}

执行以下操作建立中药筛选标准：

\begin{enumerate}
  \item \textbf{中药纳入规则制定（需同时满足所有条件）：}
    \begin{enumerate}
      \item \textbf{基原规范性：}收录于《中国药典》2025 版第一部，中文学名与拉丁学名可在 TCMSP 数据库检索到对应记录。
      \item \textbf{现代药理学基础：}在 PubMed 检索式为「[中药名] AND (nasopharyngeal carcinoma OR immune)」中，检索到 $\geq$ 2 篇近 10 年内发表的相关研究文献。
      \item \textbf{数据充足性：}在 TCMSP 数据库中，该中药的化学成分记录 $\geq$ 3 种（以确保后续网络构建的数据量）。
    \end{enumerate}

  \item \textbf{核心排除标准（任一满足即排除）：}
    \begin{enumerate}
      \item 含已确认肝肾毒性成分（如马兜铃酸、汉防己甲素）且无脱毒处理依据
      \item TCMSP 成分记录 $<$ 3 种
      \item ChP 2025 未收录，且缺乏来源可追溯的规范质量标准
    \end{enumerate}

  \item \textbf{示例——鼻咽癌核心候选中药参考列表：}

    \begin{templatebox}
    以下为【鼻咽癌场景】核心候选中药参考，研究者可根据自身研究聚焦点选择 3--5 味：
    \begin{itemize}
      \item 连翘（\textit{Forsythia suspensa}）：含连翘苷、连翘酯苷 A，已有 NPC 抗肿瘤活性报告
      \item 黄芩（\textit{Scutellaria baicalensis}）：含黄芩苷、汉黄芩素，免疫调节活性明确
      \item 半枝莲（\textit{Scutellaria barbata}）：野黄芩苷等，已报告 NPC 细胞毒性
      \item 白花蛇舌草（\textit{Hedyotis diffusa}）：车叶草苷酸等，NK 细胞激活活性
      \item 人参（\textit{Panax ginseng}）：人参皂苷 Rg3，免疫增强 + 抗肿瘤双效
    \end{itemize}
    【其他瘤种：替换为对应瘤种有活性记录的中药，保持上述筛选逻辑不变】
    \end{templatebox}
\end{enumerate}

\subsection{中药化学成分多维数据采集标准化操作}
\label{sec:module2_data_collection}

\subsubsection{步骤一：TCMSP 数据库检索操作}

\begin{enumerate}
  \item 打开浏览器，访问 \texttt{https://tcmsp-e.com}。
  \item 在主页搜索框中输入中药的中文名称（如「连翘」），点击「Search」。
  \item 在检索结果页面，点击「Herb Name」列中对应条目的链接，进入该药材详情页。
  \item 点击「Related Ingredients」标签页，获取该中药所有化学成分列表。
  \item 点击列表右上角「Download」按钮，选择「CSV」格式下载，保存为 \texttt{[药材名]\_ingredients\_raw.csv}。
  \item 对每一味中药重复上述步骤，分别下载并保存。
\end{enumerate}

\begin{warningbox}
TCMSP 网站偶发访问不稳定，建议同时注册 \database{HERB}（herb.ac.cn）作为备用数据库。HERB 数据库收录范围更广，数据更新至 2021 年，且提供 API 接口，适合批量检索。
\end{warningbox}

\subsubsection{步骤二：OB 和 DL 参数筛选操作}

\begin{enumerate}
  \item 打开 R 4.3.3，新建 R 脚本文件 \texttt{02\_adme\_filter.R}。
  \item 执行以下代码完成 ADME/T 筛选：
\end{enumerate}

\begin{lstlisting}[language=R,caption={ADME/T 筛选核心代码}]
# ==== ADME/T 筛选脚本 ====
# 作者：基于 TCMSP 标准流程
# 参考：Ru et al., J Cheminformatics, 2014

library(readr)
library(dplyr)
library(openxlsx)

# 1. 读取原始成分数据
ingredients_raw <- read_csv("data/[药材名]_ingredients_raw.csv")

# 2. 查看列名（确认 OB 和 DL 列名）
colnames(ingredients_raw)
# 预期列名包括：Mol.ID, Molecule.Name, OB (%), DL

# 3. 执行 ADME 筛选（金标准阈值）
# OB >= 30% AND DL >= 0.18
ingredients_filtered <- ingredients_raw %>%
  filter(`OB (%)` >= 30,       # 口服生物利用度阈值
         `DL` >= 0.18)          # 类药性阈值

# 4. 质控检查：筛选后成分数量
cat("原始成分数量:", nrow(ingredients_raw), "\n")
cat("筛选后成分数量:", nrow(ingredients_filtered), "\n")
cat("筛选比例:", round(nrow(ingredients_filtered)/nrow(ingredients_raw)*100, 1), "%\n")

# 5. 保存筛选结果
write_csv(ingredients_filtered, "results/02_adme_filtered_compounds.csv")
write.xlsx(ingredients_filtered, "results/02_adme_filtered_compounds.xlsx",
           rowNames = FALSE)

cat("ADME 筛选完成，结果已保存至 results/ 目录\n")
\end{lstlisting}

\begin{tipbox}
「筛选阈值说明：\textbf{OB $\geq$ 30\%} 代表该成分在口服给药后有足够的血液吸收率，确保体内能达到有效浓度；\textbf{DL $\geq$ 0.18} 参照 Lipinski 类药规则（MW $<$ 500 Da，H 键供体 $\leq$ 5，H 键受体 $\leq$ 10，LogP $<$ 5）量化成药可能性。这两个阈值由 TCMSP 开发者 Ru 等（2014）基于中药特性校正后提出，已被超过 3,000 篇网络药理学论文采用，为当前领域的金标准。」
\end{tipbox}

\subsubsection{步骤三：多数据库整合与去重标准化}

\begin{enumerate}
  \item \textbf{补充检索 ETCM、BATMAN-TCM：}
    \begin{enumerate}
      \item 访问 \texttt{https://etcm.alihealth.cn}，使用中药名检索，下载「Ingredients」数据（CSV 格式）。
      \item 访问 \texttt{https://bionet.ncpsb.org.cn/batman-tcm}，输入中药名，下载「Ingredients and Targets」数据。
      \item 将三个数据库的成分列表导入 Excel，分别在三个 Sheet 中保存，以 Molecule Name / CAS / InChIKey 为合并键。
    \end{enumerate}

  \item \textbf{成分去重标准化操作：}
\end{enumerate}

\begin{lstlisting}[language=R,caption={成分去重与标准化代码}]
# ==== 成分去重与标准化 ====
library(dplyr)
library(readr)

# 1. 读取各数据库筛选结果
tcmsp_data  <- read_csv("results/02_tcmsp_filtered.csv")
etcm_data   <- read_csv("results/02_etcm_compounds.csv")
batman_data <- read_csv("results/02_batman_compounds.csv")

# 2. 统一列名格式
# 确保每个数据框至少包含：compound_name, smiles, mol_id
# 各数据库列名需手动对应调整

# 3. 合并数据集
all_compounds <- bind_rows(
  tcmsp_data  %>% mutate(source = "TCMSP"),
  etcm_data   %>% mutate(source = "ETCM"),
  batman_data %>% mutate(source = "BATMAN-TCM")
)

# 4. 基于 InChIKey 去重（优先），备选 CAS 或化合物名
# 方法一：若有 InChIKey 列
if ("InChIKey" %in% colnames(all_compounds)) {
  all_compounds_dedup <- all_compounds %>%
    filter(!is.na(InChIKey)) %>%
    distinct(InChIKey, .keep_all = TRUE)
} else {
  # 方法二：基于化合物名去重（大小写标准化）
  all_compounds_dedup <- all_compounds %>%
    mutate(compound_name_std = tolower(trimws(compound_name))) %>%
    distinct(compound_name_std, .keep_all = TRUE)
}

# 5. 输出去重统计
cat("合并前总成分数:", nrow(all_compounds), "\n")
cat("去重后成分数:", nrow(all_compounds_dedup), "\n")
cat("去重率:", round((1 - nrow(all_compounds_dedup)/nrow(all_compounds))*100, 1), "%\n")

# 6. 保存标准化成分库
write_csv(all_compounds_dedup, "results/03_compounds_db_standardized.csv")
\end{lstlisting}

\begin{qcbox}
\textbf{模块二质控合格标准：}
\begin{itemize}
  \item TCMSP/ETCM/BATMAN-TCM 三库检索均已执行，记录每库检索结果数量 \ding{52}
  \item ADME 筛选采用固定阈值 \textbf{OB $\geq$ 30\% AND DL $\geq$ 0.18}，阈值来源已标注 \ding{52}
  \item 去重后活性成分数量 $\geq$ 20 种（少于 20 种需检查中药选择是否过少，或降低 DL 阈值至 0.15 并在论文中说明）\ding{52}
  \item 每种活性成分均有 SMILES 字符串（用于后续分子对接配体文件生成）\ding{52}
  \item 最终成分数据库（CSV 格式）已备份至 Git 仓库 \ding{52}
\end{itemize}
\end{qcbox}

\subsection{中药-化学成分对应关系库搭建}
\label{sec:module2_mapping}

执行以下操作建立「中药-成分」一对多映射关系库：

\begin{enumerate}
  \item 新建 Excel 文件 \texttt{04\_herb\_compound\_mapping.xlsx}，设置以下列：
    \begin{center}
    \begin{tabular}{llll}
    \toprule
    \textbf{列名} & \textbf{数据类型} & \textbf{示例值} & \textbf{说明} \\
    \midrule
    herb\_id & 字符串 & H001 & 中药内部编号 \\
    herb\_name\_cn & 字符串 & 连翘 & 中文名 \\
    herb\_name\_lat & 字符串 & \textit{Forsythia suspensa} & 拉丁学名 \\
    compound\_id & 字符串 & MOL001 & 成分内部编号 \\
    compound\_name & 字符串 & 连翘苷 & 成分名称 \\
    smiles & 字符串 & C...O & SMILES 结构式 \\
    ob\_percent & 数值 & 45.6 & 口服生物利用度 (\%) \\
    dl & 数值 & 0.21 & 类药性 \\
    source\_db & 字符串 & TCMSP & 数据来源 \\
    \bottomrule
    \end{tabular}
    \end{center}

  \item 对于复方中药研究，每味药的成分均单独列行，herb\_id 字段区分药材来源，实现多对多关系的规范化存储。

  \item 将 Excel 文件另存为 CSV 格式（UTF-8 编码），作为模块三靶点预测的输入文件。
\end{enumerate}

\begin{warningbox}
\begin{itemize}
  \item SMILES 字符串的正确性直接影响后续分子对接的结果，须使用 \database{PubChem Identifier Exchange}（pubchem.ncbi.nlm.nih.gov/idexchange）进行 SMILES 验证，避免使用 TCMSP 中个别成分存在的错误 SMILES 记录（已知数据库错误率约 2--5\%）。
  \item 若某成分无法检索到 SMILES，禁止手动编写 SMILES，应在备注列标注「无结构数据-排除」并从后续分析中移除。
\end{itemize}
\end{warningbox}

\moduleoutput{标准化中药-成分映射表（CSV/Excel）：herb\_id、compound\_id、SMILES、OB、DL，共 [N] 种活性成分，用于模块三的靶点预测输入}

% ============================================================
% MODULE 3: 免疫激活+抗肿瘤双维度核心靶点集构建
% ============================================================
\section{模块三：免疫激活 + 抗肿瘤双维度核心靶点集构建}
\label{sec:module3}

\begin{infobox}{模块三概览}
\textbf{操作目的：}分别构建「抗肿瘤靶点集」与「免疫激活靶点集」，通过 Venn 图交集分析获取双效候选靶点集，作为网络药理学与虚拟筛选的核心靶点库。\\
\textbf{前置输入：}模块二输出的成分-靶点预测结果（SwissTargetPrediction）、中药活性成分列表。\\
\textbf{本模块输出：}抗肿瘤靶点集（TXT）、免疫激活靶点集（TXT）、双效靶点交集（Excel）、Venn 图（PDF）。\\
\textbf{预计耗时：}1 周。
\end{infobox}

\subsection{抗肿瘤核心靶点数据集构建操作}
\label{sec:module3_tumor}

\subsubsection{步骤一：疾病靶点多库并行检索}

需对以下 5 个数据库执行检索，分别下载鼻咽癌（NPC）相关靶点：

\begin{enumerate}
  \item \textbf{GeneCards 检索操作：}
    \begin{enumerate}
      \item 访问 \texttt{https://www.genecards.org}。
      \item 在搜索框输入 「nasopharyngeal carcinoma」，点击搜索。
      \item 在结果页面点击「Genes」标签，查看基因关联列表。
      \item 按「GeneCards Score」降序排列，下载前 \textbf{500} 个基因（点击「Download」，选择「Excel」或「CSV」）。
      \item \textbf{重要参数：}Relevance Score 阈值设置为 $\geq$ 2.0（过滤低相关性基因）。
      \item 保存为 \texttt{03\_genecards\_NPC\_targets.csv}。
    \end{enumerate}

  \item \textbf{DisGeNET 检索操作：}
    \begin{enumerate}
      \item 访问 \texttt{https://www.disgenet.org}，注册免费账号并登录。
      \item 点击「Search」→「Diseases」，搜索「nasopharyngeal carcinoma」。
      \item 选择对应疾病条目（Disease Identifier: C0006413 / DOID:9261）。
      \item 进入疾病详情页，点击「Genes」标签，设置 \textbf{Score $\geq$ 0.1} 过滤。
      \item 下载所有关联基因（CSV 格式），保存为 \texttt{03\_disgenet\_NPC\_targets.csv}。
    \end{enumerate}

  \item \textbf{OMIM 检索操作：}
    \begin{enumerate}
      \item 访问 \texttt{https://www.omim.org}，搜索「nasopharyngeal carcinoma」。
      \item 获取 OMIM 条目编号（如 607107），进入详情页提取「Gene」列表。
      \item 手动或通过 API 提取相关基因，保存为 \texttt{03\_omim\_NPC\_targets.txt}。
    \end{enumerate}
\end{enumerate}

\subsubsection{步骤二：GEO 数据集差异基因分析（鼻咽癌核心差异基因筛选）}

\begin{warningbox}
\textbf{NPC 与 TCGA 数据库的重要说明：}TCGA 数据库的肿瘤类型中不包含独立的鼻咽癌（NPC）分类，NPC 数据需从以下三个 GEO 数据集获取（均为领域内公认的 NPC 标准数据集）：
\begin{itemize}
  \item \textbf{GSE12452}：NPC vs. 正常鼻咽黏膜，31 NPC + 10 正常，HG-U133A 芯片
  \item \textbf{GSE53819}：NPC vs. 正常，18 NPC + 18 正常，Agilent 芯片
  \item \textbf{GSE61218}：NPC vs. 正常，31 NPC + 10 正常，HG-U133A 芯片
\end{itemize}
三个数据集均可直接通过 GEO2R（NCBI 在线工具）或 R/GEOquery 包下载。
\end{warningbox}

\begin{enumerate}
  \item \textbf{GEOquery 下载数据集操作：}
\end{enumerate}

\begin{lstlisting}[language=R,caption={GEO 数据集下载与差异分析（limma 标准流程）}]
# ==== GEO + limma 差异基因分析 ====
# 参考：Ritchie et al., Nucleic Acids Research, 2015

# 安装必要包（首次运行时）
if (!require("BiocManager")) install.packages("BiocManager")
BiocManager::install(c("GEOquery", "limma", "affy", "hgu133a.db"))
library(GEOquery)
library(limma)

# 1. 下载 GEO 数据集 GSE12452
gse <- getGEO("GSE12452", GSEMatrix = TRUE, destdir = "data/GEO/")
expr_data <- exprs(gse[[1]])      # 表达矩阵
pheno_data <- pData(gse[[1]])     # 样本信息

# 2. 查看分组信息（确认 NPC vs. Normal 分组）
print(table(pheno_data$`characteristics_ch1`))

# 3. 定义实验分组
group <- factor(ifelse(grepl("NPC|tumor|carcinoma",
  pheno_data$`characteristics_ch1`, ignore.case = TRUE), "NPC", "Normal"))
design <- model.matrix(~0 + group)
colnames(design) <- levels(group)

# 4. 数据质控（箱线图检查）
boxplot(expr_data[,1:min(10, ncol(expr_data))],
        main = "GSE12452 Expression Distribution",
        las = 2, col = "lightblue")

# 5. limma 差异分析
fit <- lmFit(expr_data, design)
contrast_matrix <- makeContrasts(NPC - Normal, levels = design)
fit2 <- contrasts.fit(fit, contrast_matrix)
fit2 <- eBayes(fit2)

# 6. 提取差异基因（阈值设置）
DEGs <- topTable(fit2,
                 adjust.method = "BH",   # Benjamini-Hochberg 校正
                 p.value = 0.05,         # FDR < 0.05
                 lfc = 1,                # |log2FC| >= 1
                 number = Inf)           # 提取所有显著基因

cat("差异基因总数:", nrow(DEGs), "\n")
cat("上调基因:", sum(DEGs$logFC > 0), "\n")
cat("下调基因:", sum(DEGs$logFC < 0), "\n")

# 7. 保存差异基因列表
write.csv(DEGs, "results/03_DEGs_GSE12452_NPC_vs_Normal.csv")

# 8. 绘制火山图
with(DEGs, {
  plot(logFC, -log10(adj.P.Val),
       pch = 20, cex = 0.5,
       xlab = "log2 Fold Change",
       ylab = "-log10 Adjusted P-value",
       main = "GSE12452: NPC vs. Normal",
       col = ifelse(adj.P.Val < 0.05 & abs(logFC) > 1,
                    ifelse(logFC > 0, "red", "blue"), "gray"))
  abline(h = -log10(0.05), lty = 2, col = "darkgreen")
  abline(v = c(-1, 1), lty = 2, col = "orange")
})
\end{lstlisting}

\begin{itemize}
  \item 对 \textbf{GSE53819} 和 \textbf{GSE61218} 重复上述流程，分别生成差异基因列表。
  \item 取三个数据集差异基因的并集（union），保存为 \texttt{03\_NPC\_DEGs\_union.txt}（每行一个 Gene Symbol）。
\end{itemize}

\subsection{免疫激活核心靶点数据集构建操作}
\label{sec:module3_immune}

\subsubsection{步骤一：ImmPort 免疫相关基因集检索}

\begin{enumerate}
  \item 访问 \texttt{https://www.immport.org}，点击「Data Browser」。
  \item 选择「Gene Lists」→「Immunology Genes」，下载 ImmPort 提供的标准免疫基因集（格式：CSV，含基因名称与功能分类）。
  \item 重点提取以下免疫功能分类的基因（在下载文件中按「Category」列筛选）：
    \begin{itemize}
      \item T cell receptor signaling（T 细胞受体信号）
      \item B cell receptor signaling（B 细胞受体信号）
      \item Innate immune signaling（先天免疫信号）
      \item NK cell cytotoxicity（NK 细胞毒性）
      \item Cytokine/cytokine receptor（细胞因子及受体）
      \item Immune checkpoint（免疫检查点）
    \end{itemize}
  \item 保存为 \texttt{03\_immport\_immune\_genes.txt}。
\end{enumerate}

\subsubsection{步骤二：免疫检查点靶点标准基因集纳入}

需手动整合以下免疫检查点（Immune Checkpoint, ICP）相关核心靶点（基于当前已批准免疫治疗靶点）：

\begin{center}
\begin{tabular}{llll}
\toprule
\textbf{基因 Symbol} & \textbf{蛋白名称} & \textbf{功能类别} & \textbf{NPC 相关性} \\
\midrule
\gene{CD274} & PD-L1 & 抑制性检查点 & NPC 高表达，EBV-LMP1 驱动 \\
\gene{PDCD1} & PD-1 & 抑制性检查点 & T 细胞耗竭标志物 \\
\gene{CTLA4} & CTLA-4 & 抑制性检查点 & Treg 激活 \\
\gene{LAG3} & LAG-3 & 抑制性检查点 & 耗竭 T 细胞共表达 \\
\gene{HAVCR2} & TIM-3 & 抑制性检查点 & 晚期耗竭标志物 \\
\gene{CD80} & B7-1 & 共刺激分子 & DC 激活信号 \\
\gene{CD86} & B7-2 & 共刺激分子 & 巨噬细胞 M1 极化标志物 \\
\gene{IFNG} & IFN-$\gamma$ & 效应细胞因子 & Th1/CTL 核心效应因子 \\
\gene{TNF} & TNF-$\alpha$ & 效应细胞因子 & M1 极化与 NF-$\kappa$B 信号 \\
\gene{IL2} & IL-2 & T 细胞增殖因子 & CTL 激活核心因子 \\
\gene{IL6} & IL-6 & 炎症细胞因子 & 双向免疫调节 \\
\gene{TGFB1} & TGF-$\beta$1 & 免疫抑制因子 & 抑制 CTL 的核心通路 \\
\bottomrule
\end{tabular}
\end{center}

\subsubsection{步骤三：免疫激活靶点集汇总}

\begin{lstlisting}[language=R,caption={免疫靶点集整合代码}]
# ==== 免疫靶点集整合 ====
library(dplyr)
library(readr)

# 1. 读取 ImmPort 基因集
immport_genes <- read_csv("data/03_immport_immune_genes.txt",
                          col_names = c("gene_symbol", "category"))

# 2. 定义核心免疫检查点/激活基因（手动整合）
checkpoint_genes <- c("CD274", "PDCD1", "CTLA4", "LAG3", "HAVCR2",
                       "CD80", "CD86", "CD28", "ICOS", "CD40",
                       "IFNG", "TNF", "IL2", "IL6", "IL10",
                       "IL12A", "IL12B", "TGFB1", "CXCL9", "CXCL10",
                       "CCL5", "GZMB", "PRF1", "NOS2", "CD68",
                       "FOXP3", "SELL", "CD4", "CD8A", "CD19",
                       "NCAM1", "KLRB1", "PDCD1LG2", "IDO1", "VEGFA")

checkpoint_df <- data.frame(
  gene_symbol = checkpoint_genes,
  category = "Immune_Checkpoint_or_Activation"
)

# 3. 合并所有免疫靶点
immune_targets <- bind_rows(immport_genes, checkpoint_df) %>%
  distinct(gene_symbol, .keep_all = TRUE)

cat("免疫靶点集总数:", nrow(immune_targets), "\n")

# 4. 保存免疫靶点集
write_csv(immune_targets, "results/03_immune_target_set.csv")
\end{lstlisting}

\subsection{双维度靶点交集筛选操作（Venn 分析）}
\label{sec:module3_venn}

\begin{lstlisting}[language=R,caption={靶点交集分析与 Venn 图绘制}]
# ==== 双维度靶点 Venn 交集分析 ====
library(VennDiagram)
library(dplyr)
library(readr)

# 1. 读取各靶点集
npc_targets <- read_csv("results/03_NPC_DEGs_union.txt",
                        col_names = "gene_symbol")$gene_symbol
immune_targets <- read_csv("results/03_immune_target_set.csv")$gene_symbol

# 注：中药靶点预测结果（从 SwissTargetPrediction 获取）
# 访问 http://www.swisstargetprediction.ch，上传 SMILES，下载靶点结果
tcm_targets <- read_csv("results/03_tcm_predicted_targets.csv")$gene_symbol

# 2. 计算交集
# 核心交集：中药靶点 ∩ 抗肿瘤靶点（单靶点）
tcm_npc_intersection <- intersect(tcm_targets, npc_targets)
# 双维度交集：中药靶点 ∩ 抗肿瘤靶点 ∩ 免疫靶点
triple_intersection <- Reduce(intersect,
                               list(tcm_targets, npc_targets, immune_targets))

cat("中药靶点数量:", length(tcm_targets), "\n")
cat("NPC 靶点数量:", length(npc_targets), "\n")
cat("免疫靶点数量:", length(immune_targets), "\n")
cat("中药-NPC 交集:", length(tcm_npc_intersection), "\n")
cat("三维交集（双效靶点）:", length(triple_intersection), "\n")

# 3. 绘制 Venn 图（三集 Venn）
venn.diagram(
  x = list(
    TCM    = tcm_targets,
    NPC    = npc_targets,
    Immune = immune_targets
  ),
  filename = "results/03_venn_diagram.pdf",
  category.names = c("中药靶点", "NPC 靶点", "免疫靶点"),
  col = c("#1B9E77", "#D95F02", "#7570B3"),
  fill = c(alpha("#1B9E77", 0.3), alpha("#D95F02", 0.3), alpha("#7570B3", 0.3)),
  cex = 1.5,
  fontfamily = "sans",
  cat.cex = 1.2,
  cat.fontfamily = "sans",
  main = "中药-NPC-免疫靶点三维 Venn 图",
  output = TRUE
)

# 4. 保存双效靶点列表
write_csv(data.frame(gene_symbol = triple_intersection),
          "results/03_dual_effect_targets.csv")
cat("双效靶点已保存至 results/03_dual_effect_targets.csv\n")
\end{lstlisting}

\subsection{靶点数据集标准化与注释规范}
\label{sec:module3_annotation}

\begin{lstlisting}[language=R,caption={Gene ID 转换与靶点注释标准化}]
# ==== 靶点 Gene ID 标准化与注释 ====
library(AnnotationDbi)
library(org.Hs.eg.db)
library(dplyr)

# 1. 加载双效靶点列表
targets <- read_csv("results/03_dual_effect_targets.csv")$gene_symbol

# 2. Gene Symbol 转 Entrez ID（用于后续 clusterProfiler 富集分析）
target_entrez <- mapIds(org.Hs.eg.db,
                        keys = targets,
                        keytype = "SYMBOL",
                        column = "ENTREZID",
                        multiVals = "first")

# 3. 检查转换成功率（应 >= 90%）
success_rate <- sum(!is.na(target_entrez)) / length(target_entrez)
cat("Gene Symbol → Entrez ID 转换成功率:", round(success_rate*100, 1), "%\n")
if (success_rate < 0.90) warning("转换成功率低于 90%，请检查 Gene Symbol 大小写格式！")

# 4. 添加 UniProt ID（用于分子对接受体选择）
target_uniprot <- mapIds(org.Hs.eg.db,
                         keys = targets,
                         keytype = "SYMBOL",
                         column = "UNIPROT",
                         multiVals = "first")

# 5. 整合注释信息
target_annotation <- data.frame(
  gene_symbol = targets,
  entrez_id   = target_entrez,
  uniprot_id  = target_uniprot
) %>%
  filter(!is.na(entrez_id))  # 移除无法转换的基因

write_csv(target_annotation, "results/03_target_annotation_final.csv")
cat("最终注释靶点数量:", nrow(target_annotation), "\n")
\end{lstlisting}

\begin{qcbox}
\textbf{模块三质控合格标准：}
\begin{itemize}
  \item 已从 $\geq$ 3 个独立数据库（GeneCards + DisGeNET + GEO 差异分析）获取 NPC 靶点 \ding{52}
  \item GEO 差异分析采用 limma 包，FDR 校正方法为 Benjamini-Hochberg，阈值 FDR $<$ 0.05 且 $|$log2FC$|$ $\geq$ 1 \ding{52}
  \item 免疫靶点集已纳入 ImmPort 标准基因集 + 核心免疫检查点基因（共 $\geq$ 500 个基因）\ding{52}
  \item Venn 图三维交集结果已生成（PDF 格式），双效靶点数量 $\geq$ 10 个 \ding{52}
  \item Gene Symbol → Entrez ID 转换成功率 $\geq$ 90\% \ding{52}
  \item 最终双效靶点列表已含 UniProt ID（用于分子对接受体蛋白查找）\ding{52}
\end{itemize}
\end{qcbox}

\moduleoutput{双效靶点交集列表（CSV，含 Gene Symbol + Entrez ID + UniProt ID）、Venn 图（PDF），输入模块四的网络构建}

% ============================================================
% MODULE 4: 基于多维数据的中药活性成分虚拟筛选全流程
% ============================================================
\section{模块四：基于多维数据的中药活性成分虚拟筛选全流程}
\label{sec:module4}

\begin{infobox}{模块四概览}
\textbf{操作目的：}通过网络药理学与分子对接双维度虚拟筛选，在计算层面系统预测中药活性成分与双效靶点的结合能力与作用机制，确定进入湿实验验证的候选成分。\\
\textbf{前置输入：}模块二的活性成分-SMILES 数据库、模块三的双效靶点列表（含 UniProt ID）。\\
\textbf{本模块输出：}成分-靶点-通路多维网络图、分子对接打分排序表、候选活性成分最终确认列表。\\
\textbf{预计耗时：}2--3 周。
\end{infobox}

\subsection{网络药理学核心分析标准化操作}
\label{sec:module4_network}

\subsubsection{步骤一：中药-成分-靶点-疾病四级网络构建（Cytoscape 3.10.1）}

\begin{enumerate}
  \item \textbf{准备网络输入文件（Excel 格式）：}
    新建 Excel 文件，按以下格式整理「成分-靶点」映射表（两列）：
    \begin{center}
    \begin{tabular}{ll}
    \toprule
    \textbf{Source（成分）} & \textbf{Target（靶点）} \\
    \midrule
    连翘苷 (Phillyrin) & TP53 \\
    连翘苷 (Phillyrin) & AKT1 \\
    槲皮素 (Quercetin) & EGFR \\
    槲皮素 (Quercetin) & CD274 \\
    ... & ... \\
    \bottomrule
    \end{tabular}
    \end{center}
    保存为 \texttt{04\_compound\_target\_network.csv}。

  \item \textbf{导入 Cytoscape 3.10.1：}
    \begin{enumerate}
      \item 打开 \software{Cytoscape 3.10.1} 软件（cytoscape.org 下载安装）。
      \item 点击顶部菜单栏 \clickpath{File → Import → Network from File}。
      \item 在文件选择对话框中，选择 \texttt{04\_compound\_target\_network.csv}，点击「Open」。
      \item 在「Import Network from Table」对话框中：
        \begin{itemize}
          \item 将「Source」列设置为 \textbf{Source Node}（绿色箭头图标）
          \item 将「Target」列设置为 \textbf{Target Node}（蓝色箭头图标）
        \end{itemize}
      \item 点击「OK」完成导入。
    \end{enumerate}

  \item \textbf{设置网络节点属性（拓扑分析）：}
    \begin{enumerate}
      \item 点击菜单栏 \clickpath{Tools → NetworkAnalyzer → Analyze Network}，选择「Undirected」（无向网络），点击「OK」。
      \item 分析完成后，在节点属性表（Node Table）中新增以下列：
        \begin{itemize}
          \item \textbf{Degree}（度值，连接数量）
          \item Betweenness Centrality（介数中心性）
          \item Closeness Centrality（接近中心性）
        \end{itemize}
      \item 核心成分筛选标准：选择 \textbf{Degree $\geq$ 中位数 $\times$ 2 AND Betweenness Centrality $\geq$ 中位数 $\times$ 2} 的节点为关键成分。
    \end{enumerate}

  \item \textbf{网络可视化：}
    \begin{enumerate}
      \item 点击菜单栏 \clickpath{Layout → Prefuse Force Directed Layout}，选择「Weight」为 Degree（调整节点布局）。
      \item 在「Style」面板：
        \begin{itemize}
          \item 节点大小（Node Size）映射：Degree 值 → 节点大小（值越大节点越大）
          \item 节点颜色：成分节点设为橙色（\#FF8C00），靶点节点设为深蓝色（\#003366）
          \item 边线颜色：灰色（\#AAAAAA），粗细均匀
        \end{itemize}
      \item 点击 \clickpath{File → Export → Network to Image}，选择 PDF 格式，分辨率 \textbf{300 dpi}，保存。
    \end{enumerate}
\end{enumerate}

\subsubsection{步骤二：PPI 网络构建与核心靶点筛选}

\begin{enumerate}
  \item \textbf{STRING 数据库 PPI 网络生成：}
    \begin{enumerate}
      \item 访问 \texttt{https://string-db.org}。
      \item 点击「Multiple Proteins」，将模块三输出的双效靶点 Gene Symbol 列表（每行一个）粘贴至输入框。
      \item 设置参数：Organism = \textbf{Homo sapiens (9606)}，Minimum interaction score = \textbf{0.700}（高置信度阈值），点击「Continue」。
      \item 在结果页面，点击「Exports」，选择「Interactions (TSV)」，下载 PPI 数据文件，保存为 \texttt{04\_STRING\_PPI.tsv}。
    \end{enumerate}

  \item \textbf{CytoHubba 插件筛选核心靶点（Cytoscape 3.10.1）：}
    \begin{enumerate}
      \item 在 Cytoscape 中，点击 \clickpath{Apps → App Manager}，搜索「CytoHubba」，点击「Install」安装插件。
      \item 导入 STRING PPI 数据：\clickpath{File → Import → Network from File}，选择 \texttt{04\_STRING\_PPI.tsv}。
      \item 点击 \clickpath{Apps → cytoHubba → Start}，在 CytoHubba 面板中：
        \begin{itemize}
          \item 「Number of nodes」设置为 \textbf{10}（输出前 10 个 Hub 靶点）
          \item 「Hub method」选择 \textbf{MCC}（Maximum Clique Centrality，最大团中心性，为网络药理学文献公认的最优 Hub 筛选算法）
        \end{itemize}
      \item 点击「Submit」，等待计算完成。
      \item 记录输出的前 10 个 Hub 靶点（Gene Symbol + Score），保存为 \texttt{04\_hub\_targets\_top10.csv}。
    \end{enumerate}
\end{enumerate}

\subsubsection{步骤三：GO 富集分析与 KEGG 通路富集分析}

\begin{lstlisting}[language=R,caption={GO/KEGG 富集分析（clusterProfiler）}]
# ==== GO & KEGG 富集分析 ====
# 参考：Ritchie et al., Nucleic Acids Research, 2015

library(clusterProfiler)
library(org.Hs.eg.db)
library(ggplot2)
library(enrichplot)

# 1. 读取双效靶点的 Entrez ID
targets <- read_csv("results/03_target_annotation_final.csv")
entrez_ids <- targets$entrez_id[!is.na(targets$entrez_id)]
cat("输入富集分析的靶点数量:", length(entrez_ids), "\n")

# 2. GO 富集分析（BP: Biological Process 生物过程）
go_result <- enrichGO(
  gene          = entrez_ids,
  OrgDb         = org.Hs.eg.db,
  ont           = "BP",            # BP=生物过程; CC=细胞组分; MF=分子功能
  pAdjustMethod = "BH",            # Benjamini-Hochberg 校正
  pvalueCutoff  = 0.05,
  qvalueCutoff  = 0.2,
  readable      = TRUE             # 将 Entrez ID 转换为 Gene Symbol 显示
)

cat("GO-BP 显著条目数量:", nrow(go_result@result), "\n")

# 3. KEGG 通路富集分析
kegg_result <- enrichKEGG(
  gene         = entrez_ids,
  organism     = "hsa",            # 人类
  pAdjustMethod = "BH",
  pvalueCutoff = 0.05,
  qvalueCutoff = 0.2
)

cat("KEGG 显著通路数量:", nrow(kegg_result@result), "\n")

# 4. 气泡图可视化（GO）
p_go <- dotplot(go_result,
                showCategory = 20,
                font.size = 10) +
  ggtitle("GO-BP Enrichment Analysis") +
  theme_bw() +
  theme(axis.text.y = element_text(size = 9))

ggsave("results/04_GO_BP_dotplot.pdf", p_go,
       width = 10, height = 8, dpi = 300)

# 5. 气泡图可视化（KEGG）
p_kegg <- dotplot(kegg_result,
                  showCategory = 20,
                  font.size = 10) +
  ggtitle("KEGG Pathway Enrichment Analysis") +
  theme_bw()

ggsave("results/04_KEGG_dotplot.pdf", p_kegg,
       width = 10, height = 8, dpi = 300)

# 6. 保存富集分析结果
write_csv(as.data.frame(go_result), "results/04_GO_BP_enrichment.csv")
write_csv(as.data.frame(kegg_result), "results/04_KEGG_enrichment.csv")
\end{lstlisting}

\subsection{分子对接虚拟筛选全流程操作}
\label{sec:module4_docking}

\subsubsection{步骤一：配体小分子 3D 结构准备}

\begin{enumerate}
  \item \textbf{配体文件批量下载（PubChem）：}
    \begin{enumerate}
      \item 访问 \texttt{https://pubchem.ncbi.nlm.nih.gov/idexchange}。
      \item 在「Input IDs」框中粘贴所有候选成分的 SMILES 列表，「Output IDs」选择「SDF」（3D conformer）。
      \item 或使用 Python 批量下载脚本：
    \end{enumerate}
\end{enumerate}

\begin{lstlisting}[language=Python,caption={批量下载配体 3D 结构（Python）}]
# ==== 配体批量下载脚本 ====
import pubchempy as pcp
import pandas as pd
import os

# 读取候选成分列表
compounds_df = pd.read_csv("results/02_adme_filtered_compounds.csv")
smiles_list = compounds_df['smiles'].dropna().tolist()
names_list = compounds_df['compound_name'].dropna().tolist()

os.makedirs("docking/ligands_sdf", exist_ok=True)

for name, smiles in zip(names_list, smiles_list):
    try:
        # 通过 SMILES 在 PubChem 查找化合物
        compounds = pcp.get_compounds(smiles, 'smiles')
        if compounds:
            cid = compounds[0].cid
            # 下载 3D SDF 文件
            pcp.download('SDF', f"docking/ligands_sdf/{name}.sdf",
                         cid, record_type='3d', overwrite=True)
            print(f"Downloaded: {name} (CID: {cid})")
        else:
            print(f"Not found in PubChem: {name}")
    except Exception as e:
        print(f"Error for {name}: {e}")

print("批量下载完成")
\end{lstlisting}

\begin{enumerate}
  \setcounter{enumi}{1}
  \item \textbf{配体结构优化（RDKit 2023.09.1）：}
\end{enumerate}

\begin{lstlisting}[language=Python,caption={配体 3D 构象优化（RDKit）}]
# ==== 配体结构优化 ====
from rdkit import Chem
from rdkit.Chem import AllChem, Descriptors
import os
import glob

input_dir = "docking/ligands_sdf/"
output_dir = "docking/ligands_optimized/"
os.makedirs(output_dir, exist_ok=True)

sdf_files = glob.glob(f"{input_dir}*.sdf")

for sdf_file in sdf_files:
    name = os.path.basename(sdf_file).replace('.sdf', '')
    try:
        mol = Chem.MolFromMolFile(sdf_file, removeHs=False)
        if mol is None:
            print(f"Error reading: {name}")
            continue

        # 加氢
        mol_h = Chem.AddHs(mol)

        # MMFF94 力场优化 3D 构象
        result = AllChem.MMFFOptimizeMolecule(mol_h, maxIters=2000)
        if result == 0:
            print(f"Optimized: {name}")
        elif result == 1:
            print(f"Not converged (acceptable): {name}")
        else:
            # 备选: UFF 力场
            AllChem.EmbedMolecule(mol_h, AllChem.ETKDG())
            AllChem.UFFOptimizeMolecule(mol_h, maxIters=2000)
            print(f"UFF optimized: {name}")

        # 输出优化后 SDF
        writer = Chem.SDWriter(f"{output_dir}{name}_opt.sdf")
        writer.write(mol_h)
        writer.close()

    except Exception as e:
        print(f"Failed: {name} - {e}")

# SDF 转 PDBQT（用于 AutoDock Vina）
# 使用 OpenBabel 执行
import subprocess
pdbqt_dir = "docking/ligands_pdbqt/"
os.makedirs(pdbqt_dir, exist_ok=True)
for sdf_file in glob.glob(f"{output_dir}*.sdf"):
    name = os.path.basename(sdf_file).replace('_opt.sdf','')
    cmd = f"obabel {sdf_file} -O {pdbqt_dir}{name}.pdbqt"
    subprocess.run(cmd, shell=True, capture_output=True)
    print(f"Converted to PDBQT: {name}")

print("配体准备完成")
\end{lstlisting}

\subsubsection{步骤二：受体蛋白结构准备}

核心靶点受体晶体结构选择标准：

\begin{center}
\begin{tabular}{p{2.5cm}p{2cm}p{2cm}p{6cm}}
\toprule
\textbf{靶点} & \textbf{推荐 PDB ID} & \textbf{分辨率} & \textbf{选择理由} \\
\midrule
PD-L1 (CD274) & 5J89 & 1.50 Å & 与抑制剂共晶，活性口袋明确 \\
EGFR & 1IVO & 2.60 Å & 激酶域与 ATP 结合口袋 \\
TP53 & 2OCJ & 2.05 Å & DNA 结合域，中药干预位点 \\
AKT1 & 4EJN & 1.99 Å & 激酶活性区，含 ATP 类似物 \\
VEGFA & 3QTK & 2.20 Å & 与抗体结合区域明确 \\
STAT3 & 6NJS & 1.70 Å & SH2 域，小分子结合口袋 \\
\bottomrule
\end{tabular}
\end{center}

受体预处理步骤（以 PyMOL 2.5 为例）：

\begin{enumerate}
  \item 打开 \software{PyMOL 2.5}（pymol.org 下载安装）。
  \item 在命令行输入：\texttt{fetch 5J89}（自动下载 PDB 结构）。
  \item 执行以下 PyMOL 命令完成预处理：
\begin{lstlisting}[language={},caption={PyMOL 受体预处理命令}]
# 1. 移除水分子
remove solvent

# 2. 移除非必要配体（保留共晶配体用于活性口袋定位）
# 先查看配体名称
select ligands, resn UNK or resn HET
show sticks, ligands

# 3. 选择并提取蛋白链（去除共晶配体）
select protein_only, chain A and polymer.protein
save receptor_clean.pdb, protein_only

# 4. 在 AutoDockTools 中加极性氢并生成 PDBQT
# (PyMOL 仅用于可视化，PDBQT 生成需用 AutoDockTools 或 prepare_receptor4.py)
\end{lstlisting}

  \item 使用 \software{AutoDockTools 1.5.7} 完成 PDBQT 生成：
\begin{lstlisting}[language=bash,caption={AutoDockTools 受体转换命令}]
# 使用 prepare_receptor4.py 脚本（AutoDockTools 组件）
python prepare_receptor4.py \
    -r receptor_clean.pdb \    # 输入 PDB 文件
    -o receptor.pdbqt \        # 输出 PDBQT 文件
    -A hydrogens \             # 添加极性氢原子
    -U nphs \                  # 非极性氢原子合并
    -C                         # 保留原有电荷（不重新计算 Gasteiger 电荷）
\end{lstlisting}
\end{enumerate}

\subsubsection{步骤三：AutoDock Vina 1.2.3 分子对接操作}

\begin{enumerate}
  \item \textbf{活性口袋坐标确定：}
    \begin{itemize}
      \item 若 PDB 晶体结构含共晶配体：以共晶配体的几何中心坐标作为对接盒子（grid box）中心，使用 PyMOL 命令获取：\texttt{get\_position → 计算中心点坐标}
      \item 若 PDB 无共晶配体：使用在线工具 \database{FPocket}（fpocket.sourceforge.net）或 \database{CastP}（sts.bioe.uic.edu/castp）预测活性口袋
    \end{itemize}

  \item \textbf{配置文件准备（config.txt）：}
\begin{lstlisting}[language={},caption={AutoDock Vina 参数配置文件}]
# config.txt for AutoDock Vina 1.2.3
# 以 PD-L1 (PDB: 5J89) 为例

receptor = receptor.pdbqt     # 受体 PDBQT 文件路径

# 对接盒子参数（根据共晶配体位置调整）
center_x = 15.024              # 活性口袋 X 中心坐标
center_y = 5.176               # 活性口袋 Y 中心坐标
center_z = 22.831              # 活性口袋 Z 中心坐标

# 盒子尺寸（单位: Å）
size_x = 20                    # 覆盖活性口袋 + 2 Å 余量
size_y = 20
size_z = 20

# 搜索参数
exhaustiveness = 32            # 搜索精度，提高至 32（默认 8），适用于批量筛选
num_modes = 10                 # 每个配体输出前 10 个结合构象
energy_range = 4               # 能量范围（kcal/mol）

# 精度模式（Vina 1.2.3 新特性）
scoring = vina                 # 评分函数：vina 或 vinardo（对抗体靶点推荐）
\end{lstlisting}

  \item \textbf{批量对接执行脚本（Python）：}
\end{enumerate}

\begin{lstlisting}[language=Python,caption={批量分子对接脚本（AutoDock Vina 1.2.3）}]
# ==== 批量分子对接脚本 ====
import subprocess
import pandas as pd
import os
import glob
import re

ligands_dir = "docking/ligands_pdbqt/"
output_dir  = "docking/results/"
config_file = "docking/config.txt"
os.makedirs(output_dir, exist_ok=True)

results = []
ligand_files = glob.glob(f"{ligands_dir}*.pdbqt")

for ligand_file in ligand_files:
    ligand_name = os.path.basename(ligand_file).replace('.pdbqt','')
    output_file = f"{output_dir}{ligand_name}_out.pdbqt"
    log_file    = f"{output_dir}{ligand_name}_log.txt"

    # 执行 AutoDock Vina 1.2.3 对接
    cmd = [
        "vina",
        "--config",  config_file,
        "--ligand",  ligand_file,
        "--out",     output_file,
        "--log",     log_file
    ]

    print(f"Docking: {ligand_name}...")
    proc = subprocess.run(cmd, capture_output=True, text=True)

    # 解析结合自由能（第一位结合模式）
    best_score = None
    if os.path.exists(log_file):
        with open(log_file) as f:
            for line in f:
                match = re.search(r'^\s+1\s+([-\d.]+)', line)
                if match:
                    best_score = float(match.group(1))
                    break

    results.append({
        'ligand': ligand_name,
        'best_score_kcal_mol': best_score
    })
    print(f"  {ligand_name}: {best_score} kcal/mol")

# 整理结果并按打分排序
results_df = pd.DataFrame(results)
results_df = results_df.sort_values('best_score_kcal_mol')  # 分数越小越好
results_df['rank'] = range(1, len(results_df)+1)

results_df.to_csv("docking/04_docking_results_ranked.csv", index=False)
print("\n=== 对接结果排名（前 10 位）===")
print(results_df.head(10).to_string())
\end{lstlisting}

\begin{warningbox}
\textbf{对接结果评价标准（重要）：}
\begin{itemize}
  \item \textbf{强结合（优先验证）：}结合自由能 $\leq$ $-$8.0 kcal/mol，可视为与临床药物相当的结合亲和力
  \item \textbf{中等结合（可验证）：}$-$8.0 < 结合自由能 $\leq$ $-$7.0 kcal/mol，需辅以其他证据支持
  \item \textbf{弱结合（不建议进入湿实验）：}结合自由能 $>$ $-$7.0 kcal/mol
  \item \textbf{假阳性规避：}同一成分对多个靶点均显示 $\leq$ $-$7.0 kcal/mol，方可认定为多靶点候选成分；单靶点强结合需排除晶体结构质量问题
\end{itemize}
\end{warningbox}

\subsubsection{步骤四：PyMOL 2.5 结合模式可视化}

\begin{lstlisting}[language={},caption={PyMOL 结合模式可视化命令}]
# PyMOL 2.5 操作步骤（命令行模式）
# 以最优对接结果为例

# 1. 加载受体和对接输出文件
load receptor.pdb, receptor
load docking/results/compound_out.pdbqt, ligand

# 2. 可视化设置
as cartoon, receptor
as sticks, ligand
show surface, receptor
set surface_transparency, 0.7

# 3. 显示活性口袋周围残基（距配体 5 Å 范围）
select binding_site, receptor within 5 of ligand
show sticks, binding_site
color green, binding_site

# 4. 显示氢键
distance hbonds, receptor, ligand, 3.5, mode=2
color yellow, hbonds

# 5. 着色方案
color marine, receptor
color red, ligand
set stick_radius, 0.15

# 6. 导出高分辨率图片
ray 3000, 2000                 # 光线追踪渲染，3000x2000 像素
png results/04_docking_visualization.png, dpi=300
\end{lstlisting}

\subsection{AI 大模型辅助高效筛选操作流程}
\label{sec:module4_ai}

\begin{warningbox}
\textbf{AI 辅助的边界说明：}AI 大模型（包括 ChatGPT-4o、Claude 等）可作为文献调研与假说生成的辅助工具，但其输出结果必须经过人工核实与实验验证，不可直接引用为数据。所有 AI 辅助生成的内容须在论文中明确声明，并附以验证性实验结果。
\end{warningbox}

AI 辅助操作的标准化使用场景：

\begin{enumerate}
  \item \textbf{初筛提示词设计（用于文献调研加速）：}

    标准化提示词模板：
    \begin{greenbox}
    「请基于以下信息，检索并总结最近 5 年内发表的相关文献（PubMed/Google Scholar）：
    化合物：[化合物名称/SMILES]；靶点：[靶点名称]；疾病：鼻咽癌（Nasopharyngeal Carcinoma）
    请说明：①该化合物是否有对应靶点的直接结合实验证据？②该化合物在 NPC 或相关肿瘤中的活性报告？
    注意：仅引用 PubMed 可查询的真实文献，不编造内容。」
    \end{greenbox}

  \item \textbf{人工验证步骤（必须执行）：}
    \begin{enumerate}
      \item 将 AI 输出的每篇参考文献在 PubMed（pubmed.ncbi.nlm.nih.gov）中逐一核实是否真实存在。
      \item 下载并阅读全文，确认数据与 AI 描述一致。
      \item 通过人工核实的文献方可纳入手册与论文参考文献列表。
    \end{enumerate}
\end{enumerate}

\subsection{候选中药/活性成分最终确定规则}
\label{sec:module4_selection}

建立多维度打分体系，最终入选标准须同时满足：

\begin{center}
\begin{tabular}{p{4cm}p{4cm}p{2cm}p{3.5cm}}
\toprule
\textbf{评价维度} & \textbf{评价指标} & \textbf{阈值} & \textbf{分值} \\
\midrule
ADME/T 性质 & OB \& DL 双过滤 & 均通过 & 基本准入（不计分，必须满足）\\
\addlinespace
网络药理学 & 成分 Degree 排名 & 前 20\% & +2 分 \\
\addlinespace
靶点相关性 & 与双效靶点交集数量 & $\geq$ 3 个 & +2 分 \\
\addlinespace
分子对接 & 最优结合自由能 & $\leq$ $-$8.0 kcal/mol & +3 分 \\
\addlinespace
公共数据验证 & 核心靶点 GEO 差异表达 & FDR $<$ 0.05 & +2 分 \\
\addlinespace
文献支持 & PubMed 活性报告 $\geq$ 1 篇 & 近 10 年 & +1 分 \\
\bottomrule
\textbf{合格线} & & & \textbf{总分 $\geq$ 7 分} \\
\bottomrule
\end{tabular}
\end{center}

\begin{qcbox}
\textbf{模块四质控合格标准：}
\begin{itemize}
  \item 网络拓扑分析已完成，核心成分（Degree 前 20\%）已确认 \ding{52}
  \item PPI 网络经 STRING（置信度 $\geq$ 0.700）构建，Hub 靶点（MCC 法，前 10 位）已标注 \ding{52}
  \item GO/KEGG 富集分析已执行，结果以气泡图可视化（PDF，300 dpi）\ding{52}
  \item 所有候选成分的分子对接均已完成（AutoDock Vina 1.2.3，exhaustiveness = 32）\ding{52}
  \item 最优候选成分的结合模式已用 PyMOL 2.5 可视化，氢键/疏水相互作用已标注 \ding{52}
  \item 最终入选候选成分 $\geq$ 3 种，多维评分 $\geq$ 7 分 \ding{52}
\end{itemize}
\end{qcbox}

\moduleoutput{候选活性成分最终列表（$\geq$ 3 种）、网络图（PDF）、富集分析图（PDF）、对接结果排序表（CSV）、结合模式可视化图（PNG，300 dpi），输入模块五的多组学验证与模块六的湿实验}

% ============================================================
% 补充：高通量预筛选方案（70味中药单浓度初筛策略）
% ============================================================
\subsection{补充方案：70 味中药高通量单浓度预筛选 SOP}
\label{sec:module4_hts}

\begin{infobox}{高通量预筛选方案概览}
\textbf{操作目的：}在虚拟筛选候选成分进入全剂量湿实验前，以最低工作量对候选中药库（以 70 味中药为标准规模）进行系统性生物活性初筛，快速缩小实验范围，筛出具有明确生物活性的阳性命中药物，降低后续全剂量实验的盲目性与资源消耗。\\
\textbf{核心策略：}\textbf{单浓度（100 $\mu$g/mL）× 三复孔}，以 CCK-8 法检测鼻咽癌细胞增殖抑制率，同步以单浓度检测免疫细胞激活指标（TNF-$\alpha$/IL-6 分泌量），实现直接抗肿瘤与免疫激活的双终点初筛。\\
\textbf{工作量预估：}70 味中药分 3--4 轮 96 孔板完成，单轮约 1 天工作，合计约 \textbf{2--3 周}（含样品制备与数据处理）。\\
\textbf{前置输入：}模块二完成的中药提取物库（标准化浓度，DMSO 助溶至 100 mg/mL 母液）。
\end{infobox}

\subsubsection{步骤一：中药提取物标准化母液制备}

\begin{enumerate}
  \item \textbf{提取物来源规范：}优先选用符合《中国药典》2025 版的标准化提取物（应注明提取率、批次号、质量检测证书），禁止使用批次不明或质检资料缺失的粗提物。
  \item \textbf{DMSO 母液制备：}精密称取各中药提取物 \textbf{10 mg}，加入 \textbf{100 $\mu$L DMSO（分析纯，Sigma，D8418）}，充分涡旋溶解（\textbf{60 s}），配制成 \textbf{100 mg/mL}（即 100,000 $\mu$g/mL）DMSO 母液。
  \item \textbf{溶解性初筛：}目视检查母液是否澄清（轻微浑浊可接受），若出现明显沉淀，超声处理（\textbf{37$^{\circ}$C，10 min}），仍不溶者记录为「水不溶样品」，后续以 \textbf{200 $\mu$g/mL PBS 悬液}（超声 10 min + 涡旋 30 s）替代。
  \item \textbf{工作液配制：}取 DMSO 母液 \textbf{1 $\mu$L}，加入 \textbf{999 $\mu$L 完全培养基}，配制成 \textbf{100 $\mu$g/mL}（DMSO 终浓度 \textbf{0.1\%}）工作液，现配现用，禁止存放过夜。
\end{enumerate}

\begin{warningbox}
\textbf{关键控制点：} DMSO 终浓度须严格控制在 \textbf{$\leq$ 0.1\%}（超过 0.5\% 可直接导致细胞毒性假阳性）。每轮实验均需同时设置 \textbf{0.1\% DMSO 溶剂对照}，排除溶剂干扰；若使用 PBS 悬液，须同时设置等量 PBS 空白对照。
\end{warningbox}

\subsubsection{步骤二：CNE2 细胞单浓度 CCK-8 初筛操作（抗肿瘤终点）}

\begin{enumerate}
  \item \textbf{铺板（Day 0）：}取对数生长期 CNE2 细胞（passage 5--15），台盼蓝染色确认活力 $\geq$ 95\%，以 \textbf{3,000 细胞/孔}接种于 96 孔板，每孔 \textbf{100 $\mu$L} 完全培养基（RPMI-1640 + 10\% FBS）。放入 37$^{\circ}$C、5\% CO$_2$ 培养箱，过夜贴壁（\textbf{12--16 h}）。
  \item \textbf{孔位分配（Day 1）：}按表~\ref{tab:hts_plate_layout} 分配各孔用途，每轮 96 孔板最多可检测 \textbf{27 个样品}（含对照孔）：
  \begin{itemize}
    \item \textbf{阴性对照（NC）：}0.1\% DMSO + 完全培养基，设 \textbf{6 个复孔}。
    \item \textbf{阳性对照（PC）：}顺铂（Cisplatin，终浓度 \textbf{20 $\mu$g/mL}），设 \textbf{3 个复孔}。
    \item \textbf{待筛中药：}每样品设 \textbf{3 个技术复孔}（终浓度 100 $\mu$g/mL）。
    \item \textbf{空白对照（Blank）：}无细胞 + 培养基，设 \textbf{3 个孔}（扣除背景 OD）。
  \end{itemize}
  \item \textbf{加药（Day 1）：}弃去旧培养基（每孔保留约 10 $\mu$L 避免机械损伤），每孔加入 \textbf{100 $\mu$L} 对应工作液（终浓度 100 $\mu$g/mL），37$^{\circ}$C 孵育。
  \item \textbf{孵育（Day 1--3）：}37$^{\circ}$C、5\% CO$_2$ 孵育 \textbf{48 h}（第一轮推荐，可根据细胞倍增时间调整至 24--72 h）。
  \item \textbf{CCK-8 检测（Day 3）：}每孔加入 \textbf{10 $\mu$L CCK-8 液}（Dojindo，CK04），轻轻混匀（勿产生气泡），放回培养箱孵育 \textbf{2 h}（至 NC 孔 OD$_{450}$ $\approx$ 1.0--2.0），酶标仪 \textbf{450 nm}（参考波长 630 nm）读取 OD 值。
  \item \textbf{抑制率计算：}
  \[
  \text{抑制率 (\%)} = \frac{(\overline{\text{OD}}_{\text{NC}} - \overline{\text{OD}}_{\text{Blank}}) - (\overline{\text{OD}}_{\text{sample}} - \overline{\text{OD}}_{\text{Blank}})}{(\overline{\text{OD}}_{\text{NC}} - \overline{\text{OD}}_{\text{Blank}})} \times 100\%
  \]
  \item \textbf{阳性命中阈值：}100 $\mu$g/mL 下增殖抑制率 \textbf{$\geq$ 50\%}（严格标准）；适合粗提物研究可适当放宽至 \textbf{$\geq$ 30\%}（宽松标准），但需在报告中说明。
\end{enumerate}

\begin{center}
\begin{tabular}{ccccccccccccc}
\toprule
 & \textbf{1} & \textbf{2} & \textbf{3} & \textbf{4} & \textbf{5} & \textbf{6} & \textbf{7} & \textbf{8} & \textbf{9} & \textbf{10} & \textbf{11} & \textbf{12} \\
\midrule
A & NC & NC & S01 & S01 & S01 & S07 & S07 & S07 & S13 & S13 & S13 & Blank \\
B & NC & NC & S02 & S02 & S02 & S08 & S08 & S08 & S14 & S14 & S14 & Blank \\
C & NC & NC & S03 & S03 & S03 & S09 & S09 & S09 & S15 & S15 & S15 & Blank \\
D & NC & NC & S04 & S04 & S04 & S10 & S10 & S10 & S16 & S16 & S16 & Blank \\
E & PC & PC & S05 & S05 & S05 & S11 & S11 & S11 & S17 & S17 & S17 & NC \\
F & PC & PC & S06 & S06 & S06 & S12 & S12 & S12 & \ldots & \ldots & \ldots & NC \\
G & PC & \ldots & \ldots & \ldots & \ldots & \ldots & \ldots & \ldots & \ldots & \ldots & \ldots & NC \\
H & \ldots & \ldots & \ldots & \ldots & \ldots & \ldots & \ldots & \ldots & \ldots & \ldots & \ldots & \ldots \\
\bottomrule
\end{tabular}
\captionof{table}{96 孔板高通量初筛孔位分配示意（NC=阴性对照；PC=顺铂阳性对照；S01--S17=待筛中药样品；Blank=无细胞背景孔）}\label{tab:hts_plate_layout}
\end{center}

\subsubsection{步骤三：免疫激活双终点同步初筛（THP-1 平行方案）}

\begin{enumerate}
  \item \textbf{模型选择：}使用 \textbf{THP-1 M0 巨噬细胞}（PMA 100 ng/mL 诱导分化 48 h，详见模块六 \ref{sec:module6_immune}），以 TNF-$\alpha$ 分泌量（ELISA 快速定量，总耗时约 2.5 h）作为免疫激活初筛指标，操作简便高效。
  \item \textbf{加药浓度：}\textbf{100 $\mu$g/mL}（与抗肿瘤初筛相同）；若某样品 CCK-8 抑制率 $>$ 50\%（提示强细胞毒），则将免疫细胞加药浓度降至 \textbf{1/4 CCK-8 IC50}（避免免疫毒性干扰免疫激活信号）。
  \item \textbf{免疫激活阳性命中阈值：}TNF-$\alpha$ 分泌量较 0.1\% DMSO 对照组升高 \textbf{$\geq$ 2 倍}（倍变比 $\geq$ 2.0，One-way ANOVA，P $<$ 0.05）。
\end{enumerate}

\subsubsection{步骤四：双终点综合评分与复孔验证}

\begin{enumerate}
  \item \textbf{分级命中原则（Excel 汇总矩阵）：}
  \begin{itemize}
    \item \textbf{A 级命中（优先进入 IC50 剂量效应验证）：}抑制率 $\geq$ 50\% \textit{且} TNF-$\alpha$ 倍变 $\geq$ 2 倍（双终点俱优）。
    \item \textbf{B 级命中（次优，酌情进入验证）：}抑制率 30--50\%，或单独免疫激活阳性。
    \item \textbf{排除标准：}抑制率 $<$ 30\% 且免疫激活阴性。
  \end{itemize}
  \item \textbf{复孔验证（Replication Validation）：}对所有 A 级命中，使用新鲜提取物（不同批次）独立重复实验 \textbf{2 次}（累计 $n = 3$ 次独立生物重复）；仅 3 次重复均达阳性标准的中药方可进入 IC50 全剂量实验。
  \item \textbf{阳性对照每轮确认：}每轮顺铂阳性对照 CNE2 抑制率须在 \textbf{75--95\%} 范围内，否则当轮数据视为无效，须重复实验（排除细胞状态或操作误差）。
\end{enumerate}

\begin{qcbox}
\textbf{70 味中药高通量预筛质控合格标准：}
\begin{itemize}
  \item 0.1\% DMSO 溶剂对照细胞存活率 $\geq$ 95\%（溶剂无细胞毒性）\ding{52}
  \item 顺铂阳性对照 CNE2 抑制率 75--95\%（实验体系有效性验证）\ding{52}
  \item 空白孔 OD$_{450}$ $\leq$ 0.15（排除培养基显色背景）\ding{52}
  \item 每样品三技术复孔 CV $\leq$ 15\%（精密度合格）\ding{52}
  \item A 级命中中药通过 $n = 3$ 次独立实验复孔验证 \ding{52}
  \item 筛选数据以 Excel 汇总，保留原始 OD 值文件（可溯源存档）\ding{52}
\end{itemize}
\end{qcbox}

\begin{tipbox}
「科研进阶技巧：」若经费充裕，可在单浓度 CCK-8 初筛的同时，并行加做 \textbf{Calcein AM/PI 双染法}（流式细胞术，30 min/板），可区分细胞凋亡（Calcein$^{-}$/PI$^{-}$）与坏死（PI$^{+}$），比单纯 CCK-8 提供更多凋亡机制预判信息。该数据可直接用于后期论文流式凋亡实验的结果支撑。
\end{tipbox}

\moduleoutput{70 味中药双终点命中矩阵（Excel）、A 级命中中药名单（建议 3--10 味）、复孔验证原始数据（PDF + Excel），直接输入模块六 IC50 全剂量实验优先名单}

\begin{warningbox}
\textbf{Plan B --- 70 味中药初筛全部阴性的完整应急方案}

\textbf{风险描述：}若 70 味候选中药在单浓度 100 $\mu$g/mL 初筛中双终点均未达命中阈值，属于高概率风险（中药粗提物含量不稳定），须在 2 周内启动分层应急方案，避免整体研究进度停滞。

\textbf{应对方案（按优先级顺序执行）：}
\begin{enumerate}
  \item \textbf{Step 1（1--3 天）--- 技术排查优先：}首先确认实验体系有效性（顺铂阳性对照、细胞传代数、DMSO 溶解体系）；若对照正常则确认为真阴性结果，进入 Step 2。
  \item \textbf{Step 2（3--7 天）--- 提高筛选浓度至 200 $\mu$g/mL：}同时保留 100 $\mu$g/mL 数据，评估是否存在浓度依赖效应（若 200 $\mu$g/mL 有活性，提示有效浓度偏高，后续须在论文中充分讨论毒性窗口及体内吸收问题）。
  \item \textbf{Step 3（1--2 周）--- 重构候选药物库：}使用 BATMAN-TCM 2.0 基于「清热解毒」「活血化瘀」「扶正固本」功效分类，补充纳入 20--30 味具有明确文献报道抗肿瘤或免疫调节活性的中药（排除与已筛 70 味重叠者），扩大候选库并重新初筛。
  \item \textbf{Step 4（2 周）--- 切换靶向验证策略（须导师确认）：}跳过全库表型筛选，直接对多维评分 $\geq$ 7 分的前 3--5 种候选单体化合物进行 IC50 全剂量实验，聚焦「单体成分—核心靶点」机制研究路线（适用于发表单体活性成分机制类 SCI 文章，影响因子可达 IF $\geq$ 4.0）。
  \item \textbf{Step 5（1--2 天）--- 向导师汇报并更新研究计划：}以书面形式（科研日志 + PPT）汇报初筛结果、应急方案选择理由与调整后的时间轴，确保研究方向调整有据可查。
\end{enumerate}
\end{warningbox}

% ============================================================
% MODULE 5: 多组学多维数据整合验证操作流程
% ============================================================
\section{模块五：多组学多维数据整合验证操作流程}
\label{sec:module5}

\begin{infobox}{模块五概览}
\textbf{操作目的：}利用公共数据集（TCGA/GEO）与单细胞测序数据，在患者水平验证核心靶点的临床相关性与免疫调控机制，建立「成分-靶点-通路-表型」全链条生物信息学证据链。\\
\textbf{前置输入：}模块四确认的候选成分与核心靶点列表（Hub 靶点 + 双效靶点交集）。\\
\textbf{本模块输出：}差异表达热图、生存分析 KM 曲线（PDF）、免疫浸润相关性图（PDF）、单细胞 UMAP 图（PDF）。\\
\textbf{预计耗时：}1--2 周。
\end{infobox}

\subsection{公共数据集验证操作（TCGA/GEO）}
\label{sec:module5_public}

\subsubsection{步骤一：核心靶点在 NPC 与正常组织中的差异表达验证}

\begin{lstlisting}[language=R,caption={TCGA/GEO 核心靶点差异表达验证}]
# ==== 核心靶点差异表达验证 ====
library(GEOquery)
library(limma)
library(pheatmap)
library(ggplot2)
library(dplyr)

# 1. 读取 Hub 靶点列表（模块四输出）
hub_targets <- read_csv("results/04_hub_targets_top10.csv")$gene_symbol

# 2. 从 GSE12452 提取核心靶点表达矩阵
gse <- getGEO("GSE12452", GSEMatrix = TRUE, destdir = "data/GEO/")
expr_data  <- exprs(gse[[1]])
pheno_data <- pData(gse[[1]])

# 3. 获取探针-基因对应关系（GPL570 芯片注释）
feature_data <- fData(gse[[1]])
probe_gene_map <- feature_data[, c("ID", "Gene Symbol")]
colnames(probe_gene_map) <- c("probe_id", "gene_symbol")

# 4. 将探针表达值转换为基因表达值（取均值）
probe_gene_map <- probe_gene_map %>%
  filter(gene_symbol %in% hub_targets,
         !is.na(gene_symbol),
         gene_symbol != "")

if (nrow(probe_gene_map) == 0) {
  cat("警告：hub 靶点在该芯片注释中未找到，检查基因名称格式\n")
} else {
  hub_expr <- expr_data[probe_gene_map$probe_id, ]
  rownames(hub_expr) <- probe_gene_map$gene_symbol

  # 5. 定义分组
  group <- factor(ifelse(grepl("NPC|tumor", pheno_data$characteristics_ch1,
                               ignore.case = TRUE), "NPC", "Normal"))
  design <- model.matrix(~0 + group)
  colnames(design) <- levels(group)

  # 6. 差异分析
  fit <- lmFit(hub_expr, design)
  contrast <- makeContrasts(NPC - Normal, levels = design)
  fit2 <- eBayes(contrasts.fit(fit, contrast))
  DEG_hub <- topTable(fit2, number = Inf, sort.by = "logFC")

  cat("Hub 靶点差异表达结果:\n")
  print(DEG_hub[, c("logFC", "AveExpr", "adj.P.Val")])

  # 7. 热图可视化
  annotation_col <- data.frame(Group = group)
  rownames(annotation_col) <- colnames(hub_expr)

  pheatmap(hub_expr,
           annotation_col = annotation_col,
           scale = "row",
           cluster_rows = TRUE,
           cluster_cols = TRUE,
           color = colorRampPalette(c("blue","white","red"))(100),
           show_colnames = FALSE,
           fontsize = 10,
           main = "Hub Target Expression: NPC vs. Normal (GSE12452)",
           filename = "results/05_hub_targets_heatmap.pdf",
           width = 10, height = 6)

  write_csv(as.data.frame(DEG_hub), "results/05_hub_DEGs_results.csv")
}
\end{lstlisting}

\subsubsection{步骤二：核心靶点预后价值分析（KM 生存曲线）}

\begin{lstlisting}[language=R,caption={KM 生存曲线分析（以 TCGA HNSC 替代数据为例）}]
# ==== KM 生存分析 ====
# 注：NPC 需使用 GEO 数据集（GSE102349 含预后信息）
# 或使用在线工具 GEPIA2 (gepia2.cancer-pku.cn) 的可视化分析

library(survival)
library(survminer)
library(dplyr)

# 方案 A：使用 GEPIA2 在线分析（推荐新手）
# 操作步骤：
# 1. 访问 http://gepia2.cancer-pku.cn
# 2. 点击「Survival Analysis」→「Overall Survival」
# 3. 输入基因名称（一次一个 Hub 靶点），选择 Cancer type = HNSC
# 4. 设置：Group cutoff = Median，Signature scoring = Gene mean
# 5. 下载 KM 曲线图（PDF 格式）

# 方案 B：R 语言本地分析（适合有 GEO 临床数据的情况）
# 以 GSE102349 为例（包含 NPC 患者生存数据）
gse_surv <- getGEO("GSE102349", GSEMatrix = TRUE, destdir = "data/GEO/")
expr_surv  <- exprs(gse_surv[[1]])
pheno_surv <- pData(gse_surv[[1]])

# 提取生存信息（列名需根据实际数据调整）
surv_data <- data.frame(
  time   = as.numeric(pheno_surv$`overall survival time (months):ch1`),
  status = as.numeric(pheno_surv$`overall survival (censoring):ch1`)
)

# 以关键 Hub 靶点（如 TP53）表达量分组
gene <- "TP53"    # 替换为实际 Hub 靶点名

# 获取探针
feature <- fData(gse_surv[[1]])
probe_id <- feature$ID[grep(gene, feature$`Gene Symbol`)[1]]

surv_data$expr <- expr_surv[probe_id, ]
surv_data$group <- ifelse(surv_data$expr >= median(surv_data$expr),
                          "High", "Low")
surv_data <- surv_data[complete.cases(surv_data), ]

# 绘制 KM 曲线
fit <- survfit(Surv(time, status) ~ group, data = surv_data)
p <- ggsurvplot(fit,
                data = surv_data,
                pval = TRUE,
                risk.table = TRUE,
                title = paste0("Overall Survival by ", gene, " Expression"),
                xlab = "Time (months)",
                ylab = "Survival Probability",
                legend.labs = c("High", "Low"),
                palette = c("#E41A1C", "#377EB8"),
                risk.table.height = 0.3,
                ggtheme = theme_bw())

ggsave(paste0("results/05_KM_", gene, ".pdf"),
       plot = print(p, newpage = FALSE),
       width = 8, height = 7, dpi = 300)
\end{lstlisting}

\subsection{单细胞测序数据挖掘操作}
\label{sec:module5_scrna}

\begin{lstlisting}[language=R,caption={scRNA-seq 数据分析（Seurat 标准流程）}]
# ==== 单细胞测序分析 ====
# 数据来源：GSE150430（NPC TME 单细胞数据集，Chen et al. 2023）或 GSE162025
# 注意：请在 GEO 核实数据集名称与样本信息，避免使用非 NPC 数据集

library(Seurat)
library(ggplot2)
library(dplyr)

# 1. 数据加载（以 10X Genomics 格式为例）
data_10x <- Read10X("data/scRNA/NPC_10X/")
seurat_obj <- CreateSeuratObject(
  counts = data_10x,
  project = "NPC_scRNAseq",
  min.cells = 3,
  min.features = 200
)

# 2. 质量控制（QC）
seurat_obj[["percent.mt"]] <- PercentageFeatureSet(seurat_obj, pattern = "^MT-")

# 可视化 QC 指标
VlnPlot(seurat_obj,
        features = c("nFeature_RNA", "nCount_RNA", "percent.mt"),
        ncol = 3)

# 3. 过滤低质量细胞（严格标准）
seurat_obj <- subset(seurat_obj,
  subset = nFeature_RNA > 200 &     # 基因数 > 200
           nFeature_RNA < 6000 &    # 基因数 < 6000（过滤双细胞）
           percent.mt < 20)          # 线粒体基因 < 20%

cat("过滤后细胞数量:", ncol(seurat_obj), "\n")

# 4. 标准化 → PCA → 批次校正 → UMAP 降维
seurat_obj <- NormalizeData(seurat_obj)
seurat_obj <- FindVariableFeatures(seurat_obj, nfeatures = 3000)
seurat_obj <- ScaleData(seurat_obj)
seurat_obj <- RunPCA(seurat_obj, npcs = 30)
seurat_obj <- FindNeighbors(seurat_obj, dims = 1:20)
seurat_obj <- FindClusters(seurat_obj, resolution = 0.5)
seurat_obj <- RunUMAP(seurat_obj, dims = 1:20)

# 5. 可视化 UMAP（细胞聚类）
DimPlot(seurat_obj, reduction = "umap", label = TRUE) +
  ggtitle("NPC scRNA-seq: Cell Clustering")
ggsave("results/05_UMAP_clusters.pdf", width = 8, height = 6, dpi = 300)

# 6. 在各细胞亚群中可视化核心靶点表达
hub_targets <- c("CD274", "PDCD1", "IFNG", "TNF")  # 替换为实际 Hub 靶点

FeaturePlot(seurat_obj,
            features = hub_targets,
            ncol = 2,
            cols = c("lightgrey", "red"),
            pt.size = 0.5)
ggsave("results/05_hub_target_feature_plot.pdf", width = 12, height = 10, dpi = 300)

# 7. 各 Hub 靶点在细胞亚群中的表达量统计
avg_expr <- AverageExpression(seurat_obj,
                               features = hub_targets,
                               group.by = "seurat_clusters")
write_csv(as.data.frame(avg_expr$RNA),
          "results/05_hub_expr_by_cluster.csv",
          col_names = TRUE)
\end{lstlisting}

\subsection{临床相关性分析操作}
\label{sec:module5_clinical}

执行以下步骤分析核心靶点与免疫浸润的相关性：

\begin{enumerate}
  \item \textbf{TIMER 2.0 在线分析（推荐，操作简便）：}
    \begin{enumerate}
      \item 访问 \texttt{https://timer.cistrome.org}。
      \item 点击「GENE」→「Immune Association」。
      \item 在「Gene」框输入 Hub 靶点名称（如 \gene{CD274}）。
      \item 在「Cancer type」选择「HNSC」（Head and Neck Squamous Cell Carcinoma，NPC 代用）。
      \item 点击「Submit」，查看该靶点与六种免疫细胞浸润（B cell, CD4\textsuperscript{+} T, CD8\textsuperscript{+} T, Neutrophil, Macrophage, Dendritic cell）的相关系数（Spearman $r$ 与 $p$ 值）。
      \item 下载散点图（PDF 格式），保存为 \texttt{05\_[gene]\_immune\_corr.pdf}。
    \end{enumerate}

  \item \textbf{CIBERSORT 免疫浸润分析（R 语言，定量分析）：}
\end{enumerate}

\begin{lstlisting}[language=R,caption={CIBERSORT 免疫浸润分析}]
# ==== CIBERSORT 免疫浸润分析 ====
# 使用 immunedeconv 包（CIBERSORT 方法需注册 cibersort.stanford.edu）

library(immunedeconv)
library(ggplot2)
library(dplyr)
library(pheatmap)

# 1. 准备 GEO 表达矩阵（需为 TPM 或 FPKM 格式）
expr_matrix <- read_csv("data/GEO/GSE12452_expr_matrix.csv",
                         row.names = 1) %>%
  as.matrix()

# 2. 执行 CIBERSORT 推断
# 注：CIBERSORT 需先在网站注册获取许可，此处使用 TIMER 方法作为替代
deconv_result <- deconvolute(expr_matrix, method = "timer",
                              indications = rep("HNSC", ncol(expr_matrix)))

# 3. 结合临床分组
group_info <- data.frame(
  sample_id = colnames(expr_matrix),
  group = ifelse(grepl("NPC", colnames(expr_matrix)), "NPC", "Normal")
)

# 4. 可视化免疫细胞浸润比例
deconv_long <- deconv_result %>%
  pivot_longer(cols = -cell_type, names_to = "sample_id", values_to = "score") %>%
  left_join(group_info, by = "sample_id")

p_boxplot <- ggplot(deconv_long, aes(x = cell_type, y = score, fill = group)) +
  geom_boxplot(outlier.size = 0.5) +
  scale_fill_manual(values = c("NPC" = "#E41A1C", "Normal" = "#377EB8")) +
  labs(title = "Immune Cell Infiltration: NPC vs. Normal",
       x = "Immune Cell Type", y = "Deconvolution Score") +
  theme_bw() +
  theme(axis.text.x = element_text(angle = 45, hjust = 1))

ggsave("results/05_immune_infiltration_boxplot.pdf", p_boxplot,
       width = 12, height = 6, dpi = 300)

cat("免疫浸润分析完成\n")
\end{lstlisting}

\begin{qcbox}
\textbf{模块五质控合格标准：}
\begin{itemize}
  \item Hub 靶点差异表达验证完成，热图已生成（PDF，300 dpi）\ding{52}
  \item KM 生存曲线已完成（$\geq$ 3 个 Hub 靶点，P 值已标注）\ding{52}
  \item scRNA-seq 分析（若有数据集）：UMAP 降维图 + 特征基因表达图已生成 \ding{52}
  \item TIMER 2.0 免疫浸润相关性分析完成（每个 Hub 靶点 $\geq$ 1 张散点图）\ding{52}
  \item 所有 GEO 数据集 accession 编号已记录，可供论文数据可用性声明 \ding{52}
\end{itemize}
\end{qcbox}

\moduleoutput{核心靶点差异表达热图、KM 生存曲线 PDF（$\geq$ 3 个靶点）、免疫浸润散点图 PDF，作为论文生信分析章节的核心配图}

% ============================================================
% MODULE 6: 湿实验验证全流程标准化操作手册
% ============================================================
\section{模块六：湿实验验证全流程标准化操作手册（SOP）}
\label{sec:module6}

\begin{infobox}{模块六概览}
\textbf{操作目的：}在细胞水平验证候选中药活性成分的双重药效（直接抗肿瘤 + 免疫激活），为分子机制研究提供表型基础。\\
\textbf{前置输入：}模块四确认的候选活性成分（纯度 $\geq$ 98\%，已配制 10 mmol/L DMSO 储液）。\\
\textbf{实验室要求：}BSL-2 生物安全实验室；超净工作台（A2 型或以上）；CO$_2$ 培养箱。\\
\textbf{核心细胞系：}CNE1、CNE2（鼻咽癌细胞系）；Jurkat（T 细胞模型）；THP-1（单核细胞/巨噬细胞模型）。\\
\textbf{预计耗时：}8--12 周。
\end{infobox}

\subsection{体外抗肿瘤活性验证实验}
\label{sec:module6_antitumor}

\subsubsection{CNE1/CNE2 细胞系标准培养 SOP}

\paragraph{培养基配制：}

\begin{enumerate}
  \item 取 RPMI-1640 培养基（Gibco，目录号 11875119）500 mL，加热至 37$^{\circ}$C 水浴备用。
  \item 加入 \textbf{56$^{\circ}$C 热灭活胎牛血清（FBS）}（Gibco，目录号 10100147）\textbf{50 mL}（终浓度 10\%）。
  \item 加入 100$\times$ 青霉素/链霉素溶液（Pen/Strep）\textbf{5 mL}（终浓度：青霉素 100 U/mL，链霉素 100 $\mu$g/mL）。
  \item 混匀后用 0.22 $\mu$m 无菌过滤器过滤，分装于 50 mL 离心管，储存于 4$^{\circ}$C（使用有效期 4 周）。
\end{enumerate}

\paragraph{细胞传代 SOP（标准流程）：}

\begin{enumerate}
  \item 在超净工作台内，用眼观察培养瓶内细胞（倒置显微镜），确认细胞汇合率（confluency）达到 \textbf{80--90\%} 时进行传代。
  \item 弃去培养基，加入 \textbf{1$\times$ PBS（磷酸盐缓冲液，pH 7.4）5 mL}，轻摇洗涤 30 s，弃去。重复一次。
  \item 加入 \textbf{0.25\% 胰蛋白酶-EDTA（Trypsin-EDTA）2 mL}，37$^{\circ}$C 培养箱中孵育 \textbf{2--3 min}（CNE1/CNE2 消化时间偏短，过度消化损伤细胞活力）。
  \item 倒置显微镜下观察，当 $\geq$ 90\% 细胞脱壁呈圆形时，立即加入完全培养基 \textbf{5 mL} 终止消化，轻轻吹打 3--5 次使细胞充分分散。
  \item 转入 15 mL 离心管，\textbf{300 $\times$ g，5 min} 离心，弃上清。
  \item 用完全培养基重悬沉淀，按 \textbf{1:3 比例}（即 1/3 细胞接种至新瓶）接种于 T75 培养瓶，补充完全培养基至 \textbf{15 mL}。
  \item 37$^{\circ}$C，5\% CO$_2$ 培养箱中培养，\textbf{每 2--3 天传代一次}（CNE1/CNE2 倍增时间约 24--36 h）。
\end{enumerate}

\begin{warningbox}
\begin{itemize}
  \item 胰蛋白酶孵育时间超过 5 min 会导致细胞活力下降，影响后续 CCK-8 实验基线，须严格计时。
  \item CNE1/CNE2 传代代次须控制在复苏后 \textbf{5--20 代}（passage 5--20）内用于实验，超代会导致实验结果不可重复。
  \item 每次传代须记录：日期、传代代次、汇合率、培养基批次。
\end{itemize}
\end{warningbox}

\paragraph{细胞冻存与复苏 SOP：}

\textbf{冻存操作：}
\begin{enumerate}
  \item 消化细胞，计数（锥虫蓝排除法），将细胞密度调整为 \textbf{$2 \times 10^6$ 个/mL}。
  \item 配制冻存液：90\% FBS + 10\% DMSO（Sigma，D2650，细胞冻存级），轻混均匀。
  \item 将细胞悬液与等体积冻存液混合（终 DMSO 浓度为 5\%），分装至冻存管，每管 \textbf{1 mL}。
  \item 置于程序降温盒（「Mr. Frosty」或等效设备）中，-80$^{\circ}$C 降温过夜（降温速率 $\approx$ -1$^{\circ}$C/min）。
  \item 次日转入液氮罐长期保存（$-$196$^{\circ}$C）。
\end{enumerate}

\textbf{复苏操作：}
\begin{enumerate}
  \item 从液氮中取出冻存管，在 37$^{\circ}$C 水浴中快速解冻（$<$ 2 min）。
  \item 在超净台内，加入预热完全培养基 \textbf{9 mL}，轻柔混匀，离心 300 $\times$ g，5 min。
  \item 弃上清，用完全培养基重悬，接种至 T25 培养瓶，次日换液去除残余 DMSO。
\end{enumerate}

\subsubsection{CCK-8 实验 SOP（IC50 检测）}

\paragraph{实验原理：}CCK-8（Cell Counting Kit-8）试剂含有高度水溶性的 WST-8 四唑盐，可被细胞内脱氢酶还原为橙色 formazan（甲臜）产物，其 450 nm 吸光值与活细胞数量正相关。

\paragraph{操作步骤：}

\begin{enumerate}
  \item \textbf{铺板（Day 0）：}
    \begin{enumerate}
      \item 消化 CNE1/CNE2 细胞，计数，用完全培养基调整密度为 \textbf{$5 \times 10^3$ 个/mL}（200 $\mu$L/孔即每孔 1,000 个细胞，依细胞增殖速度调整，CNE2 倍增快，建议 800 个/孔）。
      \item 吸取 100 $\mu$L 细胞悬液加入 96 孔板（每孔 $5 \times 10^2$ 个细胞）。
      \item 在 37$^{\circ}$C 培养箱中培养 \textbf{24 h}，使细胞贴壁稳定。
    \end{enumerate}

  \item \textbf{加药（Day 1）：}
    \begin{enumerate}
      \item 将候选成分的 DMSO 储液（10 mmol/L）用完全培养基梯度稀释，配制 7 个浓度梯度（以 3 倍稀释为例）：
        \begin{center}
        \begin{tabular}{llllllll}
        \toprule
        \textbf{浓度梯度（$\mu$mol/L）} & 100 & 33.3 & 11.1 & 3.70 & 1.23 & 0.41 & 0.14 \\
        \bottomrule
        \end{tabular}
        \end{center}
      \item 每个浓度设 \textbf{6 个复孔}（parallel replicates）。
      \item 阴性对照：等体积 DMSO 溶液（终浓度 $\leq$ 0.1\%），每组 6 孔。
      \item 阳性对照：顺铂（cisplatin）或 5-FU（标准浓度范围需预实验确定，CNE1 顺铂 IC50 参考值约 5--10 $\mu$mol/L）。
      \item 弃 96 孔板中旧培养基，加入 \textbf{90 $\mu$L 新鲜完全培养基 + 10 $\mu$L 药物稀释液}（总体积 100 $\mu$L/孔）。
    \end{enumerate}

  \item \textbf{检测（Day 3）：}
    \begin{enumerate}
      \item 加药 \textbf{48 h} 后（亦可设 24 h 和 72 h 多时间点），每孔加入 \textbf{10 $\mu$L CCK-8 试剂}（Dojindo，目录号 CK04）。
      \item 37$^{\circ}$C 培养箱中孵育 \textbf{1--4 h}（依细胞密度和 CCK-8 批次调整，OD 值目标范围 0.5--1.5）。
      \item 用酶标仪测量 \textbf{450 nm 吸光度}（参考波长 650 nm），记录数值。
    \end{enumerate}

  \item \textbf{IC50 计算（GraphPad Prism 9.0）：}
    \begin{enumerate}
      \item 计算各浓度的细胞抑制率（Inhibition Rate, IR\%）：
        \[ \text{IR\%} = \left(1 - \frac{\text{OD}_{\text{drug}} - \text{OD}_{\text{blank}}}{\text{OD}_{\text{control}} - \text{OD}_{\text{blank}}}\right) \times 100\% \]
        其中：OD\textsubscript{blank} = 无细胞+培养基+CCK-8 孔的 OD 值。
      \item 打开 \software{GraphPad Prism 9.0}，新建「XY」类型数据表。
      \item X 轴输入药物浓度的 log10 值（如 $\log_{10}(100) = 2$），Y 轴输入对应 IR\%。
      \item 点击菜单 \clickpath{Analyze → Fit a Curve → [Inhibitor] vs. Response -- Variable Slope (4PL)}。
      \item 从拟合结果中读取 \textbf{IC50 值}（附 95\% 可信区间），记录于实验记录本。
    \end{enumerate}
\end{enumerate}

\begin{qcbox}
\textbf{CCK-8 实验质控合格标准：}
\begin{itemize}
  \item 空白孔 OD 值（无细胞）$<$ 0.1 \ding{52}
  \item 对照组（DMSO 0.1\%）OD 值 $>$ 0.5，确认细胞生长状态正常 \ding{52}
  \item DMSO 对照组与完全培养基对照组无显著差异（t 检验 P $>$ 0.05）\ding{52}
  \item R\textsuperscript{2}（剂量-效应曲线拟合优度）$\geq$ 0.95 \ding{52}
  \item 实验重复 $\geq$ \textbf{3 次独立生物学重复}（3 次实验的 IC50 值变异系数 CV $\leq$ 20\%）\ding{52}
  \item 顺铂阳性对照 IC50 值在历史参考范围内（CNE1: 2--10 $\mu$mol/L）\ding{52}
\end{itemize}
\end{qcbox}

\subsubsection{流式细胞术检测细胞凋亡与细胞周期}

\paragraph{细胞凋亡检测（Annexin V/PI 双染法）：}

\begin{enumerate}
  \item \textbf{细胞处理：}在 6 孔板中培养 CNE1/CNE2，汇合率 60--70\% 时，加入候选成分（IC50 和 2$\times$IC50 浓度），培养 \textbf{24--48 h}。
  \item \textbf{收集细胞：}弃培养基，加入 0.25\% 胰蛋白酶消化 $\leq$ 2 min，加培养基终止，收集消化液与漂浮细胞（包含早期凋亡细胞），合并转入 1.5 mL EP 管，300 $\times$ g，5 min 离心。
  \item \textbf{染色步骤：}
    \begin{enumerate}
      \item 弃上清，用 \textbf{1$\times$ 结合缓冲液（Binding Buffer）500 $\mu$L} 重悬。
      \item 加入 \textbf{Annexin V-FITC 5 $\mu$L}（BioLegend 或等效），轻混，室温避光孵育 \textbf{15 min}。
      \item 加入 \textbf{PI（碘化丙啶）溶液（50 $\mu$g/mL）5 $\mu$L}，轻混，立即上机检测。
    \end{enumerate}
  \item \textbf{流式仪器参数：}
    \begin{itemize}
      \item FITC 通道激发光：488 nm，发射波长：$\sim$530 nm
      \item PI 通道激发光：488 nm，发射波长：$\sim$617 nm
      \item 每个样品采集 $\geq$ \textbf{$1 \times 10^4$} 个细胞事件
    \end{itemize}
  \item \textbf{结果分析：}
    \begin{itemize}
      \item 早期凋亡：Annexin V\textsuperscript{+}/PI\textsuperscript{-}（右下象限）
      \item 晚期凋亡/坏死：Annexin V\textsuperscript{+}/PI\textsuperscript{+}（右上象限）
      \item 存活细胞：Annexin V\textsuperscript{-}/PI\textsuperscript{-}（左下象限）
      \item \textbf{总凋亡率 = 早期凋亡率 + 晚期凋亡率}
    \end{itemize}
\end{enumerate}

\paragraph{细胞周期检测（PI 单染法）：}

\begin{enumerate}
  \item 药物处理 CNE1/CNE2 细胞 \textbf{24--48 h}，收集细胞（同上），冷 PBS 洗涤 2 次。
  \item 用 \textbf{70\% 冰乙醇（$-$20$^{\circ}$C 预冷）}以 5:1 比例滴加固定细胞，轻混，\textbf{$-$20$^{\circ}$C 固定 $\geq$ 24 h}（可延长至 1 周）。
  \item 固定完成后，300 $\times$ g 离心 5 min，弃乙醇，PBS 洗涤 2 次（洗去乙醇）。
  \item 加入 \textbf{RNase A（100 $\mu$g/mL）50 $\mu$L}，37$^{\circ}$C 水浴 \textbf{30 min}（降解 RNA，消除干扰）。
  \item 加入 \textbf{PI 溶液（50 $\mu$g/mL）450 $\mu$L}，37$^{\circ}$C 或室温避光孵育 \textbf{30 min}。
  \item 用 $300 \mu$m（或更细）滤网过滤单细胞悬液后立即上机，激发波长 488 nm，采集 $\geq 2 \times 10^4$ 个事件。
  \item 使用 \software{FlowJo 10.9}（或等效软件）的 Dean-Jett-Fox 或 Watson（多项式）模型拟合细胞周期各期（G1、S、G2/M）比例。
\end{enumerate}

\subsection{免疫激活活性验证实验}
\label{sec:module6_immune}

\subsubsection{Jurkat T 细胞与 THP-1 巨噬细胞培养 SOP}

\paragraph{Jurkat T 细胞培养（悬浮细胞系）：}

\begin{enumerate}
  \item \textbf{培养基：}RPMI-1640 + 10\% FBS + 1\% Pen/Strep。
  \item \textbf{培养密度控制：}\textbf{$1 \times 10^5$--$1 \times 10^6$ 个/mL}（低于 $5 \times 10^4$ 或高于 $2 \times 10^6$ 均影响细胞状态）。
  \item 每 \textbf{2--3 天传代一次}：300 $\times$ g 离心 5 min，弃上清，用新鲜培养基重悬至初始密度。
  \item 注意：Jurkat 细胞为悬浮细胞，禁止加胰蛋白酶消化。
\end{enumerate}

\paragraph{THP-1 巨噬细胞诱导分化 SOP：}

\begin{enumerate}
  \item THP-1（人单核细胞系）培养条件同 Jurkat，密度范围 $2 \times 10^5$--$8 \times 10^5$ 个/mL。
  \item \textbf{M0 巨噬细胞分化：}将 THP-1 以 \textbf{$5 \times 10^5$ 个/mL} 接种于 6 孔板，加入 \textbf{佛波酯（PMA，Phorbol 12-myristate 13-acetate）100 ng/mL}（Sigma，P8139），刺激 \textbf{48 h}，细胞贴壁即完成 M0 分化。
  \item \textbf{M1 极化（经典激活型，评估中药免疫激活效果）：}
    \begin{itemize}
      \item M0 基础上加入 \textbf{LPS（脂多糖）100 ng/mL} + \textbf{IFN-$\gamma$ 20 ng/mL}，处理 \textbf{24 h}（阳性对照）。
      \item 候选中药成分处理组：加入候选成分（IC50 的 1/4 浓度，避免细胞毒干扰），与 LPS/IFN-$\gamma$ 对照平行设置。
    \end{itemize}
  \item \textbf{M1 极化标志物检测（流式细胞术）：}
    \begin{itemize}
      \item 收集细胞，PBS 洗涤，加入 Fc 受体封闭剂（5 min，防止非特异结合）。
      \item 加入 \textbf{Anti-CD86-APC（克隆 IT2.2，BioLegend，目录号 305412）2 $\mu$L/管}（$1 \times 10^5$ 细胞）。
      \item 4$^{\circ}$C 避光孵育 \textbf{30 min}，PBS 洗涤 2 次，流式检测 CD86\textsuperscript{+}\% 比例。
      \item M1 极化合格判定：LPS/IFN-$\gamma$ 组 CD86\textsuperscript{+}\% $\geq$ 70\%（相比 M0 基线）。
    \end{itemize}
\end{enumerate}

\subsubsection{ELISA 检测细胞因子分泌水平 SOP}

检测细胞因子：IFN-$\gamma$、TNF-$\alpha$、IL-2、IL-6（重要性依次排列）。

\begin{enumerate}
  \item \textbf{样品收集：}
    \begin{enumerate}
      \item 药物处理结束后，收集细胞培养上清液（取 300--400 $\mu$L，避免吸到细胞沉淀）。
      \item 1000 $\times$ g 离心 5 min，去除细胞碎片，上清分装至 EP 管。
      \item 短期使用（$\leq$ 3 天）：4$^{\circ}$C 保存；长期：$-$80$^{\circ}$C 冻存（避免反复冻融）。
    \end{enumerate}

  \item \textbf{ELISA 操作（以 IFN-$\gamma$ ELISA Kit 为例，Human IFN-$\gamma$ ELISA Kit，Invitrogen，目录号 BMS228）：}
    \begin{enumerate}
      \item 按说明书配制标准品系列（7 个浓度梯度），通常从最高浓度（如 1000 pg/mL）按 2 倍稀释。
      \item 将包被好捕获抗体的 96 孔板从 4$^{\circ}$C 取出，每孔加入 \textbf{100 $\mu$L PBS（含 1\% BSA）封闭液}，室温孵育 \textbf{1 h}（若厂家 Kit 已预封闭可跳过此步）。
      \item 洗板（5 次，每次加 PBST 300 $\mu$L，拍干）。
      \item 加标准品与样品（\textbf{100 $\mu$L/孔}），设双复孔，室温孵育 \textbf{2 h}（或 4$^{\circ}$C 过夜）。
      \item 洗板 5 次，加入检测抗体（\textbf{100 $\mu$L/孔}），室温孵育 \textbf{1 h}。
      \item 洗板 5 次，加入 \textbf{HRP-链亲和素（HRP-Streptavidin）100 $\mu$L/孔}，室温孵育 \textbf{30 min}。
      \item 洗板 7 次（此步洗涤要彻底，否则背景高），加入 \textbf{TMB 底物 100 $\mu$L/孔}，避光室温孵育 \textbf{10--15 min}（至蓝色出现）。
      \item 加入 \textbf{终止液（2 M H$_2$SO$_4$）100 $\mu$L/孔}，溶液变黄，立即在酶标仪 \textbf{450 nm} 读取吸光值（参考波长 570 nm）。
    \end{enumerate}

  \item \textbf{结果计算：}
    \begin{enumerate}
      \item 以标准品浓度（pg/mL）为 X 轴，OD$_{450}$ 为 Y 轴，用 \software{GraphPad Prism 9.0} 拟合标准曲线（四参数拟合）。
      \item 将样品 OD 值代入标准曲线，读取对应细胞因子浓度（pg/mL）。
      \item 统计分析：使用 \textbf{单因素方差分析（One-way ANOVA）+ Tukey 事后检验}，比较各组差异。
    \end{enumerate}
\end{enumerate}

\begin{warningbox}
\begin{itemize}
  \item \textbf{对照组浓度预估：}LPS/IFN-$\gamma$ 刺激的 THP-1 的 IL-6 分泌量一般在 1,000--10,000 pg/mL，无刺激对照约在 50--200 pg/mL。样品的稀释倍数需根据预实验结果调整，确保所有样品 OD 值落在标准曲线线性区间内。
  \item 若样品 OD 值超出标准曲线上限，需按 2$\times$、5$\times$ 稀释后重测，并在结果中乘以稀释倍数。
  \item ELISA 板各步骤洗涤需彻底，尤其是最后一次 TMB 加底物前，否则背景高会降低实验灵敏度。
\end{itemize}
\end{warningbox}

% ============================================================
% 新增：T细胞系统全套SOP（Jurkat预实验 + 小鼠脾脏主方案 + 人PBMC临床方案）
% ============================================================
\subsubsection{Jurkat T 细胞预实验方案（技能培训与快速预验证）}
\label{sec:jurkat_preexp}

\begin{infobox}{Jurkat 预实验定位说明}
\textbf{推荐用途（三种场景）：}（1）\textbf{新手技能培训}——Jurkat 为永生化 T 细胞系，无需原代分离，操作简便，适合零基础研究者快速掌握流式细胞术、ELISA 等核心技能；（2）\textbf{快速预验证}——在原代 T 细胞实验前，用 Jurkat 快速评估候选成分对 T 细胞活化标志物的影响方向，为原代实验方案设计提供参考；（3）作为\textbf{补充数据}，与原代 T 细胞实验平行报告（SCI 论文可接受 Jurkat 补充数据）。\\
\textbf{局限性说明：}Jurkat 细胞缺失 TCR$\zeta$链相关信号组分，对部分免疫检查点调控药物的响应与原代 T 细胞存在差异，实验结果须经原代 T 细胞验证后方可用于论文主要结论。
\end{infobox}

\begin{enumerate}
  \item \textbf{Jurkat 培养条件：}RPMI-1640 + 10\% FBS + 1\% Pen/Strep，密度维持在 $1 \times 10^5$--$1 \times 10^6$ 个/mL，37$^{\circ}$C，5\% CO$_2$（详见 \ref{sec:module6_immune}）。
  \item \textbf{PMA + Ionomycin 激活方案（标准预验证流程）：}
  \begin{enumerate}
    \item 取 Jurkat 细胞（$5 \times 10^5$ 个/mL），加入 \textbf{佛波酯（PMA）10 ng/mL} + \textbf{Ionomycin 1 $\mu$g/mL}（Sigma，I9657），激活 \textbf{6--24 h}，诱导 IL-2 分泌与 CD69 表达上调。
    \item 候选成分与 PMA + Ionomycin 同步加入，DMSO 终浓度 $\leq$ 0.1\%，浓度梯度设 \textbf{5 个点}（如 1、5、10、25、50 $\mu$g/mL）。
    \item 收集上清：检测 IL-2（Human IL-2 ELISA Kit，Invitrogen，BMS221）。
    \item 收集细胞：检测 CD69\textsuperscript{+}\%（Anti-CD69-FITC，BioLegend，目录号 310904）与 CD25\textsuperscript{+}\%（Anti-CD25-PE，BioLegend，目录号 302606）。
  \end{enumerate}
  \item \textbf{阳性判定标准：}候选成分处理组 IL-2 分泌量较 0.1\% DMSO 对照组升高 \textbf{$\geq$ 30\%}（P $<$ 0.05），且 CD69\textsuperscript{+}\% 升高 $\geq$ 15 个百分点。
  \item \textbf{结果解读与进阶说明：}Jurkat 阳性需进一步用原代 T 细胞（脾脏/PBMC）验证；若 Jurkat 阳性而原代 T 细胞阴性，须结合 PD-1/PD-L1 信号通路分析差异原因（Jurkat 高表达 PD-1，对 PD-1 阻断类成分响应与原代 T 细胞有差异）。
\end{enumerate}

\subsubsection{小鼠脾脏 T 细胞分离、培养与激活 SOP（核心推荐主方案）}
\label{sec:mouse_spleen_tcell}

\begin{infobox}{小鼠脾脏 T 细胞系统概览}
\textbf{核心推荐理由：}（1）单次可获取大量 T 细胞（$\geq 1 \times 10^8$ 个/只），批量实验材料充足；（2）免疫细胞功能完整，接近体内生理状态；（3）与鼠源荷瘤模型（C57BL/6 背景）完全兼容，为体内实验提供一致免疫学背景；（4）与含药血清预实验（大鼠血清）体系相互印证，证据链更完整。\\
\textbf{推荐小鼠背景：}C57BL/6J（或 BALB/c），6--8 周龄，同性别同批次使用（减少个体差异）。
\end{infobox}

\paragraph{A. 小鼠脾脏 T 细胞分离 SOP：}

\textbf{所需材料：}C57BL/6 小鼠脾脏（新鲜，颈椎脱臼处死）、细胞过滤网（70 $\mu$m，Falcon，目录号 352350）、5 mL 注射器活塞、ACK 裂解液（0.15 M NH$_4$Cl + 10 mM KHCO$_3$ + 0.1 mM Na$_2$EDTA，pH 7.2--7.4，0.2 $\mu$m 过滤除菌）、RPMI-1640 完全培养基。

\begin{enumerate}
  \item \textbf{取脾（须在伦理委员会批准方案下执行）：}颈椎脱臼处死小鼠，75\% 酒精全身擦拭（30 s），无菌条件下剪取脾脏（位于胃左上方，暗红色卵圆形器官约 0.1 g），置于含预冷 PBS 的 60 mm 培养皿中。
  \item \textbf{制备单细胞悬液：}70 $\mu$m 细胞过滤网置于 15 mL 离心管口，用 5 mL 注射器活塞研磨脾脏（螺旋式施压），PBS 反复冲洗过滤网直至脾脏变白（共约 5 mL/只）。
  \item \textbf{红细胞裂解（ACK 法）：}\textbf{300 $\times$ g，5 min} 离心，弃上清；加入 \textbf{1 mL ACK 裂解液}，室温作用 \textbf{2 min}（计时严格，超时会损伤 T 细胞），立即加 \textbf{10 mL PBS} 终止，\textbf{300 $\times$ g，5 min} 离心，弃上清。
  \item \textbf{清洗重悬：}PBS 洗涤 2 次，用 RPMI-1640 完全培养基重悬，台盼蓝计数，细胞活力须 \textbf{$\geq$ 90\%}。
  \item \textbf{T 细胞磁珠纯化（选做）：}若需纯度 $>$ 90\% 的 CD3\textsuperscript{+} T 细胞，使用 \textbf{Miltenyi Pan T Cell Isolation Kit（小鼠，目录号 130-095-130）}按说明书操作（约 30 min，负向磁珠分选，去除非 T 细胞群）；若使用总脾细胞（T 细胞约占 25--40\%）可直接使用，操作更简便但需在分析时考虑细胞混合背景。
  \item \textbf{接种密度：}总脾细胞：\textbf{$2 \times 10^6$ 个/mL}；纯化 T 细胞：\textbf{$1 \times 10^6$ 个/mL}，37$^{\circ}$C，5\% CO$_2$。
\end{enumerate}

\paragraph{B. 小鼠 T 细胞激活与中药活性检测 SOP：}

\begin{enumerate}
  \item \textbf{ConA 激活法（快速标准方案，48--72 h）：}加入 \textbf{刀豆蛋白 A（ConA，Sigma，C0412）终浓度 5 $\mu$g/mL}（总脾细胞）或 \textbf{2 $\mu$g/mL}（纯化 T 细胞），激活 \textbf{48--72 h}，诱导 T 细胞增殖与 IL-2/IFN-$\gamma$ 分泌。
  \item \textbf{Anti-CD3$\varepsilon$/CD28 板结合法（推荐机制研究首选）：}
  \begin{enumerate}
    \item 96 孔板预包被：每孔加入 \textbf{Anti-CD3$\varepsilon$（克隆 145-2C11，Bio X Cell，BE0001-1）50 $\mu$L，浓度 2 $\mu$g/mL}，4$^{\circ}$C 过夜。弃液后用 PBS 洗涤 2 次。
    \item 接种脾细胞（$2 \times 10^5$ 个/孔），加入可溶性 \textbf{Anti-CD28（克隆 37.51，Bio X Cell，BE0015-1）2 $\mu$g/mL}，37$^{\circ}$C，5\% CO$_2$，孵育 \textbf{72 h}。
  \end{enumerate}
  \item \textbf{候选成分加药方案：}与激活体系同步加入候选成分（DMSO $\leq$ 0.1\%，浓度梯度参考 CCK-8 IC50 的 1/8、1/4、1/2 倍），平行设置溶剂对照（0.1\% DMSO）与阳性对照（重组 IL-2，50 U/mL）。
  \item \textbf{结果检测指标：}
  \begin{itemize}
    \item \textbf{ELISA：}收集上清（激活 48 h 后），检测 IL-2（Mouse IL-2 ELISA Kit，eBioscience，BMS601）和 IFN-$\gamma$（Mouse IFN-$\gamma$ ELISA Kit，eBioscience，BMS606）。
    \item \textbf{流式细胞术：}CD4\textsuperscript{+}/CD8\textsuperscript{+} T 细胞活化标志物——CD69（克隆 H1.2F3，BioLegend）、CD25（克隆 PC61，BioLegend）、PD-1（克隆 29F.1A12，BioLegend）。
    \item \textbf{T 细胞增殖（CFSE 稀释法）：}激活前用 \textbf{CFSE 终浓度 5 $\mu$M}（37$^{\circ}$C，10 min）预标记脾细胞，PBS 洗涤 3 次，激活后 72 h 流式检测 CFSE 荧光稀释程度（每分裂一次荧光减半，可直接计算增殖指数 PI 值）。
  \end{itemize}
  \item \textbf{阳性判定标准：}候选成分组 IL-2 分泌量较 DMSO 对照升高 \textbf{$\geq$ 30\%}（P $<$ 0.05），或 CD69\textsuperscript{+}CD8\textsuperscript{+}\% 升高 $\geq$ 20 个百分点。
\end{enumerate}

\begin{qcbox}
\textbf{小鼠脾脏 T 细胞实验质控合格标准：}
\begin{itemize}
  \item 新鲜脾脏细胞活力 $\geq$ 90\%（台盼蓝，分离后 2 h 内完成加药）\ding{52}
  \item ConA/Anti-CD3 阳性激活对照 IL-2 分泌量较无刺激对照升高 $\geq$ 5 倍（激活有效性）\ding{52}
  \item 每次实验使用 $\geq$ 3 只不同小鼠（每只为 1 次生物重复，$n \geq 3$）\ding{52}
  \item DMSO 对照组 T 细胞活力 $\geq$ 85\%（确认溶剂无免疫毒性）\ding{52}
  \item 实验全程在超净台内操作，分离后 24 h 内完成加药 \ding{52}
\end{itemize}
\end{qcbox}

\begin{warningbox}
\textbf{Plan B --- 原代 T 细胞大量死亡的分层应急方案}

\textbf{风险描述：}原代 T 细胞（PBMC/脾脏 T 细胞）在分离、激活或加药过程中，因操作失误、DMSO 浓度超标、FBS 批次问题或候选成分本身毒性，可能出现大量死亡（活力 $<$ 70\%），导致整批实验数据无效。

\textbf{应急处置步骤（须在 24 h 内完成评估）：}
\begin{enumerate}
  \item \textbf{立即核查（0--2 h）：}台盼蓝染色确认死亡率；对照 CCK-8 数据判断是否为候选成分本身毒性（若成分 IC50 $<$ 5 $\mu$g/mL，则 T 细胞毒性窗口极窄，须将加药浓度降至 1 $\mu$g/mL 以下重新实验）。
  \item \textbf{DMSO 溶剂核查（0--1 h）：}单独设置 DMSO 梯度浓度（0.05\%、0.1\%、0.2\%、0.5\%）T 细胞培养实验，确认当前批次 FBS + 培养基中 DMSO $\leq$ 0.1\% 对 T 细胞无毒；若 0.1\% DMSO 即可导致 T 细胞死亡，须更换 FBS 批次（T 细胞对 FBS 批次敏感性远高于肿瘤细胞）。
  \item \textbf{FBS 批次排查（1--3 天）：}对新批次 FBS 进行 T 细胞 72 h 存活率测试（活力 $\geq$ 85\% 为合格），更换合格 FBS 后重复实验。
  \item \textbf{临时切换至 Jurkat 细胞系（应急替代，2--3 天）：}若原代 T 细胞问题持续无法排查，临时使用 Jurkat 完成免疫激活验证初步数据，获取实验方向后再择时重复原代实验；论文方法部分须说明「Jurkat 为过渡替代验证，原代 T 细胞实验为论文主要数据」。
  \item \textbf{延迟实验 + 重新取材（1--2 周）：}记录失败原因（科研日志），待下一批实验材料（小鼠或供体血液）到位后重复；向导师汇报并调整时间节点。
\end{enumerate}
\end{warningbox}

\subsubsection{人 PBMC T 细胞分离、培养与激活 SOP（临床转化佐证方案）}
\label{sec:pbmc_tcell}

\begin{infobox}{人 PBMC T 细胞系统概览}
\textbf{系统定位（临床转化辅助方案）：}人外周血单个核细胞（Peripheral Blood Mononuclear Cells，PBMC）中的 T 细胞（CD3\textsuperscript{+}，占 PBMC 60--80\%）是体外评估中药免疫激活效果临床相关性最强的细胞模型，实验结果可直接外推至人体免疫应答。\\
\textbf{使用场景：}（1）需要提升 SCI 论文临床转化说服力时（IF $\geq$ 5.0 期刊通常要求人源数据）；（2）机制研究中需对比鼠源与人源 T 细胞差异时。\\
\textbf{操作前提：}须提前 4 周以上协调伦理审批与供体招募；须通过院校伦理委员会审批（血液采集知情同意）；确认供体无传染性疾病（HIV、HBV、HCV）。
\end{infobox}

\paragraph{A. 人 PBMC 分离 SOP（密度梯度离心法）：}

\textbf{所需试剂：}新鲜抗凝全血（EDTA 管，采集后 24 h 内处理）、Ficoll-Paque PLUS（GE Healthcare，比重 1.077 g/mL，目录号 17144002）、PBS（无钙镁，pH 7.4，室温）、RPMI-1640 完全培养基。

\begin{enumerate}
  \item \textbf{全血稀释：}取全血与等量 PBS 混合（1:1），使用室温 PBS（禁用低温 PBS，防止细胞裂解）。
  \item \textbf{铺层加样：}15 mL 离心管底部缓慢加入 \textbf{3 mL Ficoll-Paque PLUS}（沿管壁，一次性吸管操作），用 5 mL 吸管沿管壁缓慢将稀释全血（6 mL）铺于 Ficoll 液面上方（切勿破坏界面）。
  \item \textbf{密度梯度离心：}室温，\textbf{400 $\times$ g，30 min，无制动（Brake = 0）}，低加速（Acceleration = 2）。
  \item \textbf{收集 PBMC 白膜层：}离心后自上而下：血浆层→PBMC 白膜层（白色环状）→Ficoll 层→红细胞层。用吸管小心吸取白膜层（约 1--2 mL），转入新 15 mL 管。
  \item \textbf{双次清洗去除血小板：}加 \textbf{10 mL PBS}，\textbf{300 $\times$ g，10 min} 离心，弃上清（第一次）；再加 10 mL PBS，\textbf{200 $\times$ g，10 min} 离心，弃上清（第二次洗涤去除血小板，使用较低离心力）。
  \item \textbf{细胞计数与质控：}PBS 重悬至 1 mL，台盼蓝计数，细胞活力须 \textbf{$\geq$ 90\%}；以 \textbf{$2 \times 10^6$ 个/mL} 接种，37$^{\circ}$C，5\% CO$_2$，2 h 静置（可选：使单核细胞贴壁，去除后获取淋巴细胞富集群）。
\end{enumerate}

\paragraph{B. 人 PBMC T 细胞激活方案：}

\begin{enumerate}
  \item \textbf{PHA 激活法（快速经济，适合初步验证，72 h 方案）：}在 PBMC 中加入 \textbf{植物血凝素（PHA，Sigma，L1668）终浓度 5 $\mu$g/mL} + \textbf{重组 IL-2 100 U/mL}，刺激 72 h（CD69、CD25 显著上调）。
  \item \textbf{Anti-CD3/CD28 法（最接近生理激活，推荐机制研究）：}
  \begin{enumerate}
    \item 96 孔板预包被 \textbf{Anti-CD3（OKT3，eBioscience，16-0037-85）2 $\mu$g/mL}，4$^{\circ}$C 过夜；弃液 PBS 洗 2 次。
    \item 每孔加 PBMC \textbf{$2 \times 10^5$ 个（100 $\mu$L）}，同时加入 \textbf{可溶性 Anti-CD28（CD28.2，eBioscience，16-0289-85）0.5 $\mu$g/mL} + IL-2（100 U/mL），37$^{\circ}$C，72 h。
    \item 候选成分与激活剂同步加入，检测 IL-2/IFN-$\gamma$ ELISA 及流式活化标志物。
  \end{enumerate}
  \item \textbf{供体差异性控制（关键质控要求）：}须使用 \textbf{$\geq$ 3 位不同供体}的 PBMC（每位供体为 1 次生物重复，$n \geq 3$），数据以各供体独立实验方式汇总，报告均值 $\pm$ SD（供体内以技术重复计算，供体间以生物重复计算）。
\end{enumerate}

\begin{warningbox}
\begin{itemize}
  \item \textbf{时效性限制：}新鲜全血须在采集后 \textbf{24 h 内处理完毕}，存放时间延长将导致 T 细胞激活响应显著下降，数据不可靠。
  \item \textbf{伦理合规：}使用人血须经院校伦理委员会批准，所有供体签署知情同意书，采集前确认 HIV、HBV、HCV 阴性（供体健康筛查报告存档）。
  \item \textbf{供体间变异：}个体间 T 细胞对刺激的响应可能存在显著差异，须 $\geq$ 3 位供体方可作统计结论；若仅 1--2 位供体，数据仅作探索性分析，不作为论文主要结论。
\end{itemize}
\end{warningbox}

% ============================================================
% 新增：共培养体系SOP（Transwell间接 + 直接共培养）
% ============================================================
\subsubsection{免疫激活-抗肿瘤协同效应共培养体系 SOP}
\label{sec:coculture}

\begin{infobox}{共培养体系科学问题与实验设计}
\textbf{核心科学问题：}候选中药成分通过「直接细胞毒性」杀伤肿瘤细胞，还是通过「激活免疫细胞间接杀伤」肿瘤，抑或两者协同发挥作用？共培养体系可定量区分两种效应的相对贡献，为「免疫激活协同抗肿瘤」核心科学主张提供直接细胞水平实验证据。\\
\textbf{两套方案概览：}（1）\textbf{Transwell 间接共培养}（首选，操作简便，可分离细胞因子效应与接触效应）；（2）\textbf{直接共培养}（选做，验证细胞接触依赖性杀伤，数据解读需流式区分两种细胞）。
\end{infobox}

\paragraph{A. Transwell 间接共培养方案（推荐首选）：}

\textbf{实验原理：}Transwell 聚碳酸酯膜（0.4 $\mu$m 孔径）允许细胞因子、小分子药物等自由通过，但\textbf{阻断细胞-细胞物理接触}，从而将「旁分泌细胞因子介导」与「直接接触依赖」效应分离，实现定量区分。

\textbf{所需材料：}24 孔板配套 Transwell 小室（聚碳酸酯膜，孔径 \textbf{0.4 $\mu$m}，Corning，目录号 3396）；CNE2 细胞（下室）；小鼠脾脏 T 细胞或 PBMC T 细胞（上室）。

\begin{enumerate}
  \item \textbf{实验分组设计（4 组对照，$n \geq 3$）：}

  \begin{center}
  \begin{tabular}{p{3.2cm}p{3.5cm}p{3.5cm}p{2.5cm}}
  \toprule
  \textbf{实验组} & \textbf{下室（CNE2）} & \textbf{上室（免疫细胞）} & \textbf{加药部位} \\
  \midrule
  单药直接杀伤组 & CNE2 + 候选成分 & 培养基（无细胞） & 仅下室 \\
  \addlinespace
  免疫介导杀伤组 & CNE2（无药） & T 细胞 + 候选成分 & 仅上室 \\
  \addlinespace
  协同杀伤组 & CNE2 + 候选成分 & T 细胞 + 候选成分 & 上下室均加 \\
  \addlinespace
  溶剂对照组 & CNE2 + 0.1\% DMSO & T 细胞 + 0.1\% DMSO & 上下室均加 \\
  \bottomrule
  \end{tabular}
  \end{center}

  \item \textbf{操作步骤：}
  \begin{enumerate}
    \item \textbf{Day 0（铺下室）：}CNE2 细胞以 \textbf{$5 \times 10^4$ 个/孔}（500 $\mu$L/孔）接种于 24 孔板下室，37$^{\circ}$C 过夜贴壁（12--16 h）。
    \item \textbf{Day 1（放入 Transwell 小室 + 加药）：}将 Transwell 小室放入各孔。上室加入脾脏 T 细胞（\textbf{$1 \times 10^6$ 个，200 $\mu$L}）或 PBMC（$2 \times 10^5$ 个，200 $\mu$L）。按分组方案在上室/下室分别加入候选成分（浓度设 \textbf{1/4 IC50} 和 \textbf{1/2 IC50} 两个浓度，各组平行操作）。
    \item \textbf{Day 1--3（孵育 48 h）：}37$^{\circ}$C，5\% CO$_2$，静置培养，勿频繁移动培养板（避免界面扰动）。
    \item \textbf{Day 3（检测）：}取出 Transwell 小室。
    \begin{itemize}
      \item \textbf{下室（CNE2 存活率）：}弃去上室培养基，每孔加 \textbf{100 $\mu$L CCK-8 工作液}（10 $\mu$L CCK-8 + 90 $\mu$L 培养基），孵育 2 h，酶标仪 450 nm 读取，计算细胞存活率。
      \item \textbf{上室培养上清（免疫激活指标）：}收集 200 $\mu$L 上清，ELISA 检测 IFN-$\gamma$（最敏感的 T 细胞激活指标）和 IL-2。
    \end{itemize}
  \end{enumerate}

  \item \textbf{结果解读框架：}
  \begin{itemize}
    \item 仅「单药直接杀伤组」CNE2 存活率下降 → 候选成分具备\textbf{直接细胞毒活性}。
    \item 仅「免疫介导杀伤组」CNE2 存活率下降（无药直接接触）→ 候选成分通过\textbf{激活免疫细胞间接抑制肿瘤}（旁分泌细胞因子介导）。
    \item 「协同杀伤组」效果 $>$ 上述两组效果之和（Bliss 超协同）→ \textbf{双机制协同抗肿瘤}（核心创新点支撑证据）。
    \item 上室 IFN-$\gamma$ 升高 + 下室 CNE2 存活率下降同步发生 → 直接证明「免疫激活—抗肿瘤」因果关联。
  \end{itemize}
\end{enumerate}

\paragraph{B. 直接共培养方案（选做，验证细胞接触依赖性杀伤）：}

\textbf{实验原理：}肿瘤细胞与效应 T 细胞直接混合培养，允许物理接触（Perforin/Granzyme B 介导的直接杀伤），配合荧光标记区分两种细胞群，用流式细胞术定量分析肿瘤细胞死亡率。

\begin{enumerate}
  \item \textbf{CNE2 预标记（CFSE 荧光标记，Day 0）：}CNE2 细胞用 \textbf{CFSE（终浓度 5 $\mu$M，37$^{\circ}$C，10 min）}荧光标记，PBS 洗涤 \textbf{3 次}（充分去除游离 CFSE），以 \textbf{$5 \times 10^4$ 个/孔}接种于 24 孔板，过夜贴壁。
  \item \textbf{效靶比（E:T Ratio）设置：}T 细胞（效应细胞）:CNE2（靶细胞）= 5:1、10:1、20:1（至少设 2 个效靶比），按此比例加入 T 细胞。
  \item \textbf{加药与共培养（Day 1）：}在混合培养体系中同步加入候选成分（浓度与 Transwell 方案一致），共培养 \textbf{24 h}。
  \item \textbf{流式检测（Day 2）：}
  \begin{enumerate}
    \item 收集细胞：加入胰酶消化贴壁 CNE2 细胞，与培养基（含悬浮 T 细胞）合并收集至 EP 管。
    \item 加入 \textbf{PI（碘化丙啶，终浓度 1 $\mu$g/mL）}区分死活细胞（PI\textsuperscript{+} = 细胞膜破损的死亡细胞）。
    \item \textbf{门控策略：}先门控 CFSE\textsuperscript{+} 细胞群（CNE2）→ 在此门内统计 PI\textsuperscript{+}（CNE2 死亡率）；CFSE\textsuperscript{-} 细胞群（T 细胞）→ 检测 CD69\textsuperscript{+}（T 细胞活化状态）。
  \end{enumerate}
  \item \textbf{特异性杀伤率计算：}
  \[
  \text{特异性杀伤率 (\%)} = \frac{\text{PI}^{+}\text{CFSE}^{+}\%_{\text{共培养}} - \text{PI}^{+}\text{CFSE}^{+}\%_{\text{仅 CNE2 对照}}}{100\% - \text{PI}^{+}\text{CFSE}^{+}\%_{\text{仅 CNE2 对照}}} \times 100\%
  \]
\end{enumerate}

\begin{qcbox}
\textbf{共培养体系质控合格标准：}
\begin{itemize}
  \item Transwell 方案：溶剂对照（0.1\% DMSO）CNE2 存活率 $\geq$ 90\%（溶剂无细胞毒性）\ding{52}
  \item Transwell 方案：免疫激活阳性对照（ConA/PHA）上清 IFN-$\gamma$ 升高 $\geq$ 5 倍（T 细胞有效激活）\ding{52}
  \item 直接共培养：CFSE 标记 CNE2 标记率 $\geq$ 95\%（流式验证），标记后细胞活力 $\geq$ 90\% \ding{52}
  \item 所有共培养实验 $n \geq 3$ 次独立生物重复 \ding{52}
  \item 协同效应定量分析须使用 \textbf{Bliss 独立性模型}（可用 SynergyFinder 软件，\texttt{https://synergyfinder.fimm.fi/}，免费在线工具）\ding{52}
\end{itemize}
\end{qcbox}

\begin{tipbox}
「科研进阶技巧：」若希望直接证明 CD8\textsuperscript{+} 细胞毒 T 细胞（CTL）的特异性杀伤作用，可在直接共培养实验前，用磁珠纯化 \textbf{CD8\textsuperscript{+} T 细胞}（Miltenyi CD8$\alpha$\textsuperscript{+} T Cell Isolation Kit，小鼠，目录号 130-104-075），与 CNE2 共培养，获得高纯度 CTL 杀伤数据。此数据直接支持「候选成分激活 CTL 特异性清除肿瘤细胞」的核心机制假说，具有显著的 SCI 发表价值（可支撑 IF $\geq$ 6.0 期刊投稿）。
\end{tipbox}

\subsection{分子机制验证实验}
\label{sec:module6_mechanism}

\subsubsection{qRT-PCR 检测核心靶点 mRNA 表达水平}

\begin{enumerate}
  \item \textbf{RNA 提取（TRIzol 法）：}
    \begin{enumerate}
      \item 弃去 6 孔板培养基，每孔加入 \textbf{TRIzol 试剂（Invitrogen）1 mL}，室温静置 5 min，用移液枪充分裂解细胞，转入 1.5 mL EP 管。
      \item 加入 \textbf{氯仿 200 $\mu$L}，剧烈振荡混匀 15 s，室温静置 3 min，\textbf{12,000 $\times$ g，4$^{\circ}$C，15 min} 离心。
      \item 小心吸取上层水相（约 400--450 $\mu$L），转入新 EP 管，禁止吸到中间层（蛋白质层）与下层（有机相）。
      \item 加入等体积 \textbf{异丙醇 400 $\mu$L}，充分混匀，室温静置 10 min，12,000 $\times$ g，4$^{\circ}$C，10 min 离心，弃上清。
      \item 加入 \textbf{75\% 乙醇（DEPC 水配制）1 mL}，7,500 $\times$ g，4$^{\circ}$C，5 min 离心洗涤，弃乙醇，室温风干 \textbf{5--10 min}（不可过干）。
      \item 加入 \textbf{DEPC 水 30--50 $\mu$L} 溶解 RNA，用 NanoDrop 检测 OD$_{260/280}$（目标值：1.8--2.1）和 RNA 浓度。
    \end{enumerate}

  \item \textbf{逆转录（cDNA 合成）：}
    使用 HiScript III RT SuperMix（Vazyme，R323，或等效产品），按说明书操作，通常反应体积 20 $\mu$L，42$^{\circ}$C，15 min；85$^{\circ}$C，10 s 灭活。

  \item \textbf{qPCR 反应：}
    \begin{enumerate}
      \item 使用 \software{CFX96 实时荧光定量 PCR 仪}（BioRad）或等效仪器，反应体系 20 $\mu$L：
        \begin{center}
        \begin{tabular}{lll}
        \toprule
        \textbf{组分} & \textbf{体积} & \textbf{终浓度} \\
        \midrule
        SYBR Green PCR Mix (2$\times$) & 10 $\mu$L & 1$\times$ \\
        正向引物（10 $\mu$mol/L） & 0.8 $\mu$L & 0.4 $\mu$mol/L \\
        反向引物（10 $\mu$mol/L） & 0.8 $\mu$L & 0.4 $\mu$mol/L \\
        cDNA 模板（10--50 ng/$\mu$L） & 2 $\mu$L & -- \\
        无核酸酶水 & 6.4 $\mu$L & -- \\
        \bottomrule
        \end{tabular}
        \end{center}
      \item \textbf{PCR 程序：}95$^{\circ}$C，3 min；95$^{\circ}$C，10 s → 60$^{\circ}$C，30 s → 72$^{\circ}$C，30 s（40 个循环）；熔解曲线 60$^{\circ}$C → 95$^{\circ}$C（验证产物特异性）。
      \item 内参基因：\textbf{$\beta$-actin（\gene{ACTB}）或 GAPDH（\gene{GAPDH}）}，两者 Ct 值应在 18--24 之间。
      \item 结果计算：使用 $2^{-\Delta\Delta Ct}$ 方法计算相对表达量，以对照组（DMSO）为基准，统计 $\geq$ 3 次独立生物学重复（每次 3 个技术重复）。
    \end{enumerate}
\end{enumerate}

\subsubsection{Western Blot 检测核心蛋白表达}

\begin{enumerate}
  \item \textbf{蛋白提取：}弃培养基，PBS 洗涤，每孔加入 \textbf{RIPA 裂解液（含蛋白酶抑制剂复合物）100--200 $\mu$L}，冰上裂解 15 min，刮细胞，12,000 $\times$ g，4$^{\circ}$C，15 min 离心，收集上清，BCA 法定量（目标浓度：1--5 mg/mL）。
  \item \textbf{SDS-PAGE 电泳：}配制相应浓度分离胶（目标蛋白 MW 决定，如 TP53 $\sim$ 53 kDa 用 10\% 胶），每泳道上样 \textbf{20--30 $\mu$g 蛋白}。80 V，30 min（浓缩胶）→ 120 V，60--90 min（分离胶）。
  \item \textbf{转膜：}湿转法，100 V，60 min（或恒流 300 mA，45 min），转至 PVDF 膜（预先甲醇激活 30 s）。
  \item \textbf{封闭与抗体孵育：}5\% 脱脂奶粉（或 5\% BSA，用于磷酸化蛋白）封闭 \textbf{1 h}；一抗 4$^{\circ}$C 孵育过夜（孵育缓冲液：5\% BSA/TBST）；TBST 洗涤 3 $\times$ 10 min；HRP 标记二抗室温孵育 \textbf{1 h}；ECL 化学发光显影。
\end{enumerate}

\begin{qcbox}
\textbf{模块六质控合格标准：}
\begin{itemize}
  \item 细胞系来源有据可查，STR 鉴定结果与数据库匹配 \ding{52}
  \item CCK-8 实验：R\textsuperscript{2} $\geq$ 0.95，IC50 生物重复 CV $\leq$ 20\%，$\geq$ 3 次独立重复 \ding{52}
  \item 凋亡检测：阴性对照凋亡率 $<$ 10\%（细胞状态正常基线）\ding{52}
  \item ELISA：标准曲线 R\textsuperscript{2} $\geq$ 0.99，所有样品 OD 值落在标准曲线线性区间 \ding{52}
  \item qRT-PCR：内参 Ct 值 18--24，扩增效率 95--105\%，熔解曲线单峰 \ding{52}
  \item WB：内参蛋白（$\beta$-actin/GAPDH）条带清晰，各泳道均匀 \ding{52}
  \item 所有实验 $\geq$ 3 次独立生物学重复，结果以「均值 $\pm$ 标准差（SD）」呈现 \ding{52}
\end{itemize}
\end{qcbox}

\moduleoutput{CCK-8 IC50 数据表、凋亡/周期流式散点图、ELISA 细胞因子浓度数据、qRT-PCR 相对表达量图、WB 图，作为论文结果章节（Results）核心实验数据}

% ============================================================
% MODULE 7: 实验结果多维数据分析与可视化规范
% ============================================================
\section{模块七：实验结果多维数据分析与可视化规范}
\label{sec:module7}

\begin{infobox}{模块七概览}
\textbf{操作目的：}规范所有实验数据的统计分析方法与图表绘制标准，确保数据处理符合学术规范，生成可直接用于 SCI 期刊投稿的发表级图表。\\
\textbf{本模块输出：}所有实验数据的统计图（PDF 格式，$\geq$ 300 dpi）、数据统计汇总表（Excel）。
\end{infobox}

\subsection{定量实验数据统计学分析方法}
\label{sec:module7_stats}

\subsubsection{统计方法选择原则}

\begin{center}
\begin{tabular}{p{3.5cm}p{3cm}p{3.5cm}p{4cm}}
\toprule
\textbf{实验类型} & \textbf{数据分布} & \textbf{推荐统计方法} & \textbf{软件操作路径} \\
\midrule
CCK-8 IC50（组间比较）& 正态 & 独立样本 t 检验（两组）& Prism 9.0: Analyze → t tests \\
\addlinespace
细胞凋亡率（多组）& 正态 & 单因素 ANOVA + Tukey & Prism: Analyze → One-way ANOVA \\
\addlinespace
细胞周期（多组）& 正态 & 单因素 ANOVA + Tukey & Prism: Analyze → One-way ANOVA \\
\addlinespace
ELISA 细胞因子 & 正态/偏态 & 先 Shapiro-Wilk 检验；正态→ ANOVA；偏态→ Kruskal-Wallis & Prism: Analyze → Normality \\
\addlinespace
qRT-PCR（$\Delta\Delta$Ct）& 对数正态 & 对 $\Delta$Ct 进行统计分析（t 检验或 ANOVA）& Prism: 对 $\Delta$Ct 数据直接分析 \\
\addlinespace
WB 灰度值分析 & 正态 & 单因素 ANOVA + Dunnett 检验（多组 vs. 对照）& Prism: One-way ANOVA \\
\bottomrule
\end{tabular}
\end{center}

\paragraph{显著性差异符号规范（统一标注）：}
\begin{itemize}
  \item $*$: $p < 0.05$（统计显著）
  \item $**$: $p < 0.01$（高度显著）
  \item $***$: $p < 0.001$（极显著）
  \item $****$: $p < 0.0001$（极极显著）
  \item ns: not significant（$p \geq 0.05$）
\end{itemize}

\subsubsection{GraphPad Prism 9.0 核心操作步骤}

\begin{enumerate}
  \item \textbf{单因素 ANOVA + Tukey 检验（多组细胞实验数据）：}
    \begin{enumerate}
      \item 打开 Prism 9.0，选择「Column」数据表类型。
      \item 在各列填入各组数据（每列为一个处理组，每行为一个生物重复值）。
      \item 点击 \clickpath{Analyze → Compare Means → One-way ANOVA (and nonparametric)}。
      \item 选择「Ordinary One-way ANOVA」，勾选「Assume Gaussian distribution」（正态假设）。
      \item 在「Multiple Comparisons」中选择「Compare the mean of each column with the mean of every other column (Tukey)」（Tukey 全配对事后检验）。
      \item 点击「OK」，查看结果并记录各组对比的 P 值与显著性标注。
    \end{enumerate}
\end{enumerate}

\subsection{生信分析结果可视化规范}
\label{sec:module7_vis}

\subsubsection{火山图绘制规范（差异基因分析）}

\begin{lstlisting}[language=R,caption={发表级火山图绘制代码（ggplot2）}]
# ==== 发表级火山图 ====
library(ggplot2)
library(ggrepel)
library(dplyr)

# 读取差异基因结果
deg_data <- read_csv("results/03_DEGs_GSE12452_NPC_vs_Normal.csv")

# 添加颜色标注分类
deg_data <- deg_data %>%
  mutate(
    label = case_when(
      adj.P.Val < 0.05 & logFC > 1  ~ "Up",
      adj.P.Val < 0.05 & logFC < -1 ~ "Down",
      TRUE                           ~ "NS"
    ),
    # 标注核心 Hub 靶点的基因名
    gene_label = ifelse(Gene.symbol %in% c("TP53","AKT1","EGFR","CD274","STAT3"),
                        Gene.symbol, "")
  )

# 绘制火山图
p <- ggplot(deg_data, aes(x = logFC, y = -log10(adj.P.Val))) +
  # 背景点（非显著）
  geom_point(data = filter(deg_data, label == "NS"),
             color = "gray70", alpha = 0.4, size = 0.8) +
  # 上调基因点
  geom_point(data = filter(deg_data, label == "Up"),
             color = "#E41A1C", alpha = 0.7, size = 1.2) +
  # 下调基因点
  geom_point(data = filter(deg_data, label == "Down"),
             color = "#377EB8", alpha = 0.7, size = 1.2) +
  # 阈值线
  geom_vline(xintercept = c(-1, 1), linetype = "dashed",
             color = "orange", linewidth = 0.5) +
  geom_hline(yintercept = -log10(0.05), linetype = "dashed",
             color = "darkgreen", linewidth = 0.5) +
  # Hub 靶点标注
  geom_label_repel(
    data = filter(deg_data, gene_label != ""),
    aes(label = gene_label),
    size = 3, box.padding = 0.5, max.overlaps = 20,
    fontface = "italic"
  ) +
  # 坐标轴与主题
  labs(
    title = "Differential Gene Expression: NPC vs. Normal (GSE12452)",
    subtitle = paste0("Total DEGs: ", sum(deg_data$label != "NS"),
                      " (Up: ", sum(deg_data$label == "Up"),
                      ", Down: ", sum(deg_data$label == "Down"), ")"),
    x = expression(log[2]~"Fold Change"),
    y = expression(-log[10]~"(Adjusted P-value)")
  ) +
  theme_bw(base_size = 12) +
  theme(
    panel.grid.major = element_line(color = "gray90", linewidth = 0.3),
    panel.grid.minor = element_blank(),
    plot.title = element_text(face = "bold", size = 12),
    legend.position = "none"
  ) +
  # 颜色注释
  annotate("text", x = max(deg_data$logFC)*0.7, y = max(-log10(deg_data$adj.P.Val))*0.95,
           label = paste0("Up: ", sum(deg_data$label == "Up")), color = "#E41A1C", size = 4) +
  annotate("text", x = min(deg_data$logFC)*0.7, y = max(-log10(deg_data$adj.P.Val))*0.95,
           label = paste0("Down: ", sum(deg_data$label == "Down")), color = "#377EB8", size = 4)

ggsave("results/07_volcano_plot.pdf", p, width = 8, height = 7, dpi = 300)
cat("火山图已保存至 results/07_volcano_plot.pdf\n")
\end{lstlisting}

\subsection{实验结果可视化规范}
\label{sec:module7_exp_vis}

\paragraph{SCI 论文配图通用标准：}

\begin{itemize}
  \item \textbf{分辨率：}提交论文时 $\geq$ \textbf{300 dpi}（通常期刊要求：300 dpi for photographs, 600 dpi for line art）
  \item \textbf{格式：}投稿首选 \textbf{PDF 或 TIFF 格式}（EPS 可选）；在线平台可先提交 PDF，修改后提交 TIFF
  \item \textbf{字体：}图内字体统一使用 \textbf{Arial 或 Helvetica}，字号 $\geq$ 8 pt（印刷后清晰可读）
  \item \textbf{颜色：}使用 CMYK 色彩模式（适合印刷），或同时保留 RGB 版本（适合在线浏览）；推荐色盲友好色板（如 ColorBrewer）
  \item \textbf{图例说明：}每图附完整图例（Figure Legend），包含实验组成、重复次数（n = X，$X$ 个生物重复）、统计方法与显著性标注说明
\end{itemize}

\begin{qcbox}
\textbf{模块七质控合格标准：}
\begin{itemize}
  \item 所有统计方法已在论文方法部分注明，P 值阈值 $<$ 0.05 \ding{52}
  \item 正态性检验（Shapiro-Wilk 检验）在 ANOVA 前已执行 \ding{52}
  \item 所有图片分辨率 $\geq$ 300 dpi，字体大小 $\geq$ 8 pt \ding{52}
  \item 火山图、热图、气泡图、KM 曲线已生成 PDF 版本 \ding{52}
\end{itemize}
\end{qcbox}

% ============================================================
% MODULE 8: 研究成果输出与学术撰写框架
% ============================================================
\section{模块八：研究成果输出与学术撰写框架}
\label{sec:module8}

\begin{infobox}{模块八概览}
\textbf{操作目的：}将前七个模块的数据与图表系统整合，输出符合 SCI 期刊/核心期刊标准的学术论文、发明专利申请材料和课题申报书，实现科研成果的全场景转化。
\end{infobox}

\subsection{学术论文撰写标准化框架（IMRaD 结构）}
\label{sec:module8_paper}

\subsubsection{SCI 论文章节结构与撰写要点}

\begin{center}
\begin{tabular}{p{2.5cm}p{4cm}p{3cm}p{4cm}}
\toprule
\textbf{章节} & \textbf{核心内容要点} & \textbf{字数范围} & \textbf{数据来源模块} \\
\midrule
Abstract & 研究目的→方法→关键发现（IC50 数值、靶点名称）→结论，无缩写，最后撰写 & 250--300 words & 全部模块 \\
\addlinespace
Introduction & ①NPC 临床困境（2--3 句）②中药研究现状（引 3--5 篇顶刊）③研究空白→本研究创新点（最关键）④目标与假说 & 500--700 words & 模块一 \\
\addlinespace
Materials \& Methods & ①数据库检索（TCMSP/ETCM，OB/DL 阈值）②靶点筛选（GEO 数据集编号）③网络构建（Cytoscape 3.10.1 版本）④对接（AutoDock Vina 1.2.3，参数）⑤细胞系（CNE1/CNE2，来源，传代代次）⑥CCK-8/流式/ELISA 方法（引试剂盒目录号）⑦统计学（软件版本，P 值阈值） & 800--1200 words & 模块二--六 \\
\addlinespace
Results & 按逻辑顺序：成分筛选→靶点交集→网络分析→富集分析→对接结果→公共数据验证→细胞实验→免疫实验→机制验证 & 1200--1800 words & 模块二--六 \\
\addlinespace
Discussion & ①主要发现总结②机制解释（对比文献）③方法局限性（数据库覆盖率不足、分子对接无法替代实验等）④临床转化潜力 & 800--1000 words & 全部模块 \\
\addlinespace
References & $\geq$ 40 篇；60--70\% 为近 5 年文献；需包含模块工具原文（TCMSP、AutoDock Vina、limma 等） & -- & 全部模块 \\
\bottomrule
\end{tabular}
\end{center}

\subsubsection{目标期刊选择建议}

\begin{center}
\begin{tabular}{llll}
\toprule
\textbf{期刊} & \textbf{IF（2024）} & \textbf{收稿偏好} & \textbf{通讯周期} \\
\midrule
Journal of Ethnopharmacology & $\sim$5.4 & 中药/天然药物，含机制数据 & 3--6 个月 \\
\addlinespace
Phytomedicine & $\sim$7.0 & 中药活性成分，细胞动物实验 & 2--4 个月 \\
\addlinespace
Frontiers in Pharmacology & $\sim$5.6 & 网络药理学+实验验证，审稿快 & 2--3 个月 \\
\addlinespace
Biomedicine \& Pharmacotherapy & $\sim$7.2 & 肿瘤药理学，实验数据完整 & 2--4 个月 \\
\addlinespace
中国中药杂志（中文核心）& N/A & 中文，中药研究全流程 & 3--6 个月 \\
\bottomrule
\end{tabular}
\end{center}

\subsection{发明专利申请核心框架}
\label{sec:module8_patent}

本研究可申请以下类型的发明专利：

\begin{enumerate}
  \item \textbf{中药组合物专利：}
    \begin{itemize}
      \item 权利要求书 Claim 1（独立权利要求）：「一种含有成分 A（含量 XX\%）、成分 B（含量 XX\%）的中药组合物，其特征在于用于治疗鼻咽癌」
      \item 权利要求书 Claim 2--5（从属权利要求）：制备方法（提取、纯化工艺）、给药途径、联合用药方案
    \end{itemize}
  \item \textbf{筛选方法专利：}
    \begin{itemize}
      \item 独立权利要求：「一种基于多维数据分析筛选中药免疫激活抗肿瘤活性成分的方法，包括步骤一…步骤 N」
      \item 说明书须附实施例（至少 3 个验证数据）和附图（流程图、对接图、实验图）
    \end{itemize}
\end{enumerate}

\begin{tipbox}
「专利申请应在论文发表前完成专利公开（申请受理），以避免论文发表后因现有技术公开而丧失新颖性。建议在获得稳定的 IC50 数据和靶点验证数据后，立即委托专业专利代理机构撰写权利要求书，以确保权利范围的最大化。」
\end{tipbox}

\subsection{课题申报书撰写核心模块}
\label{sec:module8_grant}

国家自然科学基金青年/面上项目申报核心模块：

\begin{enumerate}
  \item \textbf{研究背景（立项依据）：}
    \begin{itemize}
      \item 段落一（3--4 句）：NPC 临床现状与免疫治疗局限性（引 2--3 篇高引用率临床文献，如 JUPITER-02 试验 \cite{mai2021toripalimab}）
      \item 段落二（3--4 句）：中药在免疫调节与协同抗肿瘤中的应用潜力与研究现状
      \item 段落三（2--3 句）：现存研究空白与本研究的科学问题，明确「未解决的具体问题」
    \end{itemize}
  \item \textbf{研究内容：}按「主要研究内容」分条列出，每条对应一个模块操作阶段（3--4 条主要内容）
  \item \textbf{技术路线：}提交技术路线图（PDF 嵌入或单独附件），清晰标注「投入→过程→产出」逻辑
  \item \textbf{创新点（核心竞争力）：}
    \begin{itemize}
      \item 创新点一：建立「多维数据分析筛选中药免疫激活抗肿瘤活性成分」的全新研究范式
      \item 创新点二：以「免疫激活 + 直接抗肿瘤」双终点评价取代传统单一抗肿瘤筛选，明确中药的协同免疫机制
      \item 创新点三：围绕 NPC 这一特色病种，系统揭示特定中药活性成分对 TME 的具体重塑机制
    \end{itemize}
  \item \textbf{可行性分析：}说明前期已完成的数据库构建和预实验结果（引用自己已发表的预实验数据或前期文章）
\end{enumerate}

\subsection{研究数据归档与管理规范}
\label{sec:module8_archive}

执行以下标准化归档规范：

\begin{enumerate}
  \item \textbf{文件夹结构（本地 + 云端备份）：}
\begin{lstlisting}[language={},caption={标准化数据归档目录结构}]
Project_NPC_TCM_Screening/
├── 01_Literature/          # 文献 PDF + 阅读笔记
├── 02_Database_Raw/        # 原始数据库下载文件（不修改）
├── 03_Analysis/            # 分析过程文件
│   ├── R_scripts/          # 所有 R 脚本（版本控制）
│   ├── Python_scripts/     # 所有 Python 脚本
│   ├── Results/            # 分析结果（CSV、PDF 图）
│   └── Docking/            # 分子对接工作目录
├── 04_Wet_Lab/             # 实验原始记录
│   ├── Raw_Data/           # 仪器原始输出（酶标仪 OD 文件、流式 FCS 文件）
│   ├── Processed_Data/     # 处理后 Excel 数据表
│   └── Photos/             # WB 等原始图片（TIFF 格式）
├── 05_Figures/             # 论文图（高分辨率 PDF/TIFF）
├── 06_Manuscripts/         # 论文草稿（按版本命名：v1, v2...）
└── 07_Patent/              # 专利申请材料
\end{lstlisting}

  \item \textbf{版本控制：}所有分析脚本使用 Git 进行版本管理，上传 GitHub（Private 仓库），每次关键节点提交 commit（如「完成 GO/KEGG 富集分析」）。
  \item \textbf{数据安全：}本地数据每周备份至外置硬盘，同时使用 OneDrive/百度网盘进行云端备份（加密处理）。
\end{enumerate}

% ============================================================
% MODULE 9: 研究全流程质量控制与风险规避方案
% ============================================================
\section{模块九：研究全流程质量控制与风险规避方案}
\label{sec:module9}

\begin{infobox}{模块九概览}
\textbf{操作目的：}建立覆盖全流程 9 大模块的质量管控体系，前置核心风险，明确每个环节的判定标准、失败原因排查路径与风险规避方案，确保实验结果的可重复性与学术成果的可靠性。
\end{infobox}

\subsection{生信分析全流程质量控制}
\label{sec:module9_bioinformatics}

\begin{longtable}{p{3cm}p{3.5cm}p{2.5cm}p{4.5cm}}
\caption{生信分析关键质控节点} \label{tab:bioinformatics_qc} \\
\toprule
\textbf{环节} & \textbf{质控指标} & \textbf{合格标准} & \textbf{不合格处理方案} \\
\midrule
\endfirsthead
\toprule
\textbf{环节} & \textbf{质控指标} & \textbf{合格标准} & \textbf{不合格处理方案} \\
\midrule
\endhead
\bottomrule
\endlastfoot
ADME 筛选 & 筛选后成分数量 & $\geq$ 20 种 & 降低 DL 阈值至 0.15 并说明 \\
\addlinespace
靶点检索 & 多库覆盖 & $\geq$ 3 个数据库 & 补充 DrugBank、OMIM 检索 \\
\addlinespace
差异分析 & FDR 校正方法 & BH 法，FDR $<$ 0.05 & 禁止用 P $<$ 0.05 替代 FDR \\
\addlinespace
Venn 交集 & 双效靶点数量 & $\geq$ 10 个 & 扩大免疫靶点基因集来源 \\
\addlinespace
网络分析 & 网络连通性 & 单连通分量 $>$ 80\% & 检查输入文件成分靶点是否完整 \\
\addlinespace
分子对接 & 对接成功率 & $\geq$ 90\% & 检查 PDBQT 格式、格点盒子是否覆盖活性口袋 \\
\addlinespace
Gene ID 转换 & 转换成功率 & $\geq$ 90\% & 使用 Bioconductor org.Hs.eg.db 更新版本 \\
\end{longtable}

\subsection{湿实验全流程质量控制}
\label{sec:module9_wetlab}

\subsubsection{细胞系鉴定规范（STR Profiling）}

\begin{enumerate}
  \item 每批新复苏细胞系在传代 5 代内，送第三方 STR 鉴定机构（如中科院典型培养物保藏中心，武汉大学中国典型培养物保藏中心）。
  \item 将鉴定结果与以下参考数据库比对：
    \begin{itemize}
      \item ATCC STR 数据库：\texttt{https://www.atcc.org/resources/strprofiledatabase}
      \item DSMZ 数据库：\texttt{https://www.dsmz.de/services/services-human-and-animal-cell-lines/str-analysis}
    \end{itemize}
  \item STR 鉴定结果保存于实验档案，在 SCI 论文 Materials \& Methods 中注明（「Cell line identity was authenticated by STR profiling at [institution name]」）。
\end{enumerate}

\subsubsection{试剂批次管理规范}

\begin{itemize}
  \item 关键试剂（FBS、RPMI-1640、CCK-8 试剂盒、ELISA Kit）使用前登记批号，同一实验系列（如 IC50 测定）须使用同批次试剂。
  \item FBS 须经 56$^{\circ}$C 热灭活 30 min（灭活补体），分装 50 mL/管，$-$20$^{\circ}$C 存储，避免反复冻融（$\leq$ 3 次）。
  \item ELISA 标准品在有效期内使用，每次实验重新配制标准品系列（不可使用前次配制的标准品）。
\end{itemize}

\subsection{核心风险规避方案}
\label{sec:module9_risk}

\begin{warningbox}
\textbf{风险一：中药成分数据错误（假阳性 ADME 筛选结果）}

\textbf{风险描述：}TCMSP 数据库中已知存在约 2--5\% 的 SMILES 错误和 OB/DL 数值异常。

\textbf{规避方案：}
\begin{enumerate}
  \item 对所有筛选后的活性成分，使用 RDKit 2023.09.1 重新计算 MW（$<$ 500 Da）、LogP（$<$ 5）、HBD（$\leq$ 5）、HBA（$\leq$ 10）（Lipinski 五规则）进行独立验证，与 TCMSP 值的偏差 $>$ 20\% 的成分需人工核查。
  \item 在 PubChem 确认每个成分的 CID 号，下载标准 SMILES 替换 TCMSP SMILES。
\end{enumerate}
\end{warningbox}

\begin{warningbox}
\textbf{风险二：靶点筛选假阳性（GeneCards Score 过低的基因纳入）}

\textbf{风险描述：}GeneCards 数据库包含基于文本挖掘的低可信度靶点，若 Score $<$ 2 的靶点大量纳入会引入大量假阳性。

\textbf{规避方案：}
\begin{enumerate}
  \item 设置 GeneCards Relevance Score $\geq$ 2.0 的截止值，同时要求靶点在 DisGeNET（Score $\geq$ 0.1）中有记录。
  \item 对最终进入 Hub 靶点分析的基因，在 PubMed 检索是否有 NPC 直接相关的同行评审文献。
\end{enumerate}
\end{warningbox}

\begin{warningbox}
\textbf{风险三：分子对接假阳性（对接结果与实验不相符）}

\textbf{风险描述：}分子对接是虚拟筛选工具，结合自由能好（$\leq$ $-$8 kcal/mol）不代表体外实验一定有活性，假阳性率可达 30--50\%。

\textbf{规避方案：}
\begin{enumerate}
  \item 将分子对接结果与网络药理学 Degree 排名、公共数据集表达验证联合评分（模块四的多维打分体系），不单纯依赖对接分数筛选候选成分。
  \item 优先选择具有同类化合物（相似骨架）已有实验验证的候选成分。
  \item 若经费允许，可在湿实验前进行分子动力学（MD）模拟（100 ns 以上），验证对接构象的稳定性。
\end{enumerate}
\end{warningbox}

\begin{warningbox}
\textbf{风险四：细胞实验污染（支原体污染）}

\textbf{风险描述：}支原体污染是细胞实验最常见的隐患，可导致细胞生长异常、实验结果偏差，且肉眼无法察觉。

\textbf{规避方案：}
\begin{enumerate}
  \item 每月执行一次支原体检测（PCR 法或荧光染色法，推荐 EZ-PCR Mycoplasma Test Kit，Biological Industries）。
  \item 发现污染立即销毁污染细胞，不可用抗生素「治疗」。
  \item 所有试剂操作在超净台内进行，废液收集于废液瓶（10\% 次氯酸钠灭活后按规定处理）。
\end{enumerate}
\end{warningbox}

\begin{warningbox}
\textbf{风险五：实验结果不可重复}

\textbf{风险描述：}由于细胞传代数、试剂批次、DMSO 浓度控制不当等因素，同一实验在不同时间重复时 IC50 偏差 $>$ 50\%。

\textbf{规避方案：}
\begin{enumerate}
  \item 所有 CCK-8/ELISA 实验须保留原始 OD 值记录（以 Excel 文件存档），并在每次实验中同时运行阳性对照药（顺铂），以历史顺铂 IC50 范围验证实验体系稳定性。
  \item 严格控制细胞传代代次（实验细胞 passage 5--20）和 DMSO 终浓度（$\leq$ 0.1\%）。
  \item 生物学重复（biological replicate）与技术重复（technical replicate）须区分记录，最终数据分析仅使用生物学重复计算统计。
\end{enumerate}
\end{warningbox}

% ============================================================
% 新增：完整Plan B风险应急方案（针对关键场景）
% ============================================================
\begin{warningbox}
\textbf{风险六（核心 Plan B）--- 70 味候选中药高通量初筛双终点全部阴性的完整应急方案}

\textbf{风险描述：}若 70 味候选中药在 100 $\mu$g/mL 单浓度初筛（模块四补充方案）中，抗肿瘤（CCK-8 抑制率 $<$ 30\%）与免疫激活（TNF-$\alpha$ 倍变 $<$ 2.0）双终点均未达命中阈值，属于高概率风险（中药粗提物质量不稳定、成分含量低是主要原因）。须在 \textbf{2 周内} 启动本应急方案，确保研究进度不中断。

\textbf{风险发生概率评估：}根据同类中药高通量筛选研究（Chen et al., 2022; Zhang et al., 2023），70 味中药中通常有 10--20\% 在 100 $\mu$g/mL 单浓度下表现出明显活性，全部阴性概率约 5--10\%。

\textbf{分层应急方案（按优先级严格执行）：}

\begin{enumerate}
  \item \textbf{Step 1（第 1--3 天）--- 优先排查实验体系有效性：}
  \begin{itemize}
    \item 确认顺铂阳性对照 CNE2 抑制率是否在 75--95\% 范围内（若偏低，说明细胞状态异常而非药物问题）。
    \item 检查 DMSO 工作液配制是否准确（重新检算稀释比例，并用 DMSO 对照检测 0.1\% 无毒性）。
    \item 检查细胞传代代次（passage $>$ 25 的 CNE2 对化学物质敏感性显著降低，须换低代次细胞）。
    \item 若实验体系异常，修正后重新实验，不启动后续 Step。
  \end{itemize}

  \item \textbf{Step 2（第 3--7 天）--- 提高筛选浓度至 200 $\mu$g/mL：}
  \begin{itemize}
    \item 同时保留 100 $\mu$g/mL 数据，以 200 $\mu$g/mL 重复初筛，确认是否存在浓度依赖效应。
    \item 若 200 $\mu$g/mL 有活性（抑制率 $\geq$ 30\%），在论文中充分讨论高浓度有效性的临床转化局限（需体内实验验证吸收分布），但研究可继续推进。
    \item 若 200 $\mu$g/mL 仍全部阴性，进入 Step 3。
  \end{itemize}

  \item \textbf{Step 3（第 1--2 周）--- 重构候选中药库（扩库补筛）：}
  \begin{itemize}
    \item 重新检索 BATMAN-TCM 2.0 和 ETCM，基于「清热解毒」「活血化瘀」「扶正固本」三大功效类别，补充筛选 \textbf{20--30 味新候选中药}（排除与原 70 味重叠者）。
    \item 优先纳入已有明确 \textit{in vitro} 抗肿瘤或免疫激活文献报道的中药（PubMed 检索近 5 年，以 CNE1/CNE2/鼻咽癌为筛选条件）。
    \item 使用相同 SOP 对新增候选中药重新执行高通量初筛。
  \end{itemize}

  \item \textbf{Step 4（第 2 周，须导师确认）--- 切换研究策略：靶向单体验证路线：}
  \begin{itemize}
    \item 跳过全库表型筛选，直接对多维评分 $\geq$ 7 分（模块四候选成分评分体系）的前 \textbf{3--5 种候选单体化合物}进行 IC50 全剂量实验。
    \item 研究框架从「中药粗提物筛选」转为「基于生信预测的单体化合物靶向验证」，聚焦单体成分—核心靶点—机制研究（适合发表机制类 SCI 文章，IF 可达 4.0--8.0）。
    \item 此策略须书面更新研究计划（更新技术路线图、研究内容描述），提交导师审批后方可执行。
  \end{itemize}

  \item \textbf{Step 5（第 3 周）--- 向导师与相关人员正式汇报调整方案：}
  \begin{itemize}
    \item 以 PPT 形式汇报原 70 味初筛结果（客观呈现，不选择性呈现数据）。
    \item 说明选择 Step 3 或 Step 4 的理由及预期输出。
    \item 更新时间轴（模块九进度规划表），确认关键节点是否需要调整。
  \end{itemize}
\end{enumerate}
\end{warningbox}

\begin{warningbox}
\textbf{风险七（核心 Plan B）--- 细胞大规模交叉污染（多细胞系同时操作场景）完整应急方案}

\textbf{风险描述：}本研究需同时维护 CNE1/CNE2（鼻咽癌）、Jurkat（T 细胞）、THP-1（巨噬细胞）等多个细胞系，多细胞系同时操作时支原体或真菌/细菌交叉污染风险显著提升。一旦发现全批污染，须立即启动应急流程，最小化实验损失并重建细胞库。

\textbf{污染识别标准：}
\begin{itemize}
  \item \textbf{真菌污染：}培养基出现浑浊、颜色变化、镜下可见菌丝，通常在 24--48 h 内明显可见。
  \item \textbf{细菌污染：}培养基迅速变黄（pH 下降），镜下可见运动性颗粒，通常 12--24 h 内即明显。
  \item \textbf{支原体污染（隐性）：}镜检无异常，但细胞生长减慢、染色体异常；须用 PCR 法主动检测（建议每月 1 次例行检查）。
\end{itemize}

\textbf{分级应急处置方案（48 h 内完成）：}
\begin{enumerate}
  \item \textbf{污染确认（0--6 h）：}立即隔离污染瓶（转至独立区域，4$^{\circ}$C 保存待核查），对相邻培养瓶逐一镜检；使用支原体 PCR 检测试剂盒（EZ-PCR Mycoplasma Test Kit，Biological Industries，目录号 20-700-20）检测，3--5 h 出结果。
  \item \textbf{污染处置（6--24 h）：}所有确认污染的细胞瓶加入 \textbf{终浓度 1\% 次氯酸钠}（接触时间 $\geq$ 30 min），统一按医疗废弃物处理（禁止直接倒入下水道）；对超净台进行深度消毒（75\% 酒精擦拭 + 紫外照射 60 min）；检查所有共用试剂（PBS、培养基分装瓶）是否为污染源（取样镜检）。
  \item \textbf{溯源排查（24--48 h）：}逐项排查污染来源：超净台 HEPA 滤网状态（建议每年更换）→ 操作规范（是否在超净台外分液）→ 试剂开封状态（是否存在共用污染）→ 人员操作习惯（新加入人员须重新培训无菌操作）。
  \item \textbf{重建细胞库（1--2 周）：}
  \begin{itemize}
    \item 首选方案：从本实验室液氮库调取最早期冻存管（passage $\leq$ 5），复苏后扩培，复苏细胞第 1 天即进行支原体检测（PCR 阴性确认后方可继续实验）。
    \item 备选方案（液氮库耗尽）：CNE1/CNE2 可从中国科学院典型培养物保藏委员会细胞库（CCTCC）或国家细胞资源共享平台（NCRSP）重新购置；Jurkat/THP-1 可从 ATCC 购置。
    \item 新细胞复苏后立即建立新的冻存梯度（passage 5、8、10 各存 $\geq$ 3 管），冻存前支原体 PCR 阴性确认。
  \end{itemize}
  \item \textbf{进度影响评估与恢复（第 2--3 周）：}向导师汇报污染事件（书面记录，含原因分析和预防措施），对照时间轴评估影响的实验周期（通常损失 2--4 周），压缩可并行执行的非细胞实验（生信分析、数据可视化）以弥补进度。
\end{enumerate}

\textbf{污染预防长效机制（日常规范）：}
\begin{itemize}
  \item \textbf{物理隔离：}不同细胞系使用不同标记的分液器和吸管盒，禁止混用。
  \item \textbf{例行检测：}每月进行 1 次支原体 PCR 检测（所有在用细胞系）。
  \item \textbf{冻存梯度：}每新建或复苏一个细胞系，在 passage 5、10、15 各存 $\geq$ 5 管（液氮 + 气相备份两套）。
\end{itemize}
\end{warningbox}

\begin{warningbox}
\textbf{风险八（核心 Plan B）--- 原代 T 细胞实验系统性失败的完整降级应急预案}

\textbf{风险描述：}若小鼠脾脏 T 细胞分离后活力持续低于 70\%，或激活后免疫指标（IL-2/IFN-$\gamma$）无法达到 ConA 阳性对照的 5 倍升高标准，提示实验体系存在系统性问题（不仅仅是单次操作失误），须执行本系统性应急预案。

\textbf{应急降级路径（三级降级方案）：}

\begin{enumerate}
  \item \textbf{一级应急（首选，2--5 天）--- 排查并修复原代系统问题：}
  \begin{itemize}
    \item 确认 ACK 裂解液配制是否正确（pH 须在 7.2--7.4，超时裂解 $>$ 3 min 会损伤 T 细胞）。
    \item 确认 ConA 批次是否有效（ConA 须 4$^{\circ}$C 保存，有效期内使用；可用已知阳性批次 ConA 平行测试）。
    \item 确认 FBS 批次：换用经过 T 细胞培养认证的 FBS（每批次进货时进行 72 h T 细胞存活率测试）。
    \item 修复后重新取材（新鲜小鼠脾脏）重复实验。
  \end{itemize}

  \item \textbf{二级应急（5--14 天）--- 切换激活方案：}
  \begin{itemize}
    \item 将 ConA 激活切换至 \textbf{Anti-CD3$\varepsilon$/CD28 板结合法}（排除 ConA 批次问题）。
    \item 或引入 \textbf{T-Activator CD3/CD28 Dynabeads}（Thermo Fisher，11456D），该方案对操作规范性要求更低，激活效率更稳定。
    \item 同步用 PBMC 系统验证（排除动物原因）。
  \end{itemize}

  \item \textbf{三级应急（降级至 Jurkat/THP-1 替代验证，即时可执行）：}
  \begin{itemize}
    \item 临时全面切换至 Jurkat T 细胞 + THP-1 巨噬细胞系统（均为永生化细胞系，操作成熟稳定）完成免疫激活验证。
    \item 同步联系有原代 T 细胞实验经验的实验室寻求技术支持或合作（技术合作须明确挂名条件）。
    \item 在论文方法部分明确说明「原代 T 细胞实验为计划验证内容，受条件限制以 Jurkat/THP-1 细胞系作为替代验证，后续研究将补充原代数据」，此为本领域可接受的方法描述方式（审稿人通常接受此类补充说明，尤其是 IF $\leq$ 5.0 的期刊）。
  \end{itemize}
\end{enumerate}

\textbf{风险预警指标（日常监控）：}
\begin{itemize}
  \item 每次分离后立即台盼蓝计数，活力 $<$ 85\% 即触发 Step 1 排查，无需等待实验结束。
  \item ConA 阳性对照 IL-2 升高 $<$ 3 倍（低于合格标准 5 倍但仍有响应），须记录并执行第二次确认实验（更换 FBS 后重复），两次均低于 5 倍方可判定为系统性问题并启动降级方案。
\end{itemize}
\end{warningbox}

\subsection{研究进度管理}
\label{sec:module9_timeline}

\begin{longtable}{p{1.5cm}p{3cm}p{5cm}p{4cm}}
\caption{研究全流程进度规划（参考时间轴，硕士 2--3 年期）} \label{tab:timeline} \\
\toprule
\textbf{阶段} & \textbf{时间节点} & \textbf{关键里程碑（Milestone）} & \textbf{关键风险与应对} \\
\midrule
\endfirsthead
\toprule
\textbf{阶段} & \textbf{时间节点} & \textbf{关键里程碑（Milestone）} & \textbf{关键风险与应对} \\
\midrule
\endhead
\bottomrule
\endlastfoot
Phase 0 & 第 1--4 周 & 完成模块一（研究方案设计 + 工具配置）& 伦理审批延迟：提前 2 个月递交 \\
\addlinespace
Phase 1 & 第 5--8 周 & 完成模块二（数据库构建）+ 模块三（靶点集筛选）& 数据库访问异常：准备 HERB 等备用 \\
\addlinespace
Phase 2 & 第 9--14 周 & 完成模块四（虚拟筛选）+ 模块五（多组学验证）& 分子对接结果不理想：扩大候选靶点范围 \\
\addlinespace
Phase 3 & 第 15--26 周 & 完成模块六（湿实验）第一轮（CCK-8 + 凋亡）& 细胞实验失败：预留 4 周重复实验时间 \\
\addlinespace
Phase 4 & 第 27--36 周 & 完成模块六（免疫实验 + 机制验证）& ELISA 结果不显著：调整阳性对照方案 \\
\addlinespace
Phase 5 & 第 37--44 周 & 完成模块七--八（数据分析 + 论文撰写第一稿）& 导师修改意见多：预留 4 周修改时间 \\
\addlinespace
Phase 6 & 第 45--52 周 & 完成投稿、答复审稿意见、修改后接收 & 拒稿：准备备用期刊（降一档期刊）\\
\bottomrule
\end{longtable}

\begin{qcbox}
\textbf{全流程质控检查清单（结题前执行）：}
\begin{itemize}
  \item 模块一：研究方案技术路线图已存档，导师签字确认 \ding{52}
  \item 模块二：中药-成分数据库已建立，OB/DL 阈值符合金标准，SMILES 已 PubChem 验证 \ding{52}
  \item 模块三：双效靶点交集 $\geq$ 10 个，Venn 图已生成，Gene Symbol 标准化完成 \ding{52}
  \item 模块四：网络图 + 富集图 + 对接图已完成（300 dpi PDF），候选成分多维评分记录完整 \ding{52}
  \item 模块五：差异表达、生存分析、免疫浸润分析已完成，GEO 数据集编号已记录 \ding{52}
  \item 模块六：所有细胞实验 $\geq$ 3 次生物重复，STR 鉴定结果存档，原始数据文件完整 \ding{52}
  \item 模块七：所有统计方法已记录，所有图表分辨率 $\geq$ 300 dpi \ding{52}
  \item 模块八：论文已按 IMRaD 格式撰写，参考文献 $\geq$ 40 篇，格式符合目标期刊要求 \ding{52}
  \item 模块九：所有风险清单已执行，实验记录本完整，数据已双备份（本地 + 云端）\ding{52}
\end{itemize}
\end{qcbox}

% ============================================================
% CONCLUSION & REFERENCES
% ============================================================
\section*{结语}
\addcontentsline{toc}{section}{结语}

本手册依托「研究设计→多维数据库构建→双维度靶点集筛选→虚拟筛选→多组学验证→湿实验全流程验证→数据分析可视化→成果输出→全流程质控风控」的九模块闭环科研框架，以鼻咽癌（NPC，CNE1/CNE2 细胞系）为核心示例，系统呈现了从中药多维数据整合到「免疫激活 + 直接抗肿瘤」双终点实验验证的全链条标准化操作规程。

手册的核心价值体现在以下三个维度：

\begin{enumerate}
  \item \textbf{操作可落地性：}所有操作步骤均拆解至「点击级」精度，明确软件版本（Cytoscape 3.10.1、AutoDock Vina 1.2.3、R 4.3.3、GraphPad Prism 9.0 等）、参数阈值（OB $\geq$ 30\%、DL $\geq$ 0.18、STRING 置信度 $\geq$ 0.700、对接强结合阈值 $\leq$ $-$8.0 kcal/mol），零基础研究者可依照手册 1:1 复现实验流程。

  \item \textbf{质控前置性：}每个模块均配套独立的质控合格标准与避坑指南，将中药筛选领域的高频风险（数据假阳性、细胞污染、实验不可重复）前置于每个操作节点，保障研究结果的科学性与可靠性。

  \item \textbf{成果导向性：}所有操作节点的输出结果均明确标注可直接用于 SCI 论文、发明专利或课题申报书的对应章节，实现从科研操作到成果转化的全场景适配。
\end{enumerate}

研究者在使用本手册时，应结合自身研究体系（中药类型、瘤种、细胞系）按照【模板替换位】标注进行个性化适配，并在每个关键操作节点严格执行质控标准，确保科研闭环的完整性与结果的可重复性。

\vspace{12pt}

\begin{center}
{\large\bfseries\color{darkblue} ——全流程完毕，预祝研究顺利——}
\end{center}

% ============================================================
% BIBLIOGRAPHY
% ============================================================
\clearpage
\addcontentsline{toc}{section}{参考文献}
\bibliographystyle{unsrtnat}
\bibliography{../references/references}

\clearpage

% ============================================================
% APPENDIX
% ============================================================
\appendix
\section{附录 A：核心 R 包安装脚本}
\label{appendix:r_packages}

\begin{lstlisting}[language=R,caption={R 4.3.3 核心包一键安装脚本}]
# ==== R 4.3.3 核心包一键安装脚本 ====
# 运行环境: R 4.3.3, Bioconductor 3.18
# 安装时间: 约 20-40 分钟（取决于网络速度）

# 1. 安装 CRAN 包
cran_packages <- c(
  "dplyr", "ggplot2", "readr", "openxlsx",
  "VennDiagram", "pheatmap", "ggrepel", "survival",
  "survminer", "tidyr", "stringr", "reshape2",
  "RColorBrewer", "gridExtra", "cowplot", "viridis",
  "pROC", "caret", "randomForest"
)

install.packages(cran_packages, repos = "https://cran.r-project.org/")

# 2. 安装 Bioconductor 包
if (!require("BiocManager")) install.packages("BiocManager")
BiocManager::install(version = "3.18")

bioc_packages <- c(
  "GEOquery",          # GEO 数据下载
  "limma",             # 差异基因分析
  "affy",              # 芯片数据处理
  "clusterProfiler",   # GO/KEGG 富集分析
  "org.Hs.eg.db",      # 人类基因注释数据库
  "AnnotationDbi",     # 基因 ID 转换
  "Seurat",            # 单细胞测序分析
  "enrichplot",        # 富集分析可视化
  "fgsea",             # GSEA 分析
  "immunedeconv",      # 免疫浸润分析
  "DESeq2",            # RNA-seq 差异分析（备选）
  "WGCNA"              # 加权基因共表达网络
)

BiocManager::install(bioc_packages)

# 3. 验证安装
cat("=== 安装验证 ===\n")
for (pkg in c(cran_packages, bioc_packages)) {
  if (requireNamespace(pkg, quietly = TRUE)) {
    cat(sprintf("✓ %s\n", pkg))
  } else {
    cat(sprintf("✗ %s - 安装失败，请手动重试\n", pkg))
  }
}
\end{lstlisting}

\section{附录 B：SwissTargetPrediction 靶点预测操作（替代在线操作）}
\label{appendix:stp}

若无法访问 \database{SwissTargetPrediction} 在线平台（\texttt{http://www.swisstargetprediction.ch}），可使用 Python API 批量预测：

\begin{lstlisting}[language=Python,caption={SwissTargetPrediction API 批量靶点预测}]
# ==== SwissTargetPrediction 批量靶点预测 ====
import requests
import pandas as pd
import time
import json

def predict_targets_stp(smiles: str, organism: str = "Homo sapiens") -> pd.DataFrame:
    """
    通过 SwissTargetPrediction API 预测靶点
    参数: smiles - 化合物 SMILES 字符串
          organism - 物种（默认人类）
    返回: 靶点预测结果 DataFrame
    """
    url = "http://www.swisstargetprediction.ch/predict.php"
    data = {
        "smiles": smiles,
        "organism": organism
    }

    try:
        response = requests.post(url, data=data, timeout=60)
        if response.status_code == 200:
            result = response.json()
            df = pd.DataFrame(result)
            # 筛选概率 >= 0.1 的靶点
            df = df[df['probability'] >= 0.1]
            return df
        else:
            print(f"API 返回错误: {response.status_code}")
            return pd.DataFrame()
    except Exception as e:
        print(f"请求失败: {e}")
        return pd.DataFrame()

# 批量执行
compounds_df = pd.read_csv("results/02_adme_filtered_compounds.csv")
all_targets = []

for idx, row in compounds_df.iterrows():
    print(f"预测靶点: {row['compound_name']} ({idx+1}/{len(compounds_df)})")
    targets_df = predict_targets_stp(row['smiles'])
    if not targets_df.empty:
        targets_df['compound_name'] = row['compound_name']
        targets_df['compound_id'] = row.get('compound_id', f"MOL{idx:03d}")
        all_targets.append(targets_df)
    time.sleep(2)  # 避免请求过于频繁

# 合并所有预测结果
if all_targets:
    results_df = pd.concat(all_targets, ignore_index=True)
    results_df.to_csv("results/03_tcm_predicted_targets.csv", index=False)
    print(f"靶点预测完成，共 {len(results_df)} 条记录")
else:
    print("未获取到靶点预测结果，请检查 API 访问")
\end{lstlisting}

\section{附录 C：常见错误与解决方案（FAQ）}
\label{appendix:faq}

\begin{longtable}{p{4cm}p{4cm}p{6cm}}
\caption{全流程常见报错与解决方案} \label{tab:faq} \\
\toprule
\textbf{报错描述} & \textbf{发生环节} & \textbf{解决方案} \\
\midrule
\endfirsthead
\toprule
\textbf{报错描述} & \textbf{发生环节} & \textbf{解决方案} \\
\midrule
\endhead
\bottomrule
\endlastfoot
\texttt{ERROR: GEO file download failed} & R/GEOquery & 检查网络连接；设置 \texttt{options(timeout=300)}；改用 \texttt{destdir} 参数指定本地目录 \\
\addlinespace
\texttt{object 'adj.P.Val' not found} & R/limma & 检查 topTable 调用参数，使用 \texttt{adjust.method="BH"} 而非 \texttt{p.adjust.method} \\
\addlinespace
Cytoscape 导入后网络为空 & Cytoscape 3.10.1 & 检查 CSV 列名是否含中文或特殊字符；Source/Target 列名须为纯英文 \\
\addlinespace
AutoDock Vina 对接分数为正值 & 分子对接 & 正值通常表示配体体积过大超出格点盒子；扩大 size\_x/y/z $\geq$ 30 Å；或配体 PDBQT 格式错误 \\
\addlinespace
PyMOL 加载 PDBQT 文件显示异常 & PyMOL 可视化 & PDBQT 格式与 PDB 不完全兼容；改用 \texttt{load compound\_out.pdbqt, format="pdbqt"} 命令 \\
\addlinespace
CCK-8 OD 值异常（所有孔 OD $>$ 3.0 或 $<$ 0.1）& 细胞实验 & OD $>$ 3.0：铺板密度过高或 CCK-8 孵育时间过长；OD $<$ 0.1：铺板密度过低或细胞死亡，检查细胞活力 \\
\addlinespace
ELISA 标准曲线 R² $<$ 0.99 & ELISA 实验 & 标准品配制步骤出错（移液失误）；重新配制标准品；检查是否使用 2 倍梯度稀释 \\
\addlinespace
WB 条带无信号 & Western Blot & 一抗浓度过低（提高 2-5 倍）；电转效率差（检查 PVDF 膜激活）；靶蛋白 MW 与 PAGE 胶浓度不匹配 \\
\addlinespace
qRT-PCR 熔解曲线多峰 & qRT-PCR & 引物二聚体或非特异扩增；更换引物序列；适当提高退火温度（$+$2$^{\circ}$C）\\
\bottomrule
\end{longtable}

% End of document
\end{document}
